
\Data\ID{1003020301}
\Updated{2020-07-15 03:07:12}
\Level{A}
\SubArea{Systems of equations and inequalities}
\LANGen{
\Question{
Find the solution set of the equation
\[ 2x-\frac{x+2y}3=2+\frac83y \]
in \( \mathbb{R}\times\mathbb{R} \).}
{\( \left\{\left[2y+\frac65;y\right],y\in\mathbb{R}\right\} \)}{\( \left\{\left[2y+\frac65;\frac x2-\frac35\right],x\in\mathbb{R},y\in\mathbb{R}\right\} \)}{\( \left\{\left[\frac{6+6y}5;y\right],y\in\mathbb{R}\right\} \)}{\( \emptyset \)}{}{}}
\LANGpl{
\Question{
Wyznacz zbiór rozwiązań równania
\[ 2x-\frac{x+2y}3=2+\frac83y \]
w \( \mathbb{R}\times\mathbb{R} \).}
{\( \left\{\left[2y+\frac65;y\right],y\in\mathbb{R}\right\} \)}{\( \left\{\left[2y+\frac65;\frac x2-\frac35\right],x\in\mathbb{R},y\in\mathbb{R}\right\} \)}{\( \left\{\left[\frac{6+6y}5;y\right],y\in\mathbb{R}\right\} \)}{\( \emptyset \)}{}{}}
\LANGcs{
\Question{
Určete množinu řešení rovnice
\[ 2x-\frac{x+2y}3=2+\frac83y \]
v oboru \( \mathbb{R}\times\mathbb{R} \).}
{\( \left\{\left[2y+\frac65;y\right],y\in\mathbb{R}\right\} \)}{\( \left\{\left[2y+\frac65;\frac x2-\frac35\right],x\in\mathbb{R},y\in\mathbb{R}\right\} \)}{\( \left\{\left[\frac{6+6y}5;y\right],y\in\mathbb{R}\right\} \)}{\( \emptyset \)}{}{}}
\LANGsk{
\Question{
Určte množinu riešení rovnice
\[ 2x-\frac{x+2y}3=2+\frac83y \]
v obore \( \mathbb{R}\times\mathbb{R} \).}
{\( \left\{\left[2y+\frac65;y\right],y\in\mathbb{R}\right\} \)}{\( \left\{\left[2y+\frac65;\frac x2-\frac35\right],x\in\mathbb{R},y\in\mathbb{R}\right\} \)}{\( \left\{\left[\frac{6+6y}5;y\right],y\in\mathbb{R}\right\} \)}{\( \emptyset \)}{}{}}
\LANGes{
\Question{
Calcula el conjunto de soluciones de la siguiente ecuación
\[ 2x-\frac{x+2y}3=2+\frac83y \]
en \( \mathbb{R}\times\mathbb{R} \).}
{\( \left\{\left[2y+\frac65;y\right],y\in\mathbb{R}\right\} \)}{\( \left\{\left[2y+\frac65;\frac x2-\frac35\right],x\in\mathbb{R},y\in\mathbb{R}\right\} \)}{\( \left\{\left[\frac{6+6y}5;y\right],y\in\mathbb{R}\right\} \)}{\( \emptyset \)}{}{}}
\End

\Data\ID{1003020302}
\Updated{2020-07-15 03:07:12}
\Level{A}
\SubArea{Systems of equations and inequalities}
\LANGen{
\Question{
Assuming \( [x;y]\in\mathbb{R}\times\mathbb{R} \), solve the equation
\[ x-y-\frac{x-y}2=\frac{x-y}3 \]
Decide which of the answers below does not express the set of roots.}
{\( \left\{[x;y],x\in\mathbb{R},y\in\mathbb{R}\right\} \)}{\( \left\{[x;x],x\in\mathbb{R}\right\} \)}{\( \left\{[y;y],y\in\mathbb{R}\right\} \)}{\( \left\{[t;t],t\in\mathbb{R}\right\} \)}{}{}}
\LANGpl{
\Question{
Zakładając, że \( [x;y]\in\mathbb{R}\times\mathbb{R} \), rozwiąż równanie 
\[ x-y-\frac{x-y}2=\frac{x-y}3 \]
Zdecyduj, która z poniższych odpowiedzi nie przedstawia zbioru pierwiastków.}
{\( \left\{[x;y],x\in\mathbb{R},y\in\mathbb{R}\right\} \)}{\( \left\{[x;x],x\in\mathbb{R}\right\} \)}{\( \left\{[y;y],y\in\mathbb{R}\right\} \)}{\( \left\{[t;t],t\in\mathbb{R}\right\} \)}{}{}}
\LANGcs{
\Question{
Pro \( [x;y]\in\mathbb{R}\times\mathbb{R}\) řešte rovnici
\[ x-y-\frac{x-y}2=\frac{x-y}3 \]
a poté rozhodněte, která z následujících variant nevyjadřuje množinu kořenů.}
{\( \left\{[x;y],x\in\mathbb{R},y\in\mathbb{R}\right\} \)}{\( \left\{[x;x],x\in\mathbb{R}\right\} \)}{\( \left\{[y;y],y\in\mathbb{R}\right\} \)}{\( \left\{[t;t],t\in\mathbb{R}\right\} \)}{}{}}
\LANGsk{
\Question{
Pre \( [x;y]\in\mathbb{R}\times\mathbb{R} \) riešte rovnicu
\[ x-y-\frac{x-y}2=\frac{x-y}3 \]
Rozhodnite, ktorá z nasledujúcich odpovedí nevyjadruje množinu koreňov.}
{\( \left\{[x;y],x\in\mathbb{R},y\in\mathbb{R}\right\} \)}{\( \left\{[x;x],x\in\mathbb{R}\right\} \)}{\( \left\{[y;y],y\in\mathbb{R}\right\} \)}{\( \left\{[t;t],t\in\mathbb{R}\right\} \)}{}{}}
\LANGes{
\Question{
Suponiendo que \( [x;y]\in\mathbb{R}\times\mathbb{R} \), resuelve la siguiente ecuación
\[ x-y-\frac{x-y}2=\frac{x-y}3 \]
A continuación, decide cuál de las respuestas no expresa el conjunto de las raíces.}
{\( \left\{[x;y],x\in\mathbb{R},y\in\mathbb{R}\right\} \)}{\( \left\{[x;x],x\in\mathbb{R}\right\} \)}{\( \left\{[y;y],y\in\mathbb{R}\right\} \)}{\( \left\{[t;t],t\in\mathbb{R}\right\} \)}{}{}}
\End

\Data\ID{1003020303}
\Updated{2020-07-15 03:07:12}
\Level{A}
\SubArea{Systems of equations and inequalities}
\LANGen{
\Question{
Find the solution set of the equation
\[1-\frac{x-2y}4=x+\frac{y+2}2\]
in \( \mathbb{R}\times\mathbb{R} \).}
{\( \left\{[0;y],y\in\mathbb{R}\right\} \)}{\( \left\{[x;y],x\in\mathbb{R},y\in\mathbb{R}\right\} \)}{\( \left\{\left[\frac{-4y}5;y\right],y\in\mathbb{R}\right\} \)}{\( \{0\} \)}{\( \emptyset \)}{}}
\LANGpl{
\Question{
Wyznacz zbiór rozwiązań równania
\[1-\frac{x-2y}4=x+\frac{y+2}2\]
w \( \mathbb{R}\times\mathbb{R} \).}
{\( \left\{[0;y],y\in\mathbb{R}\right\} \)}{\( \left\{[x;y],x\in\mathbb{R},y\in\mathbb{R}\right\} \)}{\( \left\{\left[\frac{-4y}5;y\right],y\in\mathbb{R}\right\} \)}{\( \{0\} \)}{\( \emptyset \)}{}}
\LANGcs{
\Question{
Určete množinu řešení rovnice
\[1-\frac{x-2y}4=x+\frac{y+2}2\]
v oboru \( \mathbb{R}\times\mathbb{R} \).}
{\( \left\{[0;y],y\in\mathbb{R}\right\} \)}{\( \left\{[x;y],x\in\mathbb{R},y\in\mathbb{R}\right\} \)}{\( \left\{\left[\frac{-4y}5;y\right],y\in\mathbb{R}\right\} \)}{\( \{0\} \)}{\( \emptyset \)}{}}
\LANGsk{
\Question{
Určte množinu riešení rovnice
\[1-\frac{x-2y}4=x+\frac{y+2}2\]
v obore \( \mathbb{R}\times\mathbb{R} \).}
{\( \left\{[0;y],y\in\mathbb{R}\right\} \)}{\( \left\{[x;y],x\in\mathbb{R},y\in\mathbb{R}\right\} \)}{\( \left\{\left[\frac{-4y}5;y\right],y\in\mathbb{R}\right\} \)}{\( \{0\} \)}{\( \emptyset \)}{}}
\LANGes{
\Question{
Calcula el conjunto de soluciones de la siguiente ecuación
\[1-\frac{x-2y}4=x+\frac{y+2}2\]
en \( \mathbb{R}\times\mathbb{R} \).}
{\( \left\{[0;y],y\in\mathbb{R}\right\} \)}{\( \left\{[x;y],x\in\mathbb{R},y\in\mathbb{R}\right\} \)}{\( \left\{\left[\frac{-4y}5;y\right],y\in\mathbb{R}\right\} \)}{\( \{0\} \)}{\( \emptyset \)}{}}
\End

\Data\ID{1003020304}
\Updated{2020-07-15 03:07:12}
\Level{A}
\SubArea{Systems of equations and inequalities}
\LANGen{
\Question{
Find the solution set of the equation
\[1-\left[4x+3\cdot(x-y)\right]=\frac{1-14x}2-\frac{3-6y}2\]
in \( \mathbb{R}\times\mathbb{R} \).}
{\( \emptyset \)}{\( \left\{[x;y],x\in\mathbb{R},y\in\mathbb{R}\right\} \)}{\( \left\{\left[x;\frac13+x\right],x\in\mathbb{R}\right\} \)}{\( \left\{\left[x;-\frac13\right],x\in\mathbb{R}\right\} \)}{}{}}
\LANGpl{
\Question{
Wyznacz zbiór rozwiązań równania
\[1-\left[4x+3\cdot(x-y)\right]=\frac{1-14x}2-\frac{3-6y}2\]
w \( \mathbb{R}\times\mathbb{R} \).}
{\( \emptyset \)}{\( \left\{[x;y],x\in\mathbb{R},y\in\mathbb{R}\right\} \)}{\( \left\{\left[x;\frac13+x\right],x\in\mathbb{R}\right\} \)}{\( \left\{\left[x;-\frac13\right],x\in\mathbb{R}\right\} \)}{}{}}
\LANGcs{
\Question{
Určete množinu řešení rovnice
\[1-\left[4x+3\cdot(x-y)\right]=\frac{1-14x}2-\frac{3-6y}2\]
v oboru \( \mathbb{R}\times\mathbb{R} \).}
{\( \emptyset \)}{\( \left\{[x;y],x\in\mathbb{R},y\in\mathbb{R}\right\} \)}{\( \left\{\left[x;\frac13+x\right],x\in\mathbb{R}\right\} \)}{\( \left\{\left[x;-\frac13\right],x\in\mathbb{R}\right\} \)}{}{}}
\LANGsk{
\Question{
Určte množinu riešenia rovnice
\[1-\left[4x+3\cdot(x-y)\right]=\frac{1-14x}2-\frac{3-6y}2\]
v obore \( \mathbb{R}\times\mathbb{R} \).}
{\( \emptyset \)}{\( \left\{[x;y],x\in\mathbb{R},y\in\mathbb{R}\right\} \)}{\( \left\{\left[x;\frac13+x\right],x\in\mathbb{R}\right\} \)}{\( \left\{\left[x;-\frac13\right],x\in\mathbb{R}\right\} \)}{}{}}
\LANGes{
\Question{
Calcula el conjunto de soluciones de la siguiente ecuación
\[1-\left[4x+3\cdot(x-y)\right]=\frac{1-14x}2-\frac{3-6y}2\]
en \( \mathbb{R}\times\mathbb{R} \).}
{\( \emptyset \)}{\( \left\{[x;y],x\in\mathbb{R},y\in\mathbb{R}\right\} \)}{\( \left\{\left[x;\frac13+x\right],x\in\mathbb{R}\right\} \)}{\( \left\{\left[x;-\frac13\right],x\in\mathbb{R}\right\} \)}{}{}}
\End

\Data\ID{1003083001}
\Updated{2020-07-15 03:07:12}
\Level{A}
\SubArea{Systems of equations and inequalities}
\LANGen{
\Question{
Identify which of the following systems of equations has infinitely many solutions.}
{\( \begin{aligned} \frac13x-4y&=2\\ -\frac{x}4+3y&=-\frac32 \end{aligned} \)}{\( \begin{aligned} \frac13 x-4y&=2 \\ -x+12y&=6 \end{aligned} \)}{\( \begin{aligned} \frac13 x-4y&=2 \\ \frac x4-6y&=6 \end{aligned} \)}{\( \begin{aligned} \frac13 x-4y&=2 \\ \frac x3-4y&=0 \end{aligned} \)}{}{}}
\LANGpl{
\Question{
Określ, który z poniższych układów równań ma nieskończenie wiele rozwiązań.}
{\( \begin{aligned} \frac13x-4y&=2\\ -\frac{x}4+3y&=-\frac32 \end{aligned} \)}{\( \begin{aligned} \frac13 x-4y&=2 \\ -x+12y&=6 \end{aligned} \)}{\( \begin{aligned} \frac13 x-4y&=2 \\ \frac x4-6y&=6 \end{aligned} \)}{\( \begin{aligned} \frac13 x-4y&=2 \\ \frac x3-4y&=0 \end{aligned} \)}{}{}}
\LANGcs{
\Question{
Která z následujících soustav má nekonečně mnoho kořenů?}
{\( \begin{aligned} \frac13x-4y&=2\\ -\frac{x}4+3y&=-\frac32 \end{aligned} \)}{\( \begin{aligned} \frac13 x-4y&=2 \\ -x+12y&=6 \end{aligned} \)}{\( \begin{aligned} \frac13 x-4y&=2 \\ \frac x4-6y&=6 \end{aligned} \)}{\( \begin{aligned} \frac13 x-4y&=2 \\ \frac x3-4y&=0 \end{aligned} \)}{}{}}
\LANGsk{
\Question{
Ktorá z nasledujúcich sústav má nekonečne veľa riešení?}
{\( \begin{aligned} \frac13x-4y&=2\\ -\frac{x}4+3y&=-\frac32 \end{aligned} \)}{\( \begin{aligned} \frac13 x-4y&=2 \\ -x+12y&=6 \end{aligned} \)}{\( \begin{aligned} \frac13 x-4y&=2 \\ \frac x4-6y&=6 \end{aligned} \)}{\( \begin{aligned} \frac13 x-4y&=2 \\ \frac x3-4y&=0 \end{aligned} \)}{}{}}
\LANGes{
\Question{
Identifica cuál de los siguientes sistemas de ecuaciones tiene infinitas soluciones.}
{\( \begin{aligned} \frac13x-4y&=2\\ -\frac{x}4+3y&=-\frac32 \end{aligned} \)}{\( \begin{aligned} \frac13 x-4y&=2 \\ -x+12y&=6 \end{aligned} \)}{\( \begin{aligned} \frac13 x-4y&=2 \\ \frac x4-6y&=6 \end{aligned} \)}{\( \begin{aligned} \frac13 x-4y&=2 \\ \frac x3-4y&=0 \end{aligned} \)}{}{}}
\End

\Data\ID{1003083002}
\Updated{2020-07-15 03:07:12}
\Level{A}
\SubArea{Systems of equations and inequalities}
\LANGen{
\Question{
Identify which of the sets is not the solution set of the following system of equations.
\[ \begin{aligned} \frac12 x-y&=3 \\ \frac x3 - \frac23 y &=2 \end{aligned} \]}
{\( \left\{\left[6+2y;\frac{x-6}2\right]\colon  x\in\mathbb{R}\text{, }y\in\mathbb{R}\right\} \)}{\( \left\{\left[x; \frac{x-6}2\right]\colon  x\in\mathbb{R}\right\} \)}{\( \left\{\left[6+2y;y\right]\colon  y\in\mathbb{R}\right\} \)}{\( \left\{\left[2t;t-3\right]\colon  t\in\mathbb{R}\right\} \)}{}{}}
\LANGpl{
\Question{
Określ, który ze zbiorów nie jest rozwiązaniem podanego układu równań.
\[ \begin{aligned} \frac12 x-y&=3 \\ \frac x3 - \frac23 y &=2 \end{aligned} \]}
{\( \left\{\left[6+2y;\frac{x-6}2\right]\colon  x\in\mathbb{R}\text{, }y\in\mathbb{R}\right\} \)}{\( \left\{\left[x; \frac{x-6}2\right]\colon  x\in\mathbb{R}\right\} \)}{\( \left\{\left[6+2y;y\right]\colon  y\in\mathbb{R}\right\} \)}{\( \left\{\left[2t;t-3\right]\colon  t\in\mathbb{R}\right\} \)}{}{}}
\LANGcs{
\Question{
Z následujících množin vyberte tu, která nepředstavuje množinu kořenů následující soustavy rovnic.
\[ \begin{aligned} \frac12 x-y&=3 \\ \frac x3 - \frac23 y &=2 \end{aligned} \]}
{\( \left\{\left[6+2y;\frac{x-6}2\right]\colon  x\in\mathbb{R}\text{, }y\in\mathbb{R}\right\} \)}{\( \left\{\left[x; \frac{x-6}2\right]\colon  x\in\mathbb{R}\right\} \)}{\( \left\{\left[6+2y;y\right]\colon  y\in\mathbb{R}\right\} \)}{\( \left\{\left[2t;t-3\right]\colon  t\in\mathbb{R}\right\} \)}{}{}}
\LANGsk{
\Question{
Z nasledujúcich množín vyberte tú, ktorá nepredstavuje množinu riešení nasledujúcej sústavy rovníc.
\[ \begin{aligned} \frac12 x-y&=3 \\ \frac x3 - \frac23 y &=2 \end{aligned} \]}
{\( \left\{\left[6+2y;\frac{x-6}2\right]\colon  x\in\mathbb{R}\text{, }y\in\mathbb{R}\right\} \)}{\( \left\{\left[x; \frac{x-6}2\right]\colon  x\in\mathbb{R}\right\} \)}{\( \left\{\left[6+2y;y\right]\colon  y\in\mathbb{R}\right\} \)}{\( \left\{\left[2t;t-3\right]\colon  t\in\mathbb{R}\right\} \)}{}{}}
\LANGes{
\Question{
Identifica cuál de los conjuntos no es el conjunto de soluciones del siguiente sistema.
\[ \begin{aligned} \frac12 x-y&=3 \\ \frac x3 - \frac23 y &=2 \end{aligned} \]}
{\( \left\{\left[6+2y;\frac{x-6}2\right]\colon  x\in\mathbb{R}\text{, }y\in\mathbb{R}\right\} \)}{\( \left\{\left[x; \frac{x-6}2\right]\colon  x\in\mathbb{R}\right\} \)}{\( \left\{\left[6+2y;y\right]\colon  y\in\mathbb{R}\right\} \)}{\( \left\{\left[2t;t-3\right]\colon  t\in\mathbb{R}\right\} \)}{}{}}
\End

\Data\ID{9000007206}
\Updated{2020-07-15 03:07:34}
\Level{A}
\SubArea{Systems of equations and inequalities}
\LANGen{
\Question{
Consider the linear system 
\[ 
\begin{aligned}2x - 3y - 12& = 0,&
\\\text{???}\quad & = 0.
\\ \end{aligned}
\]
In the following list identify a missing second equation if we know that the system
does not have a solution.}
{\(- 6x + 9y - 9 = 0\)}{\(2x + 3y - 6 = 0\)}{\(- 4x + 6y + 24 = 0\)}{\(x + 2y - 12 = 0\)}{}{}}
\LANGpl{
\Question{
Rozważ układ liniowy
\[ 
\begin{aligned}2x - 3y - 12& = 0,&
\\\text{???}\quad & = 0.
\\ \end{aligned}
\]
Z poniższej listy wybierz brakujące równanie, jeżeli wiemy, że dany układ równań jest sprzeczny.}
{\(- 6x + 9y - 9 = 0\)}{\(2x + 3y - 6 = 0\)}{\(- 4x + 6y + 24 = 0\)}{\(x + 2y - 12 = 0\)}{}{}}
\LANGcs{
\Question{
Uvažujme soustavu dvou rovnic \[ 
\begin{aligned}2x - 3y - 12& = 0,&
\\\text{???}\quad & = 0.
\\ \end{aligned}
\]
Z nabízených možností vyberte chybějící druhou rovnici soustavy tak, aby
výsledná soustava neměla řešení.}
{\(- 6x + 9y - 9 = 0\)}{\(2x + 3y - 6 = 0\)}{\(- 4x + 6y + 24 = 0\)}{\(x + 2y - 12 = 0\)}{}{}}
\LANGsk{
\Question{
Uvažujme sústavu 2 rovníc \[ 
\begin{aligned}2x - 3y - 12& = 0,&
\\\text{???}\quad & = 0.
\\ \end{aligned}
\]
Z ponúknutých možností vyberte chýbajúcu druhú rovnicu sústavy tak, aby výsledná sústava nemala riešenie.}
{\(- 6x + 9y - 9 = 0\)}{\(2x + 3y - 6 = 0\)}{\(- 4x + 6y + 24 = 0\)}{\(x + 2y - 12 = 0\)}{}{}}
\LANGes{
\Question{
Dado el sistema lineal: 
\[ 
\begin{aligned}2x - 3y - 12& = 0,&
\\\text{???}\quad & = 0.
\\ \end{aligned}
\]
Identifica la segunda ecuación que falta sabiendo que el sistema no tiene ninguna solución.}
{\(- 6x + 9y - 9 = 0\)}{\(2x + 3y - 6 = 0\)}{\(- 4x + 6y + 24 = 0\)}{\(x + 2y - 12 = 0\)}{}{}}
\End

\Data\ID{9000023901}
\Updated{2020-08-23 08:08:38}
\Level{A}
\SubArea{Systems of equations and inequalities}
\LANGen{
\Question{
Solve the following system and write the solution as the ordered pair
\([x,y]\).
\[\begin{aligned}
 x + y & = -1 & &
 \\x - y & = 5 & &
 \end{aligned}\]}
{\([2;-3]\)}{\([-2;1]\)}{\([3;-2]\)}{\([-3;2]\)}{}{}}
\LANGpl{
\Question{
Rozwiąż następujący układ równań i zapisz rozwiązanie jako parę 
\([x,y]\).
\[\begin{aligned}
 x + y & = -1 & &
 \\x - y & = 5 & &
 \end{aligned}\]}
{\([2;-3]\)}{\([-2;1]\)}{\([3;-2]\)}{\([-3;2]\)}{}{}}
\LANGcs{
\Question{
Vyřešte následující soustavu rovnic a řešení zapište jako uspořádanou dvojici $[x;y]$.
\begin{align*}
x + y &= -1 \\
x - y &= 5
\end{align*}
Řešením této soustavy je uspořádaná dvojice:}
{\([2;-3]\)}{\([-2;1]\)}{\([3;-2]\)}{\([-3;2]\)}{}{}}
\LANGsk{
\Question{
Vyriešte danú sústavu rovníc. Riešením tejto sústavy je usporiadaná dvojica
\([x,y]\).
\[\begin{aligned}
 x + y & = -1 & &
 \\x - y & = 5 & &
 \end{aligned}\]}
{\([2;-3]\)}{\([-2;1]\)}{\([3;-2]\)}{\([-3;2]\)}{}{}}
\LANGes{
\Question{
Resuelve el siguiente sistema y escribe la solución como par ordenado \([x,y]\).
\[\begin{aligned}
 x + y & = -1 & &
 \\x - y & = 5 & &
 \end{aligned}\]}
{\([2;-3]\)}{\([-2;1]\)}{\([3;-2]\)}{\([-3;2]\)}{}{}}
\End

\Data\ID{9000023902}
\Updated{2020-08-23 08:08:59}
\Level{A}
\SubArea{Systems of equations and inequalities}
\LANGen{
\Question{
Solve the following system and write the solution as the ordered pair
\([x,y]\).
\[\begin{aligned}
 2x + y & = 2 & &
 \\x + 2y & = 7 & &
 \end{aligned}\]}
{\([-1;4]\)}{\([2;-2]\)}{\([1;3]\)}{\([3;-4]\)}{}{}}
\LANGpl{
\Question{
Rozwiąż następujący układ i zapisz rozwiązanie jako parę
\([x,y]\).
\[\begin{aligned}
 2x + y & = 2 & &
 \\x + 2y & = 7 & &
 \end{aligned}\]}
{\([-1;4]\)}{\([2;-2]\)}{\([1;3]\)}{\([3;-4]\)}{}{}}
\LANGcs{
\Question{
Vyřešte následující soustavu rovnic a řešení zapište jako uspořádanou dvojici $[x;y]$.
\begin{align*}
2x + y &= 2 \\
 x + 2y &= 7
\end{align*}}
{\([-1;4]\)}{\([2;-2]\)}{\([1;3]\)}{\([3;-4]\)}{}{}}
\LANGsk{
\Question{
Vyriešte danú sústavu rovníc. Riešením tejto sústavy je usporiadaná dvojica 
\([x,y]\).
\[\begin{aligned}
 2x + y & = 2 & &
 \\x + 2y & = 7 & &
 \end{aligned}\]}
{\([-1;4]\)}{\([2;-2]\)}{\([1;3]\)}{\([3;-4]\)}{}{}}
\LANGes{
\Question{
Resuelve el siguiente sistema y escribe la solución como par ordenado
\([x,y]\).
\[\begin{aligned}
 2x + y & = 2 & &
 \\x + 2y & = 7 & &
 \end{aligned}\]}
{\([-1;4]\)}{\([2;-2]\)}{\([1;3]\)}{\([3;-4]\)}{}{}}
\End

\Data\ID{9000023903}
\Updated{2020-08-23 08:08:42}
\Level{A}
\SubArea{Systems of equations and inequalities}
\LANGen{
\Question{
Let \([x;y]\) 
be the solution of the system
\[\begin{aligned}
 2x + 3y & = -2, & &
 \\3x - 2y & = 10. & &
 \end{aligned}\]
Find \(5x + y\).}
{\(8\)}{\(- 12\)}{\(- 8\)}{\(12\)}{}{}}
\LANGpl{
\Question{
Niech \([x;y]\) 
będzie rozwiązaniem układu
\[\begin{aligned}
 2x + 3y & = -2, & &
 \\3x - 2y & = 10. & &
 \end{aligned}\]
Oblicz \(5x + y\).}
{\(8\)}{\(- 12\)}{\(- 8\)}{\(12\)}{}{}}
\LANGcs{
\Question{
Vyřešte následující soustavu rovnic a řešení zapište jako uspořádanou dvojici $[x;y]$.
\begin{align*}
2x + 3y &= -2 \\ 
3x - 2y &= 10
\end{align*}
Výraz $5x+y$ je roven:}
{\(8\)}{\(- 12\)}{\(- 8\)}{\(12\)}{}{}}
\LANGsk{
\Question{
Vyriešte danú sústavu rovníc. Riešením tejto sústavy je usporiadaná dvojica \([x;y]\).
\[\begin{aligned}
 2x + 3y & = -2 & &
 \\3x - 2y & = 10 & &
 \end{aligned}\]
Výraz \( 5x+y \) sa rovná:}
{\(8\)}{\(- 12\)}{\(- 8\)}{\(12\)}{}{}}
\LANGes{
\Question{
Sea \([x;y]\)  la solución del sistema \[\begin{aligned}
 2x + 3y & = -2, & &
 \\3x - 2y & = 10. & &
 \end{aligned}\]
Halla \(5x + y\).}
{\(8\)}{\(- 12\)}{\(- 8\)}{\(12\)}{}{}}
\End

\Data\ID{9000023904}
\Updated{2020-08-23 08:08:02}
\Level{A}
\SubArea{Systems of equations and inequalities}
\LANGen{
\Question{
Let \([x;y]\) 
be the solution of the system
\[\begin{aligned}
 4x - 3y & = -3, & &
 \\x + 2y & = 13. & &
 \end{aligned}\]
Find \(2x - 7y\).}
{\(- 29\)}{\(- 41\)}{\(- 11\)}{\(31\)}{}{}}
\LANGpl{
\Question{
Niech \([x;y]\) 
będzie rozwiązaniem układu
\[\begin{aligned}
 4x - 3y & = -3, & &
 \\x + 2y & = 13. & &
 \end{aligned}\]
Oblicz \(2x - 7y\).}
{\(- 29\)}{\(- 41\)}{\(- 11\)}{\(31\)}{}{}}
\LANGcs{
\Question{
Vyřešte následující soustavu rovnic a řešení zapište jako uspořádanou dvojici $[x;y]$.
\begin{align*}
4x - 3y &= -3 \\ 
x + 2y &= 13
\end{align*}
Výraz
\(2x - 7y\) je
roven:}
{\(- 29\)}{\(- 41\)}{\(- 11\)}{\(31\)}{}{}}
\LANGsk{
\Question{
Vyriešte danú sústavu rovníc. Riešením tejto sústavy je usporiadaná dvojica \([x;y]\).
\[\begin{aligned}
 4x - 3y & = -3 & &
 \\x + 2y & = 13 & &
 \end{aligned}\]
Výraz  \(2x - 7y\) sa rovná:}
{\(- 29\)}{\(- 41\)}{\(- 11\)}{\(31\)}{}{}}
\LANGes{
\Question{
Sea \([x;y]\)  la solución del sistema
\[\begin{aligned}
 4x - 3y & = -3, & &
 \\x + 2y & = 13. & &
 \end{aligned}\]
Halla \(2x - 7y\).}
{\(- 29\)}{\(- 41\)}{\(- 11\)}{\(31\)}{}{}}
\End

\Data\ID{9000023905}
\Updated{2020-08-23 08:08:19}
\Level{A}
\SubArea{Systems of equations and inequalities}
\LANGen{
\Question{
Let \([x;y]\) 
be the solution of the system
\[\begin{aligned}
 2x + 5y & = 7, & &
 \\ - 4x - 3y & = 7. & &
 \end{aligned}\]
In the following list identify a true statement.}
{\(x^{2} + y^{2} = 25\)}{\(x^{2} + y^{2} = 7\)}{\(x^{2} - y^{2} = -7\)}{\(y^{2} - x^{2} = 7\)}{}{}}
\LANGpl{
\Question{
Niech \([x;y]\) 
będzie rozwiązaniem układu 
\[\begin{aligned}
 2x + 5y & = 7, & &
 \\ - 4x - 3y & = 7. & &
 \end{aligned}\]
Która z poniższych odpowiedzi jest poprawna?}
{\(x^{2} + y^{2} = 25\)}{\(x^{2} + y^{2} = 7\)}{\(x^{2} - y^{2} = -7\)}{\(y^{2} - x^{2} = 7\)}{}{}}
\LANGcs{
\Question{
Vyřešte následující soustavu rovnic a řešení zapište jako uspořádanou dvojici $[x;y]$.
\begin{align*}
2x + 5y &= 7 \\
-4x - 3y &= 7 \end{align*}
Které z následujících tvrzení je správné?}
{\(x^{2} + y^{2} = 25\)}{\(x^{2} + y^{2} = 7\)}{\(x^{2} - y^{2} = -7\)}{\(y^{2} - x^{2} = 7\)}{}{}}
\LANGsk{
\Question{
Vyriešte danú sústavu rovníc. Riešením tejto sústavy je usporiadaná dvojica \([x;y]\).
\[\begin{aligned}
 2x + 5y & = 7 & &
 \\ - 4x - 3y & = 7 & &
 \end{aligned}\]
Ktoré z nasledujúcich tvrdení je správne?}
{\(x^{2} + y^{2} = 25\)}{\(x^{2} + y^{2} = 7\)}{\(x^{2} - y^{2} = -7\)}{\(y^{2} - x^{2} = 7\)}{}{}}
\LANGes{
\Question{
Sea \([x;y]\) la solución del sistema
\[\begin{aligned}
 2x + 5y & = 7, & &
 \\ - 4x - 3y & = 7. & &
 \end{aligned}\]
Identifica la proposición lógica.}
{\(x^{2} + y^{2} = 25\)}{\(x^{2} + y^{2} = 7\)}{\(x^{2} - y^{2} = -7\)}{\(y^{2} - x^{2} = 7\)}{}{}}
\End

\Data\ID{9000023906}
\Updated{2020-08-23 08:08:33}
\Level{A}
\SubArea{Systems of equations and inequalities}
\LANGen{
\Question{
Let \([x;y]\) 
be the solution of the system
\[\begin{aligned}
 2x + 3y & = 4, & &
 \\4x + 6y & = 9. & &
 \end{aligned}\]
In the following list identify a true statement.}
{The system does not have any solution.}{\(x < y\)}{\(x > y\)}{\(x = y\)}{}{}}
\LANGpl{
\Question{
Niech \([x;y]\) 
będzie rozwiązaniem układu
\[\begin{aligned}
 2x + 3y & = 4, & &
 \\4x + 6y & = 9. & &
 \end{aligned}\]
Które z poniższych zdań jest prawdziwe?}
{Układ nie ma rozwiązania.}{\(x < y\)}{\(x > y\)}{\(x = y\)}{}{}}
\LANGcs{
\Question{
Vyřešte následující soustavu rovnic a řešení zapište jako uspořádanou dvojici $[x;y]$.
\begin{align*}
2x + 3y &= 4 \\
4x + 6y &= 9\end{align*}
Které z následujících tvrzení je správné?}
{Daná soustava nemá řešení.}{\(x < y\)}{\(x > y\)}{\(x = y\)}{}{}}
\LANGsk{
\Question{
Vyriešte danú sústavu rovníc. Riešením tejto sústavy je usporiadaná dvojica \([x;y]\).
\[\begin{aligned}
 2x + 3y & = 4 & &
 \\4x + 6y & = 9 & &
 \end{aligned}\]
Ktoré z nasledujúcich tvrdení je správne?}
{Daná sústava nemá riešenie.}{\(x < y\)}{\(x > y\)}{\(x = y\)}{}{}}
\LANGes{
\Question{
Sea \([x;y]\) la solución del sistema
\[\begin{aligned}
 2x + 3y & = 4, & &
 \\4x + 6y & = 9. & &
 \end{aligned}\]
Identifica la proposición lógica.}
{El sistema no tiene solución.}{\(x < y\)}{\(x > y\)}{\(x = y\)}{}{}}
\End

\Data\ID{9000023907}
\Updated{2020-08-23 08:08:01}
\Level{A}
\SubArea{Systems of equations and inequalities}
\LANGen{
\Question{
Let \([x;y]\) 
be the solution of the system
\[\begin{aligned}
 - x + 2y & = 6, & &
 \\2x + 3y & = 2. & &
 \end{aligned}\]
In the following list identify a true statement.}
{\(|x| = |y|\)}{\(|x| < |y|\)}{\(|x| > |y|\)}{The system does not have any solution.}{}{}}
\LANGpl{
\Question{
Niech \([x;y]\) 
będzie rozwiązaniem układu 
\[\begin{aligned}
 - x + 2y & = 6, & &
 \\2x + 3y & = 2. & &
 \end{aligned}\]
Która z poniższych odpowiedzi jest poprawna?}
{\(|x| = |y|\)}{\(|x| < |y|\)}{\(|x| > |y|\)}{Układ nie ma rozwiązania.}{}{}}
\LANGcs{
\Question{
Vyřešte následující soustavu rovnic a řešení zapište jako uspořádanou dvojici $[x;y]$.
\begin{align*}
- x + 2y &= 6 \\
2x + 3y &= 2 \end{align*}
Které z následujících tvrzení je správné?}
{\(|x| = |y|\)}{\(|x| < |y|\)}{\(|x| > |y|\)}{Daná soustava nemá řešení.}{}{}}
\LANGsk{
\Question{
Vyriešte danú sústavu rovníc. Riešením tejto sústavy je usporiadaná dvojica \([x;y]\).
\[\begin{aligned}
 - x + 2y & = 6 & &
 \\2x + 3y & = 2 & &
 \end{aligned}\]
Ktoré z nasledujúcich tvrdení je správne?}
{\(|x| = |y|\)}{\(|x| < |y|\)}{\(|x| > |y|\)}{Daná sústava nemá riešenie.}{}{}}
\LANGes{
\Question{
Sea \([x;y]\) la solución del sistema
\[\begin{aligned}
 - x + 2y & = 6, & &
 \\2x + 3y & = 2. & &
 \end{aligned}\]
Identifica la proposición lógica.}
{\(|x| = |y|\)}{\(|x| < |y|\)}{\(|x| > |y|\)}{El sistema no tiene solución.}{}{}}
\End

\Data\ID{9000023908}
\Updated{2020-08-23 08:08:16}
\Level{A}
\SubArea{Systems of equations and inequalities}
\LANGen{
\Question{
Let \([x;y]\) 
be the solution of the system
\[\begin{aligned}
 2x - y & = -1, & &
 \\4x - y & = 1. & &
 \end{aligned}\]
In the following list identify a true statement.}
{\(y\) is
a prime number.}{\(x\) is
a prime number.}{\(x + y\) is
a prime number.}{\(x - y\) is
a prime number.}{}{}}
\LANGpl{
\Question{
Niech \([x;y]\) 
będzie rozwiązaniem układu
\[\begin{aligned}
 2x - y & = -1, & &
 \\4x - y & = 1. & &
 \end{aligned}\]
Która z poniższych odpowiedzi jest poprawna?}
{\(y\) jest 
liczbą pierwszą.}{\(x\) jest 
liczbą pierwszą.}{\(x + y\) jest 
liczbą pierwszą.}{\(x - y\) jest 
liczbą pierwszą.}{}{}}
\LANGcs{
\Question{
Vyřešte následující soustavu rovnic a řešení zapište jako uspořádanou dvojici $[x;y]$.
\begin{align*}
2x - y &= -1 \\
4x - y &= 1\end{align*}
Které z následujících tvrzení je správné?}
{\(y\) je
prvočíslo.}{\(x\) je
prvočíslo.}{\(x + y\) je
prvočíslo.}{\(x - y\) je
prvočíslo.}{}{}}
\LANGsk{
\Question{
Vyriešte danú sústavu rovníc. Riešením tejto sústavy je usporiadaná dvojica \([x;y]\).
\[\begin{aligned}
 2x - y & = -1 & &
 \\4x - y & = 1 & &
 \end{aligned}\]
Ktoré z nasledujúcich tvrdení je správne?}
{\(y\) je prvočíslo.}{\(x\) je prvočíslo.}{\(x + y\) je prvočíslo.}{\(x - y\) je prvočíslo.}{}{}}
\LANGes{
\Question{
Sea \([x;y]\) la solución del sistema
\[\begin{aligned}
 2x - y & = -1, & &
 \\4x - y & = 1. & &
 \end{aligned}\]
Identifica la proposición lógica.}
{\(y\) es un número primo.}{\(x\) es un número primo.}{\(x + y\) es un número primo.}{\(x - y\) es un número primo.}{}{}}
\End

\Data\ID{9000023909}
\Updated{2020-08-23 08:08:33}
\Level{A}
\SubArea{Systems of equations and inequalities}
\LANGen{
\Question{
Let \([x;y]\) 
be the solution of the system
\[\begin{aligned}
 2x + 3y & = 0, & &
 \\3x + 2y & = 5. & &
 \end{aligned}\]
In the following list identify a true statement.}
{\(- 3\leq y\leq - 1\)}{\(- 1\leq x\leq 2\)}{\(- 2\leq x\leq 2\)}{\(2\leq y\leq 4\)}{}{}}
\LANGpl{
\Question{
Niech \([x;y]\) 
będzie rozwiązaniem układu
\[\begin{aligned}
 2x + 3y & = 0, & &
 \\3x + 2y & = 5. & &
 \end{aligned}\]
Która z poniższych odpowiedzi jest poprawna?}
{\(- 3\leq y\leq - 1\)}{\(- 1\leq x\leq 2\)}{\(- 2\leq x \leq 2\)}{\(2\leq y\leq 4\)}{}{}}
\LANGcs{
\Question{
Vyřešte následující soustavu rovnic a řešení zapište jako uspořádanou dvojici $[x;y]$.
\begin{align*}
2x + 3y &= 0 \\
3x + 2y &= 5\end{align*}
Které z následujících tvrzení je správné?}
{\(- 3\leq y\leq - 1\)}{\(- 1\leq x\leq 2\)}{\(- 2\leq x \leq 2\)}{\(2\leq y\leq 4\)}{}{}}
\LANGsk{
\Question{
Vyriešte danú sústavu rovníc. Riešením tejto sústavy je usporiadaná dvojica \([x;y]\).
\[\begin{aligned}
 2x + 3y & = 0 & &
 \\3x + 2y & = 5 & &
 \end{aligned}\]
Ktoré z nasledujúcich tvrdení je správne?}
{\(- 3\leq y\leq - 1\)}{\(- 1\leq x\leq 2\)}{\(- 2\leq x \leq 2\)}{\(2\leq y\leq 4\)}{}{}}
\LANGes{
\Question{
Sea \([x;y]\) la solución del sistema
\[\begin{aligned}
 2x + 3y & = 0, & &
 \\3x + 2y & = 5. & &
 \end{aligned}\]
Identifica la proposición lógica.}
{\(- 3\leq y\leq - 1\)}{\(- 1\leq x\leq 2\)}{\(- 2\leq x\leq 2\)}{\(2\leq y\leq 4\)}{}{}}
\End

\Data\ID{9000023910}
\Updated{2020-08-23 08:08:56}
\Level{A}
\SubArea{Systems of equations and inequalities}
\LANGen{
\Question{
Let \([x;y]\) 
be the solution of the system
\[\begin{aligned}
 3x - y & = 1, & &
 \\2x - y & = -1. & &
 \end{aligned}\]
In the following list identify a true statement.}
{\(x\) 
divides \(6\).}{\(x\) 
divides \(3\).}{\(y\) 
divides \(4\).}{\(y\) 
divides \(6\).}{}{}}
\LANGpl{
\Question{
Niech \([x;y]\) 
będzie rozwiązaniem układu 
\[\begin{aligned}
 3x - y & = 1, & &
 \\2x - y & = -1. & &
 \end{aligned}\]
Która z poniższych odpowiedzi jest poprawna?}
{\(x\) 
jest podzielne przez \(6\).}{\(x\) 
jest podzielne przez \(3\).}{\(y\) 
jest podzielne przez \(4\).}{\(y\) 
jest podzielne przez \(6\).}{}{}}
\LANGcs{
\Question{
Vyřešte následující soustavu rovnic a řešení zapište jako uspořádanou dvojici $[x;y]$.
\begin{align*}
3x - y &= 1 \\
2x - y &= -1\end{align*}
Které z následujících tvrzení je správné?}
{\(x\) je dělitelem
čísla \(6\).}{\(x\) je dělitelem
čísla \(3\).}{\(y\) je dělitelem
čísla \(4\).}{\(y\) je dělitelem
čísla \(6\).}{}{}}
\LANGsk{
\Question{
Vyriešte danú sústavu rovníc. Riešením tejto sústavy je usporiadaná dvojica \([x;y]\).
\[\begin{aligned}
 3x - y & = 1, & &
 \\2x - y & = -1. & &
 \end{aligned}\]
Ktoré z nasledujúcich tvrdení je správne?}
{\(x\) 
je deliteľom \(6\).}{\(x\) 
je deliteľom \(3\).}{\(y\) 
je deliteľom \(4\).}{\(y\) 
je deliteľom \(6\).}{}{}}
\LANGes{
\Question{
Sea \([x;y]\) la solución del sistema
\[\begin{aligned}
 3x - y & = 1, & &
 \\2x - y & = -1. & &
 \end{aligned}\]
Identifica la proposición lógica.}
{\(x\) 
es divisor del número \(6\).}{\(x\) 
es divisor del número \(3\).}{\(y\) 
es divisor del número \(4\).}{\(y\) 
es divisor del número \(6\).}{}{}}
\End

\Data\ID{1003020101}
\Updated{2020-07-15 03:07:12}
\Level{B}
\SubArea{Systems of equations and inequalities}
\LANGen{
\Question{
How many solutions has the inequality
\[ -4y\geq4x-8 \]
in \( \mathbb{N}\times\mathbb{N} \)? (Hint: Use the graphic solution of the inequality.)}
{\( 1 \)}{\( 4 \)}{infinity}{none}{}{}}
\LANGpl{
\Question{
Ile rozwiązań ma nierówność
\[ -4y\geq4x-8 \]
w \( \mathbb{N}\times\mathbb{N} \)? (Wskazówka: Skorzystaj z graficznego rozwiązania tej nierówności.)}
{\( 1 \)}{\( 4 \)}{nieskończoność}{nie ma rozwiązania}{}{}}
\LANGcs{
\Question{
Kolik řešení má nerovnice
\[ -4y\geq4x-8 \]
v \( \mathbb{N}\times\mathbb{N} \)? (Nápověda: Využijte grafické řešení nerovnice.)}
{\( 1 \)}{\( 4 \)}{nekonečně mnoho}{žádné}{}{}}
\LANGsk{
\Question{
Koľko riešení má nerovnica
\[ -4y\geq4x-8 \]
v \( \mathbb{N}\times\mathbb{N} \)? (Pomôcka: Využite grafické riešenie nerovnice.)}
{\( 1 \)}{\( 4 \)}{nekonečne veľa}{žiadne}{}{}}
\LANGes{
\Question{
¿Cuántas soluciones tiene la siguiente inecuación
\[ -4y\geq4x-8 \]
en \( \mathbb{N}\times\mathbb{N} \)? (Sugerencia: Usa la solución gráfica de la inecuación.)}
{\( 1 \)}{\( 4 \)}{infinitas}{ninguna}{}{}}
\End

\Data\ID{1003020103}
\Updated{2020-07-15 03:07:12}
\Level{B}
\SubArea{Systems of equations and inequalities}
\LANGen{
\Question{
How many solutions has the inequality
\[ \frac{3x-7y}2 < x-4y-\frac12 \]
in \( \mathbb{N}\times\mathbb{N} \)? (Hint: Use the graphic solution of the inequality.)}
{none}{infinity}{\( 4 \)}{\( 2 \)}{}{}}
\LANGpl{
\Question{
Ile rozwiązań ma nierówność
\[ \frac{3x-7y}2 < x-4y-\frac12 \]
w \( \mathbb{N}\times\mathbb{N} \)? (Wskazówka: Skorzystaj z graficznego rozwiązania tej nierówności.)}
{nie ma rozwiązania}{nieskończoność}{\( 4 \)}{\( 2 \)}{}{}}
\LANGcs{
\Question{
Kolik řešení má nerovnice
\[ \frac{3x-7y}2 < x-4y-\frac12 \]
v \( \mathbb{N}\times\mathbb{N} \)? (Nápověda: Využijte grafické řešení nerovnice.)}
{žádné}{nekonečně mnoho}{\( 4 \)}{\( 2 \)}{}{}}
\LANGsk{
\Question{
Koľko riešení má nerovnica
\[ \frac{3x-7y}2 < x-4y-\frac12 \]
v \( \mathbb{N}\times\mathbb{N} \)? (Pomôcka: Využite grafické riešenie nerovnice.)}
{žiadne}{nekonečne veľa}{\( 4 \)}{\( 2 \)}{}{}}
\LANGes{
\Question{
¿Cuántas soluciones tiene la siguiente inecuación
\[ \frac{3x-7y}2 < x-4y-\frac12 \]
en \( \mathbb{N}\times\mathbb{N} \)? (Sugerencia: Usa la solución gráfica de la inecuación.)}
{ninguna}{infinitas}{\( 4 \)}{\( 2 \)}{}{}}
\End

\Data\ID{1003083003}
\Updated{2020-07-15 03:07:12}
\Level{B}
\SubArea{Systems of equations and inequalities}
\LANGen{
\Question{
Find the solution set of the following system of equations.
\[ \begin{aligned}\frac23 x-\frac12y&=1 \\ -2x+\frac32y&=-3 \end{aligned} \]}
{\( \left\{\left[x; \frac{4x-6}3\right]\colon  x\in\mathbb{R}\right\} \)}{\( \left\{\left[x; y\right]\colon  x\in\mathbb{R}\text{, } y\in\mathbb{R}\right\} \)}{\( \emptyset \)}{\( \left\{[0; -2]\right\} \)}{}{}}
\LANGpl{
\Question{
Znajdź zbiór rozwiązań podanego układu równań.
\[ \begin{aligned}\frac23 x-\frac12y&=1 \\ -2x+\frac32y&=-3 \end{aligned} \]}
{\( \left\{\left[x; \frac{4x-6}3\right]\colon  x\in\mathbb{R}\right\} \)}{\( \left\{\left[x; y\right]\colon  x\in\mathbb{R}\text{, } y\in\mathbb{R}\right\} \)}{\( \emptyset \)}{\( \left\{[0; -2]\right\} \)}{}{}}
\LANGcs{
\Question{
Najděte množinu všech řešení následující soustavy rovnic.
\[ \begin{aligned}\frac23 x-\frac12y&=1 \\ -2x+\frac32y&=-3 \end{aligned} \]}
{\( \left\{\left[x; \frac{4x-6}3\right]\colon  x\in\mathbb{R}\right\} \)}{\( \left\{\left[x; y\right]\colon  x\in\mathbb{R}\text{, } y\in\mathbb{R}\right\} \)}{\( \emptyset \)}{\( \left\{[0; -2]\right\} \)}{}{}}
\LANGsk{
\Question{
Nájdite množinu všetkých riešení nasledujúcej sústavy rovníc.
\[ \begin{aligned}\frac23 x-\frac12y&=1 \\ -2x+\frac32y&=-3 \end{aligned} \]}
{\( \left\{\left[x; \frac{4x-6}3\right]\colon  x\in\mathbb{R}\right\} \)}{\( \left\{\left[x; y\right]\colon  x\in\mathbb{R}\text{, } y\in\mathbb{R}\right\} \)}{\( \emptyset \)}{\( \left\{[0; -2]\right\} \)}{}{}}
\LANGes{
\Question{
Calcula el conjunto de soluciones del siguiente sistema de ecuaciones.
\[ \begin{aligned}\frac23 x-\frac12y&=1 \\ -2x+\frac32y&=-3 \end{aligned} \]}
{\( \left\{\left[x; \frac{4x-6}3\right]\colon  x\in\mathbb{R}\right\} \)}{\( \left\{\left[x; y\right]\colon  x\in\mathbb{R}\text{, } y\in\mathbb{R}\right\} \)}{\( \emptyset \)}{\( \left\{[0; -2]\right\} \)}{}{}}
\End

\Data\ID{1103020102}
\Updated{2020-07-15 03:07:12}
\Level{B}
\SubArea{Systems of equations and inequalities}
\IMG{1103020102.pdf}
\LANGen{
\Question{
How many solutions has the inequality
\[ y\leq-\frac12x+2 \]
in \( \mathbb{N}\times\mathbb{N} \)?}
{\( 2 \)}{\( 4 \)}{infinity}{none}{}{}}
\LANGpl{
\Question{
Ile rozwiązań ma podana nierówność
\[ y\leq-\frac12x+2 \]
w \( \mathbb{N}\times\mathbb{N} \)?}
{\( 2 \)}{\( 4 \)}{nieskończenie wiele}{nie ma rozwiązania}{}{}}
\LANGcs{
\Question{
Kolik řešení má nerovnice
\[ y\leq-\frac12x+2 \]
v \( \mathbb{N}\times\mathbb{N} \)?}
{\( 2 \)}{\( 4 \)}{nekonečně mnoho}{žádné}{}{}}
\LANGsk{
\Question{
Koľko riešení má nerovnica 
\[ y\leq-\frac12x+2 \]
v \( \mathbb{N}\times\mathbb{N} \)?}
{\( 2 \)}{\( 4 \)}{nekonečne veľa}{žiadne}{}{}}
\LANGes{
\Question{
¿Cuántas soluciones tiene la inecuación
\[ y\leq-\frac12x+2 \]
en \( \mathbb{N}\times\mathbb{N} \)?}
{\( 2 \)}{\( 4 \)}{infinitas}{ninguna}{}{}}
\End

\Data\ID{1103020104}
\Updated{2020-07-15 03:07:10}
\Level{B}
\SubArea{Systems of equations and inequalities}
\LANGen{
\Question{
Among the pictures below choose the one that shows the graphic solution of the inequality
\[ 1+2y\geq y-\frac{x-9}3 \]
in \( \mathbb{R}\times\mathbb{R} \).}
{}{}{}{}{}{}}
\LANGpl{
\Question{
Spośród poniższych rysunków wybierz ten, który przedstawia rozwiązania nierówności 
\[ 1+2y\geq y-\frac{x-9}3 \]
in \( \mathbb{R}\times\mathbb{R} \).}
{}{}{}{}{}{}}
\LANGcs{
\Question{
Vyberte obrázek, na němž je šedou barvou znázorněno grafické řešení nerovnice
\[ 1+2y\geq y-\frac{x-9}3 \]
v \( \mathbb{R}\times\mathbb{R} \).}
{}{}{}{}{}{}}
\LANGsk{
\Question{
Vyberte obrázok, ktorý znázorňuje grafické riešenie nerovnice  
\[ 1+2y\geq y-\frac{x-9}3 \]
v \( \mathbb{R}\times\mathbb{R} \).}
{}{}{}{}{}{}}
\LANGes{
\Question{
NA}
{}{}{}{}{}{}}

\IMGanswer{1103020104a.pdf}
\IMGanswer{1103020104b.pdf}
\IMGanswer{1103020104c.pdf}
\IMGanswer{1103020104d.pdf}\End

\Data\ID{1103020105}
\Updated{2020-07-15 03:07:10}
\Level{B}
\SubArea{Systems of equations and inequalities}
\LANGen{
\Question{
Among the pictures below choose the one that shows the graphic solution of the inequality
\[ x-\frac{y-5}2 \leq \frac{7+2x}2 \]
in \( \mathbb{R}\times\mathbb{R} \).}
{}{}{}{}{}{}}
\LANGpl{
\Question{
Spośród poniższych rysunków wybierz graficzne rozwiązanie nierówności
\[ x-\frac{y-5}2 \leq \frac{7+2x}2 \]
in \( \mathbb{R}\times\mathbb{R} \).}
{}{}{}{}{}{}}
\LANGcs{
\Question{
Vyberte obrázek, na němž je šedou barvou znázorněno grafické řešení nerovnice
\[ x-\frac{y-5}2 \leq \frac{7+2x}2 \]
v \( \mathbb{R}\times\mathbb{R} \).}
{}{}{}{}{}{}}
\LANGsk{
\Question{
Vyberte obrázok, ktorý znázorňuje grafické riešenie nerovnice 
\[ x-\frac{y-5}2 \leq \frac{7+2x}2 \]
v \( \mathbb{R}\times\mathbb{R} \).}
{}{}{}{}{}{}}
\LANGes{
\Question{
NA}
{}{}{}{}{}{}}

\IMGanswer{1103020105a_0.pdf}
\IMGanswer{1103020105b_0.pdf}
\IMGanswer{1103020105c_0.pdf}
\IMGanswer{1103020105d_0.pdf}\End

\Data\ID{1103020106}
\Updated{2020-07-15 03:07:12}
\Level{B}
\SubArea{Systems of equations and inequalities}
\IMG{1103020106.pdf}
\LANGen{
\Question{
Which inequality fits the graphic solution showed  in the picture?}
{\( y > \frac32x+\frac{13}2 \)}{\( y \geq \frac32x+\frac{13}2 \)}{\( y < \frac32x+\frac{13}2\)}{\( y \leq \frac32x+\frac{13}2\)}{}{}}
\LANGpl{
\Question{
Rozwiązania, której nierówności przedstawiono na rysunku?}
{\( y > \frac32x+\frac{13}2 \)}{\( y \geq \frac32x+\frac{13}2 \)}{\( y < \frac32x+\frac{13}2\)}{\( y \leq \frac32x+\frac{13}2\)}{}{}}
\LANGcs{
\Question{
Která nerovnice odpovídá řešení znázorněnému na obrázku šedou barvou?}
{\( y > \frac32x+\frac{13}2 \)}{\( y \geq \frac32x+\frac{13}2 \)}{\( y < \frac32x+\frac{13}2\)}{\( y \leq \frac32x+\frac{13}2\)}{}{}}
\LANGsk{
\Question{
Ktorá nerovnica zodpovedá grafickému riešeniu znázornenému na obrázku?}
{\( y > \frac32x+\frac{13}2 \)}{\( y \geq \frac32x+\frac{13}2 \)}{\( y < \frac32x+\frac{13}2\)}{\( y \leq \frac32x+\frac{13}2\)}{}{}}
\LANGes{
\Question{
¿Qué inecuación equivale a la solución gráfica en la imagen?}
{\( y > \frac32x+\frac{13}2 \)}{\( y \geq \frac32x+\frac{13}2 \)}{\( y < \frac32x+\frac{13}2\)}{\( y \leq \frac32x+\frac{13}2\)}{}{}}
\End

\Data\ID{9000026001}
\Updated{2020-07-15 03:07:34}
\Level{B}
\SubArea{Systems of equations and inequalities}
\IMG{000260_01-figure0.pdf}
\LANGen{
\Question{
Which part of the plane describes the solution of the following system of
inequalities?
\[\begin{aligned}
 x +\phantom{ 3}y\leq &3 & &
 \\y - 2x < & - 1 & &
\end{aligned}\]}
{blue}{red}{green}{yellow}{}{}}
\LANGpl{
\Question{
Która część płaszczyzny przedstawia zbiór rozwiązań następującego układu nierówności?
\[\begin{aligned}
 x +\phantom{ 3}y\leq &3 & &
 \\y - 2x < & - 1 & &
\end{aligned}\]}
{niebieska}{czerwona}{zielona}{żółta}{}{}}
\LANGcs{
\Question{
Která část roviny znázorňuje řešení dané soustavy nerovnic?
\[\begin{aligned}
 x +\phantom{ 3}y\leq &3 & &
 \\y - 2x < & - 1 & &
\end{aligned}\]}
{modrá}{červená}{zelená}{žlutá}{}{}}
\LANGsk{
\Question{
Ktorá časť roviny znázorňuje riešenie danej sústavy nerovníc?
\[\begin{aligned}
 x +\phantom{ 3}y\leq &3 & &
 \\y - 2x < & - 1 & &
\end{aligned}\]}
{modrá}{červená}{zelená}{žltá}{}{}}
\LANGes{
\Question{
¿Qué parte del plano representa la solución del siguiente sistema de inecuaciones?
\[\begin{aligned}
 x +\phantom{ 3}y\leq &3 & &
 \\y - 2x < & - 1 & &
\end{aligned}\]}
{azul}{roja}{verde}{amarilla}{}{}}
\End

\Data\ID{9000026002}
\Updated{2020-07-15 03:07:34}
\Level{B}
\SubArea{Systems of equations and inequalities}
\IMG{000260_02-figure0.pdf}
\LANGen{
\Question{
Which part of the plane describes the solution of the following system of
inequalities?
\[\begin{aligned}
 x + y > &2 + x & &
 \\y + 1\leq &x + 1 & &
\end{aligned}\]}
{green}{red}{yellow}{blue}{}{}}
\LANGpl{
\Question{
Która część płaszczyzny przedstawia rozwiązanie następującego układu nierówności?
\[\begin{aligned}
 x + y > &2 + x & &
 \\y + 1\leq &x + 1 & &
\end{aligned}\]}
{zielona}{czerwona}{żółta}{niebieska}{}{}}
\LANGcs{
\Question{
Která část roviny znázorňuje řešení dané soustavy nerovnic?
\[\begin{aligned}
 x + y > &2 + x & &
 \\y + 1\leq &x + 1 & &
\end{aligned}\]}
{zelená}{červená}{žlutá}{modrá}{}{}}
\LANGsk{
\Question{
Ktorá časť roviny znázorňuje riešenie danej sústavy nerovníc? 
\[\begin{aligned}
 x + y > &2 + x & &
 \\y + 1\leq &x + 1 & &
\end{aligned}\]}
{zelená}{červená}{žltá}{modrá}{}{}}
\LANGes{
\Question{
¿Qué parte del plano representa la solución del siguiente sistema de inecuaciones?
\[\begin{aligned}
 x + y > &2 + x & &
 \\y + 1\leq &x + 1 & &
\end{aligned}\]}
{verde}{roja}{amarilla}{azul}{}{}}
\End

\Data\ID{9000026003}
\Updated{2020-07-15 03:07:34}
\Level{B}
\SubArea{Systems of equations and inequalities}
\IMG{000260_03-figure0.pdf}
\LANGen{
\Question{
Which part of the plane describes the solution of the following system of
inequalities?
\[\begin{aligned}
 2x - y\geq &2 & &
 \\2x + y\geq & - 2 & &
\end{aligned}\]}
{yellow}{red}{green}{blue}{}{}}
\LANGpl{
\Question{
Która część płaszczyzny przedstawia rozwiązanie następującego układu nierówności?
\[\begin{aligned}
 2x - y\geq &2 & &
 \\2x + y\geq & - 2 & &
\end{aligned}\]}
{żółta}{czerwona}{zielona}{niebieska}{}{}}
\LANGcs{
\Question{
Která část roviny znázorňuje řešení dané soustavy nerovnic?
\[\begin{aligned}
 2x - y\geq &2 & &
 \\2x + y\geq & - 2 & &
\end{aligned}\]}
{žlutá}{červená}{zelená}{modrá}{}{}}
\LANGsk{
\Question{
Ktorá časť roviny znázorňuje riešenie danej sústavy nerovníc?
\[\begin{aligned}
 2x - y\geq &2 & &
 \\2x + y\geq & - 2 & &
\end{aligned}\]}
{žltá}{červená}{zelená}{modrá}{}{}}
\LANGes{
\Question{
¿Qué parte del plano representa la solución del siguiente sistema de inecuaciones?
\[\begin{aligned}
 2x - y\geq &2 & &
 \\2x + y\geq & - 2 & &
\end{aligned}\]}
{amarilla}{roja}{verde}{azul}{}{}}
\End

\Data\ID{9000026004}
\Updated{2020-07-15 03:07:34}
\Level{B}
\SubArea{Systems of equations and inequalities}
\IMG{000260_04-figure0.pdf}
\LANGen{
\Question{
Which part of the plane describes the solution of the following system of
inequalities?
\[\begin{aligned}
 2y - x\geq &4 & &
 \\2y - x\geq & - 2 & &
\end{aligned}\]}
{blue}{red}{yellow}{}{}{}}
\LANGpl{
\Question{
Która część płaszczyzny przedstawia rozwiązanie następującego układu nierówności?
\[\begin{aligned}
 2y - x\geq &4 & &
 \\2y - x\geq & - 2 & &
\end{aligned}\]}
{niebieska}{czerwona}{żółta}{}{}{}}
\LANGcs{
\Question{
Která část roviny znázorňuje řešení dané soustavy nerovnic?
\[\begin{aligned}
 2y - x\geq &4 & &
 \\2y - x\geq & - 2 & &
\end{aligned}\]}
{modrá}{červená}{žlutá}{}{}{}}
\LANGsk{
\Question{
Ktorá časť roviny znázorňuje riešenie danej sústavy nerovníc? 
\[\begin{aligned}
 2y - x\geq &4 & &
 \\2y - x\geq & - 2 & &
\end{aligned}\]}
{modrá}{červená}{žltá}{}{}{}}
\LANGes{
\Question{
¿Qué parte del plano representa la solución del siguiente sistema de inecuaciones?
\[\begin{aligned}
 2y - x\geq &4 & &
 \\2y - x\geq & - 2 & &
\end{aligned}\]}
{azul}{roja}{amarilla}{}{}{}}
\End

\Data\ID{9000026005}
\Updated{2020-07-15 03:07:34}
\Level{B}
\SubArea{Systems of equations and inequalities}
\IMG{000260_05-figure0.pdf}
\LANGen{
\Question{
Which part of the plane describes the solution of the following system of
inequalities?
\[\begin{aligned}
 x\leq &3 & &
 \\5x > &9 - 3y & &
\end{aligned}\]}
{yellow}{red}{green}{blue}{}{}}
\LANGpl{
\Question{
Która część płaszczyzny przedstawia rozwiązanie następującego układu nierówności?
\[\begin{aligned}
 x\leq &3 & &
 \\5x > &9 - 3y & &
\end{aligned}\]}
{żółta}{czerwona}{zielona}{niebieska}{}{}}
\LANGcs{
\Question{
Která část roviny znázorňuje řešení dané soustavy nerovnic?
\[\begin{aligned}
 x\leq &3 & &
 \\5x > &9 - 3y & &
\end{aligned}\]}
{žlutá}{červená}{zelená}{modrá}{}{}}
\LANGsk{
\Question{
Ktorá časť roviny znázorňuje riešenie danej sústavy nerovníc?
\[\begin{aligned}
 x\leq &3 & &
 \\5x > &9 - 3y & &
\end{aligned}\]}
{žltá}{červená}{zelená}{modrá}{}{}}
\LANGes{
\Question{
¿Qué parte del plano representa la solución del siguiente sistema de inecuaciones?
\[\begin{aligned}
 x\leq &3 & &
 \\5x > &9 - 3y & &
\end{aligned}\]}
{amarilla}{roja}{verde}{azul}{}{}}
\End

\Data\ID{9000026006}
\Updated{2020-07-15 03:07:34}
\Level{B}
\SubArea{Systems of equations and inequalities}
\IMG{000260_06-figure0.pdf}
\LANGen{
\Question{
Which system of inequalities defines the set from the picture?}
{\(\begin{aligned}x +\phantom{ 2}y&\geq 3 &
\\y - 2x& < -1
\\ \end{aligned}\)}{\(\begin{aligned}x +\phantom{ 2}y& > 3 &
\\y - 2x& < -1
\\ \end{aligned}\)}{\(\begin{aligned}x +\phantom{ 2}y&\leq 3 &
\\y - 2x& < -1
\\ \end{aligned}\)}{\(\begin{aligned}x +\phantom{ 2}y& < 3 &
\\y - 2x& > -1
\\ \end{aligned}\)}{}{}}
\LANGpl{
\Question{
Który układ nierówności odpowiada zbiorowi przedstawionemu na rysunku poniżej?}
{\(\begin{aligned}x +\phantom{ 2}y&\geq 3 &
\\y - 2x& < -1
\\ \end{aligned}\)}{\(\begin{aligned}x +\phantom{ 2}y& > 3 &
\\y - 2x& < -1
\\ \end{aligned}\)}{\(\begin{aligned}x +\phantom{ 2}y&\leq 3 &
\\y - 2x& < -1
\\ \end{aligned}\)}{\(\begin{aligned}x +\phantom{ 2}y& < 3 &
\\y - 2x& > -1
\\ \end{aligned}\)}{}{}}
\LANGcs{
\Question{
Která soustava nerovnic odpovídá řešení, znázorněnému na obrázku
červenou barvou?}
{\(\begin{aligned}x +\phantom{ 2}y&\geq 3 &
\\y - 2x& < -1
\\ \end{aligned}\)}{\(\begin{aligned}x +\phantom{ 2}y& > 3 &
\\y - 2x& < -1
\\ \end{aligned}\)}{\(\begin{aligned}x +\phantom{ 2}y&\leq 3 &
\\y - 2x& < -1
\\ \end{aligned}\)}{\(\begin{aligned}x +\phantom{ 2}y& < 3 &
\\y - 2x& > -1
\\ \end{aligned}\)}{}{}}
\LANGsk{
\Question{
Ktorá sústava nerovníc zodpovedá riešeniu, ktoré je  znázornené na obrázku červenou farbou?}
{\(\begin{aligned}x +\phantom{ 2}y&\geq 3 &
\\y - 2x& < -1
\\ \end{aligned}\)}{\(\begin{aligned}x +\phantom{ 2}y& > 3 &
\\y - 2x& < -1
\\ \end{aligned}\)}{\(\begin{aligned}x +\phantom{ 2}y&\leq 3 &
\\y - 2x& < -1
\\ \end{aligned}\)}{\(\begin{aligned}x +\phantom{ 2}y& < 3 &
\\y - 2x& > -1
\\ \end{aligned}\)}{}{}}
\LANGes{
\Question{
¿Qué sistema de inecuaciones define el conjunto de la imagen?}
{\(\begin{aligned}x +\phantom{ 2}y&\geq 3 &
\\y - 2x& < -1
\\ \end{aligned}\)}{\(\begin{aligned}x +\phantom{ 2}y& > 3 &
\\y - 2x& < -1
\\ \end{aligned}\)}{\(\begin{aligned}x +\phantom{ 2}y&\leq 3 &
\\y - 2x& < -1
\\ \end{aligned}\)}{\(\begin{aligned}x +\phantom{ 2}y& < 3 &
\\y - 2x& > -1
\\ \end{aligned}\)}{}{}}
\End

\Data\ID{9000026007}
\Updated{2020-07-15 03:07:34}
\Level{B}
\SubArea{Systems of equations and inequalities}
\IMG{000260_07-figure0.pdf}
\LANGen{
\Question{
Which system of inequalities defines the set from the picture?}
{\(\begin{aligned}y & < 2 &
\\y + 1&\geq x + 1
\\ \end{aligned}\)}{\(\begin{aligned}y &\geq 2 &
\\y + 1& < x + 1
\\ \end{aligned}\)}{\(\begin{aligned}y & > 2 &
\\y + 1&\leq x + 1
\\ \end{aligned}\)}{\(\begin{aligned}y&\leq 2 &
\\y& > x
\\ \end{aligned}\)}{}{}}
\LANGpl{
\Question{
Który układ nierówności odpowiada zbiorowi przedstawionemu na rysunku poniżej?}
{\(\begin{aligned}y & < 2 &
\\y + 1&\geq x + 1
\\ \end{aligned}\)}{\(\begin{aligned}y &\geq 2 &
\\y + 1& < x + 1
\\ \end{aligned}\)}{\(\begin{aligned}y & > 2 &
\\y + 1&\leq x + 1
\\ \end{aligned}\)}{\(\begin{aligned}y&\leq 2 &
\\y& > x
\\ \end{aligned}\)}{}{}}
\LANGcs{
\Question{
Která soustava nerovnic odpovídá řešení, znázorněnému na obrázku
červenou barvou?}
{\(\begin{aligned}y & < 2 &
\\y + 1&\geq x + 1
\\ \end{aligned}\)}{\(\begin{aligned}y &\geq 2 &
\\y + 1& < x + 1
\\ \end{aligned}\)}{\(\begin{aligned}y & > 2 &
\\y + 1&\leq x + 1
\\ \end{aligned}\)}{\(\begin{aligned}y&\leq 2 &
\\y& > x
\\ \end{aligned}\)}{}{}}
\LANGsk{
\Question{
Ktorá sústava nerovníc zodpovedá riešeniu, ktoré je znázornené na obrázku červenou farbou?}
{\(\begin{aligned}y & < 2 &
\\y + 1&\geq x + 1
\\ \end{aligned}\)}{\(\begin{aligned}y &\geq 2 &
\\y + 1& < x + 1
\\ \end{aligned}\)}{\(\begin{aligned}y & > 2 &
\\y + 1&\leq x + 1
\\ \end{aligned}\)}{\(\begin{aligned}y&\leq 2 &
\\y& > x
\\ \end{aligned}\)}{}{}}
\LANGes{
\Question{
¿Qué sistema de inecuaciones define el conjunto de la imagen?}
{\(\begin{aligned}y & < 2 &
\\y + 1&\geq x + 1
\\ \end{aligned}\)}{\(\begin{aligned}y &\geq 2 &
\\y + 1& < x + 1
\\ \end{aligned}\)}{\(\begin{aligned}y & > 2 &
\\y + 1&\leq x + 1
\\ \end{aligned}\)}{\(\begin{aligned}y&\leq 2 &
\\y& > x
\\ \end{aligned}\)}{}{}}
\End

\Data\ID{9000026008}
\Updated{2020-07-15 03:07:34}
\Level{B}
\SubArea{Systems of equations and inequalities}
\IMG{000260_08-figure0.pdf}
\LANGen{
\Question{
Which system of inequalities defines the set from the picture?}
{\(\begin{aligned}2x - y&\leq 2 &
\\2x + y&\geq - 2
\\ \end{aligned}\)}{\(\begin{aligned}2x - y&\geq 2 &
\\2x + y&\geq - 2
\\ \end{aligned}\)}{\(\begin{aligned}2x - y&\leq 2 &
\\2x + y&\leq - 2
\\ \end{aligned}\)}{\(\begin{aligned}2x - y&\geq 2 &
\\2x + y&\leq - 2
\\ \end{aligned}\)}{}{}}
\LANGpl{
\Question{
Który układ nierówności odpowiada zbiorowi przedstawionemu na rysunku poniżej?}
{\(\begin{aligned}2x - y&\leq 2 &
\\2x + y&\geq - 2
\\ \end{aligned}\)}{\(\begin{aligned}2x - y&\geq 2 &
\\2x + y&\geq - 2
\\ \end{aligned}\)}{\(\begin{aligned}2x - y&\leq 2 &
\\2x + y&\leq - 2
\\ \end{aligned}\)}{\(\begin{aligned}2x - y&\geq 2 &
\\2x + y&\leq - 2
\\ \end{aligned}\)}{}{}}
\LANGcs{
\Question{
Která soustava nerovnic odpovídá řešení, znázorněnému na obrázku
červenou barvou?}
{\(\begin{aligned}2x - y&\leq 2 &
\\2x + y&\geq - 2
\\ \end{aligned}\)}{\(\begin{aligned}2x - y&\geq 2 &
\\2x + y&\geq - 2
\\ \end{aligned}\)}{\(\begin{aligned}2x - y&\leq 2 &
\\2x + y&\leq - 2
\\ \end{aligned}\)}{\(\begin{aligned}2x - y&\geq 2 &
\\2x + y&\leq - 2
\\ \end{aligned}\)}{}{}}
\LANGsk{
\Question{
Ktorá sústava nerovníc zodpovedá riešeniu, ktoré je znázornené na obrázku červenou farbou?}
{\(\begin{aligned}2x - y&\leq 2 &
\\2x + y&\geq - 2
\\ \end{aligned}\)}{\(\begin{aligned}2x - y&\geq 2 &
\\2x + y&\geq - 2
\\ \end{aligned}\)}{\(\begin{aligned}2x - y&\leq 2 &
\\2x + y&\leq - 2
\\ \end{aligned}\)}{\(\begin{aligned}2x - y&\geq 2 &
\\2x + y&\leq - 2
\\ \end{aligned}\)}{}{}}
\LANGes{
\Question{
¿Qué sistema de inecuaciones define el conjunto de la imagen?}
{\(\begin{aligned}2x - y&\leq 2 &
\\2x + y&\geq - 2
\\ \end{aligned}\)}{\(\begin{aligned}2x - y&\geq 2 &
\\2x + y&\geq - 2
\\ \end{aligned}\)}{\(\begin{aligned}2x - y&\leq 2 &
\\2x + y&\leq - 2
\\ \end{aligned}\)}{\(\begin{aligned}2x - y&\geq 2 &
\\2x + y&\leq - 2
\\ \end{aligned}\)}{}{}}
\End

\Data\ID{9000026009}
\Updated{2020-07-15 03:07:34}
\Level{B}
\SubArea{Systems of equations and inequalities}
\IMG{000260_09-figure0.pdf}
\LANGen{
\Question{
Which system of inequalities defines the set from the picture?}
{\(\begin{aligned}2y -\phantom{ 2}x& < 4&
\\x - 2y & < 2
\\ \end{aligned}\)}{\(\begin{aligned}2y -\phantom{ 2}x& < 4&
\\x - 2y & > 2
\\ \end{aligned}\)}{\(\begin{aligned}2y - x& > 4 &
\\2y - x& < -2
\\ \end{aligned}\)}{\(\begin{aligned}2y - x& > 4 &
\\2y - x& > -2
\\ \end{aligned}\)}{}{}}
\LANGpl{
\Question{
Który układ nierówności odpowiada zbiorowi przedstawionemu na rysunku poniżej?}
{\(\begin{aligned}2y -\phantom{ 2}x& < 4&
\\x - 2y & < 2
\\ \end{aligned}\)}{\(\begin{aligned}2y -\phantom{ 2}x& < 4&
\\x - 2y & > 2
\\ \end{aligned}\)}{\(\begin{aligned}2y - x& > 4 &
\\2y - x& < -2
\\ \end{aligned}\)}{\(\begin{aligned}2y - x& > 4 &
\\2y - x& > -2
\\ \end{aligned}\)}{}{}}
\LANGcs{
\Question{
Která soustava nerovnic odpovídá řešení, znázorněnému na obrázku
červenou barvou?}
{\(\begin{aligned}2y -\phantom{ 2}x& < 4&
\\x - 2y & < 2
\\ \end{aligned}\)}{\(\begin{aligned}2y -\phantom{ 2}x& < 4&
\\x - 2y & > 2
\\ \end{aligned}\)}{\(\begin{aligned}2y - x& > 4 &
\\2y - x& < -2
\\ \end{aligned}\)}{\(\begin{aligned}2y - x& > 4 &
\\2y - x& > -2
\\ \end{aligned}\)}{}{}}
\LANGsk{
\Question{
Ktorá sústava nerovníc zodpovedá riešeniu, ktoré je znázornené na obrázku červenou farbou?}
{\(\begin{aligned}2y -\phantom{ 2}x& < 4&
\\x - 2y & < 2
\\ \end{aligned}\)}{\(\begin{aligned}2y -\phantom{ 2}x& < 4&
\\x - 2y & > 2
\\ \end{aligned}\)}{\(\begin{aligned}2y - x& > 4 &
\\2y - x& < -2
\\ \end{aligned}\)}{\(\begin{aligned}2y - x& > 4 &
\\2y - x& > -2
\\ \end{aligned}\)}{}{}}
\LANGes{
\Question{
¿Qué sistema de inecuaciones define el conjunto de la imagen?}
{\(\begin{aligned}2y -\phantom{ 2}x& < 4&
\\x - 2y & < 2
\\ \end{aligned}\)}{\(\begin{aligned}2y -\phantom{ 2}x& < 4&
\\x - 2y & > 2
\\ \end{aligned}\)}{\(\begin{aligned}2y - x& > 4 &
\\2y - x& < -2
\\ \end{aligned}\)}{\(\begin{aligned}2y - x& > 4 &
\\2y - x& > -2
\\ \end{aligned}\)}{}{}}
\End

\Data\ID{9000026010}
\Updated{2020-07-15 03:07:34}
\Level{B}
\SubArea{Systems of equations and inequalities}
\IMG{000260_10-figure0.pdf}
\LANGen{
\Question{
Which system of inequalities defines the set from the picture?}
{\(\begin{aligned}x &\leq 3 &
\\5x& < 9 - 3y
\\ \end{aligned}\)}{\(\begin{aligned}x & < 3 &
\\5x& < 9 - 3y
\\ \end{aligned}\)}{\(\begin{aligned}x & > 3 &
\\5x& < 9 - 3y
\\ \end{aligned}\)}{\(\begin{aligned}x &\leq 3 &
\\5x& > 9 - 3y
\\ \end{aligned}\)}{}{}}
\LANGpl{
\Question{
Który układ nierówności odpowiada zbiorowi przedstawionemu na rysunku poniżej?}
{\(\begin{aligned}x &\leq 3 &
\\5x& < 9 - 3y
\\ \end{aligned}\)}{\(\begin{aligned}x & < 3 &
\\5x& < 9 - 3y
\\ \end{aligned}\)}{\(\begin{aligned}x & > 3 &
\\5x& < 9 - 3y
\\ \end{aligned}\)}{\(\begin{aligned}x &\leq 3 &
\\5x& > 9 - 3y
\\ \end{aligned}\)}{}{}}
\LANGcs{
\Question{
Která soustava nerovnic odpovídá řešení, znázorněnému na obrázku
červenou barvou?}
{\(\begin{aligned}x &\leq 3 &
\\5x& < 9 - 3y
\\ \end{aligned}\)}{\(\begin{aligned}x & < 3 &
\\5x& < 9 - 3y
\\ \end{aligned}\)}{\(\begin{aligned}x & > 3 &
\\5x& < 9 - 3y
\\ \end{aligned}\)}{\(\begin{aligned}x &\leq 3 &
\\5x& > 9 - 3y
\\ \end{aligned}\)}{}{}}
\LANGsk{
\Question{
Ktorá sústava nerovníc zodpovedá riešeniu, ktoré je znázornené na obrázku červenou farbou?}
{\(\begin{aligned}x &\leq 3 &
\\5x& < 9 - 3y
\\ \end{aligned}\)}{\(\begin{aligned}x & < 3 &
\\5x& < 9 - 3y
\\ \end{aligned}\)}{\(\begin{aligned}x & > 3 &
\\5x& < 9 - 3y
\\ \end{aligned}\)}{\(\begin{aligned}x &\leq 3 &
\\5x& > 9 - 3y
\\ \end{aligned}\)}{}{}}
\LANGes{
\Question{
¿Qué sistema de inecuaciones define el conjunto de la imagen?}
{\(\begin{aligned}x &\leq 3 &
\\5x& < 9 - 3y
\\ \end{aligned}\)}{\(\begin{aligned}x & < 3 &
\\5x& < 9 - 3y
\\ \end{aligned}\)}{\(\begin{aligned}x & > 3 &
\\5x& < 9 - 3y
\\ \end{aligned}\)}{\(\begin{aligned}x &\leq 3 &
\\5x& > 9 - 3y
\\ \end{aligned}\)}{}{}}
\End

\Data\ID{9000031103}
\Updated{2020-08-23 08:08:44}
\Level{B}
\SubArea{Systems of equations and inequalities}
\LANGen{
\Question{
Solve the following system of equations and identify a correct statement.
\[\begin{aligned}
 x - 2y + 5 = 0 & &
 \\x^{2} + y^{2} = 9 & &
\end{aligned}\]}
{The system has two solutions.}{The system does not have any solution.}{The system has a unique solution.}{The system has more than two solutions.}{}{}}
\LANGpl{
\Question{
Rozwiąż następujący układ równań i wybierz poprawną odpowiedź. 
\[\begin{aligned}
 x - 2y + 5 = 0 & &
 \\x^{2} + y^{2} = 9 & &
\end{aligned}\]}
{Dwa rozwiązania.}{Nie ma rozwiązania.}{Tylko jedno rozwiązanie.}{Więcej niż dwa rozwiązania.}{}{}}
\LANGcs{
\Question{
Je dána soustava rovnic:
\[\begin{aligned}
 x - 2y + 5 = 0, & &
 \\x^{2} + y^{2} = 9. & &
\end{aligned}\]
Vyberte správné tvrzení.}
{Soustava má právě dvě řešení.}{Soustava nemá řešení.}{Soustava má právě jedno řešení.}{Soustava má více než dvě řešení.}{}{}}
\LANGsk{
\Question{
Je daná sústava rovníc:
\[\begin{aligned}
 x - 2y + 5 = 0, & &
 \\x^{2} + y^{2} = 9. & &
\end{aligned}\] Vyberte správne tvrdenie.}
{Sústava má práve dve riešenia.}{Sústava nemá riešenie.}{Sústava má práve jedno riešenie.}{Sústava má viac než dve riešenia.}{}{}}
\LANGes{
\Question{
Dado el sistema de ecuaciones:
\[\begin{aligned}
 x - 2y + 5 = 0 & &
 \\x^{2} + y^{2} = 9 & &
\end{aligned}\]
Identifica la proposición lógica verdadera.}
{El sistema tiene dos soluciones.}{El sistema no tiene soluciones.}{El sistema tiene solo una solución.}{El sistema tiene más de dos soluciones.}{}{}}
\End

\Data\ID{9000031104}
\Updated{2020-08-23 08:08:01}
\Level{B}
\SubArea{Systems of equations and inequalities}
\LANGen{
\Question{
Solve the following system of equations and identify a correct statement.
\[\begin{aligned}
 \frac{x} 
{y + 1} - \frac{2} 
{x + 1} & = 0 & &
 \\\frac{y}
{x} + \frac{2} 
{x} & = -1 & &
\end{aligned}\]}
{The system has a unique solution.}{The system does not have any solution.}{The system has two solutions.}{The system has infinitely many solutions.}{}{}}
\LANGpl{
\Question{
Rozwiąż następujący układ równań i wybierz poprawną odpowiedź. 
\[\begin{aligned}
 \frac{x} 
{y + 1} - \frac{2} 
{x + 1} & = 0 & &
 \\\frac{y}
{x} + \frac{2} 
{x} & = -1 & &
\end{aligned}\]}
{Tylko jedno rozwiązanie.}{Nie ma rozwiązania.}{Dwa rozwiązania.}{Nieskończenie wiele rozwiązań.}{}{}}
\LANGcs{
\Question{
Je dána soustava rovnic:
\[\begin{aligned}
 \frac{x} 
{y + 1} - \frac{2} 
{x + 1} & = 0, & &
 \\\frac{y}
{x} + \frac{2} 
{x} & = -1. & &
\end{aligned}\]
Vyberte správné tvrzení.}
{Soustava má právě jedno řešení.}{Soustava nemá řešení.}{Soustava má právě dvě řešení.}{Soustava má nekonečně mnoho řešení.}{}{}}
\LANGsk{
\Question{
Je daná sústava rovníc: 
\[\begin{aligned}
 \frac{x} 
{y + 1} - \frac{2} 
{x + 1} & = 0, & &
 \\\frac{y}
{x} + \frac{2} 
{x} & = -1. & &
\end{aligned}\] Vyberte správne tvrdenie.}
{Sústava má práve jedno riešenie.}{Sústava nemá riešenie.}{Sústava má práve dve riešenia.}{Sústava má nekonečne veľa riešení.}{}{}}
\LANGes{
\Question{
Dado el sistema de ecuaciones:
\[\begin{aligned}
 \frac{x} 
{y + 1} - \frac{2} 
{x + 1} & = 0 & &
 \\\frac{y}
{x} + \frac{2} 
{x} & = -1 & &
\end{aligned}\]
Identifica la proposición lógica verdadera.}
{El sistema tiene solo una solución.}{El sistema no tiene soluciones.}{El sistema tiene dos soluciones.}{El sistema tiene infinitas soluciones.}{}{}}
\End

\Data\ID{9000031105}
\Updated{2020-08-23 08:08:20}
\Level{B}
\SubArea{Systems of equations and inequalities}
\LANGen{
\Question{
Solve the following system of equations and identify a correct statement.
\[\begin{aligned}
 \frac{1} 
{x + 1} -\frac{1} 
{y} = 0 & &
 \\y^{2} = 1 & &
\end{aligned}\]}
{The system has two solutions.}{The system does not have any solution.}{The system has a unique solution.}{The system has infinitely many solutions.}{}{}}
\LANGpl{
\Question{
Rozwiąż następujący układ równań i wybierz poprawną odpowiedź. 
\[\begin{aligned}
 \frac{1} 
{x + 1} -\frac{1} 
{y} = 0 & &
 \\y^{2} = 1 & &
\end{aligned}\]}
{Dwa rozwiązania.}{Nie ma rozwiązania.}{Tylko jedno rozwiązanie.}{Nieskończenie wiele rozwiązań.}{}{}}
\LANGcs{
\Question{
Je dána soustava rovnic:
\[\begin{aligned}
 \frac{1} 
{x + 1} -\frac{1} 
{y} = 0, & &
 \\y^{2} = 1. & &
\end{aligned}\]
Vyberte správné tvrzení.}
{Soustava má právě dvě řešení.}{Soustava nemá řešení.}{Soustava má právě jedno řešení.}{Soustava má nekonečně mnoho řešení.}{}{}}
\LANGsk{
\Question{
Je daná sústava rovníc: 
\[\begin{aligned}
 \frac{1} 
{x + 1} -\frac{1} 
{y} = 0, & &
 \\y^{2} = 1. & &
\end{aligned}\] Vyberte správne tvrdenie.}
{Sústava má práve dve riešenia.}{Sústava nemá riešenie.}{Sústava má práve jedno riešenie.}{Sústava má nekonečne veľa riešení.}{}{}}
\LANGes{
\Question{
Dado el sistema de ecuaciones:
\[\begin{aligned}
 \frac{1} 
{x + 1} -\frac{1} 
{y} = 0 & &
 \\y^{2} = 1 & &
\end{aligned}\]
Identifica la proposición lógica verdadera.}
{El sistema tiene dos soluciones.}{El sistema no tiene soluciones.}{El sistema tiene solo una solución.}{El sistema tiene infinitas soluciones.}{}{}}
\End

\Data\ID{9000039001}
\Updated{2020-07-15 03:07:34}
\Level{B}
\SubArea{Systems of equations and inequalities}
\LANGen{
\Question{
Assuming \(x\in \mathbb{R}\),
find the solution set of the following system.
\[\begin{aligned}
 - 4x & < -6 - 2x & &
 \\2(x + 5) &\leq 16 & &
\end{aligned}\]}
{\(\emptyset \)}{\((3;+\infty )\)}{\((-\infty ;3)\)}{\(\{3\}\)}{}{}}
\LANGpl{
\Question{
Zakładając, że \(x\in \mathbb{R}\),
wyznacz zbiór rozwiązań następującego układu.
\[\begin{aligned}
 - 4x & < -6 - 2x & &
 \\2(x + 5) &\leq 16 & &
\end{aligned}\]}
{\(\emptyset \)}{\((3;+\infty )\)}{\((-\infty ;3)\)}{\(\{3\}\)}{}{}}
\LANGcs{
\Question{
Je dána soustava nerovnic
\[\begin{aligned}
 - 4x & < -6 - 2x, & &
 \\2(x + 5) &\leq 16. & &
\end{aligned}\]
Množinou řešení této soustavy v
\(\mathbb{R}\) je:}
{$\emptyset$}{\((3;+\infty )\)}{\((-\infty ;3)\)}{\(\{3\}\)}{}{}}
\LANGsk{
\Question{
Nájdite množinu riešenia danej sústavy pre \(x\in \mathbb{R}\). 
\[\begin{aligned}
 - 4x & < -6 - 2x & &
 \\2(x + 5) &\leq 16 & &
\end{aligned}\]}
{\(\emptyset \)}{\((3;+\infty )\)}{\((-\infty ;3)\)}{\(\{3\}\)}{}{}}
\LANGes{
\Question{
Suponiendo que \(x\in \mathbb{R}\),
halla la solución del siguiente sistema.
\[\begin{aligned}
 - 4x & < -6 - 2x & &
 \\2(x + 5) &\leq 16 & &
\end{aligned}\]}
{\(\emptyset \)}{\((3;+\infty )\)}{\((-\infty ;3)\)}{\(\{3\}\)}{}{}}
\End

\Data\ID{1003034501}
\Updated{2020-08-23 08:08:36}
\Level{C}
\SubArea{Systems of equations and inequalities}
\LANGen{
\Question{
Two different aquarium fish shops have special price on Congo Tetra. The price is \( 42\,\mathrm{CZK} \) for one fish. Further, the discount of \( 50\,\mathrm{CZK} \) from a purchase over \( 300\,\mathrm{CZK} \) is offered in shop A.  In shop B a customer is given \( 5\% \) discount of the price of any purchase. How many of Congo Tetra must one buy to make the total price in shop A lower than the total price in shop B?}
{greater than \( 7 \) and less than \( 24 \)}{less than \( 24 \)}{greater than \( 23 \)}{less than \( 7 \)}{}{}}
\LANGpl{
\Question{
Dwa różne sklepy z rybkami akwariowymi mają specjalną cenę rybek Congo Tetra. Cena wynosi \( 42\,\mathrm{PLN} \) za jedną rybkę. W sklepie A oferowany jest rabat w wysokości \( 50\,\mathrm{PLN} \) przy zakupie za ponad \( 300\,\mathrm{PLN} \).  W sklepie B klient otrzymuje zniżkę w wysokości  \( 5\% \) przy każdym zakupie. Ile rybek Congo Tetra należy kupić, aby ostateczna cena za zakupy w sklepie A była niższa niż ostateczna cena w sklepie B?}
{więcej niż  \( 7 \) ale mniej niż  \( 24 \)}{mniej niż  \( 24 \)}{więcej niż \( 23 \)}{mniej niż \( 7 \)}{}{}}
\LANGcs{
\Question{
Dva prodejci akvarijních rybek nabízí v akci Tetru konžskou za  \( 42\,\mathrm{CZK} \) za kus. Prodejce A nabízí slevu \( 50\,\mathrm{CZK} \) při nákupu nad \( 300\,\mathrm{CZK} \).  Prodejce B nabízí  \( 5\% \) slevu z jakéhokoliv nákupu. Kolik daných rybek musíme koupit, má-li být výhodnější nákup u prodejce A?}
{více než \( 7 \) a méně než \( 24 \)}{méně než \( 24 \)}{více než \( 23 \)}{méně než \( 7 \)}{}{}}
\LANGsk{
\Question{
Dvaja predajcovia akvarijných rybiek ponúkajú v akcii Tetru konžskú za \( 42\,\mathrm{CZK} \) za kus. Predajca A ponúka zľavu \( 50\,\mathrm{CZK} \) pri nákupe nad \( 300\,\mathrm{CZK} \).  Predajca B ponúka \( 5\% \) zľavu z akéhokoľvek nákupu. Koľko daných rybiek musíme kúpiť, ak má byť  výhodnejší nákup u predajcu A?}
{viac než \( 7 \) a menej než \( 24 \)}{menej než \( 24 \)}{viac než \( 23 \)}{menej než \( 7 \)}{}{}}
\LANGes{
\Question{
Dos vendedores de peces de acuario venden un Congo tetra por \( 42\,\mathrm{CZK} \). Además, el vendedor A ofrece UN descuento de \( 50\,\mathrm{CZK} \) de una compra superior a \( 300\,\mathrm{CZK} \). El vendedor B ofrece a sus clientes un \( 5\% \) de descuento del precio de cualquier compra. ¿Cuántos peces de Congo tetra hay que comprar para que el precio total del vendedor A sea inferior al precio total del vendedor B?}
{más de \( 7 \) y menos de \( 24 \)}{menos de \( 24 \)}{más de \( 23 \)}{menos de \( 7 \)}{}{}}
\End

\Data\ID{1003034502}
\Updated{2020-08-23 08:08:56}
\Level{C}
\SubArea{Systems of equations and inequalities}
\LANGen{
\Question{
Petr would like to buy a new smartphone. If he starts a temporary job at the electro shop he gets payed \( 120\,\mathrm{CZK} \) per hour plus he gets \( 20\% \) discount on the smartphone bough in the shop. He calculated that for \( 24 \) hours of work he would not earn even half of the phone price. Another employer pays \( 150\,\mathrm{CZK} \) per hour. If Petr gets a temporary job with the other employer he is not eligible for the discount in the electro shop anymore, however he can  buy the smartphone from e-shop for the price by \( 600\,\mathrm{CZK} \) lower than in the electro shop and for \( 20 \) hours of work he earns more than one third of the electro shop smartphone price. Determine as closely as possible the smartphone price in the electronics shop.}
{greater than \( 7\,200\,\mathrm{CZK} \) and less than \( 9\,600\,\mathrm{CZK} \)}{greater than \( 7\,200\,\mathrm{CZK} \) and less than \( 10\,800\,\mathrm{CZK} \)}{greater than \( 4\,800\,\mathrm{CZK} \) and less than \( 9\,600\,\mathrm{CZK} \)}{greater than \( 4\,800\,\mathrm{CZK} \) and less than \( 10\,800\,\mathrm{CZK} \)}{}{}}
\LANGpl{
\Question{
Piotrek chciałby kupić nowy smartfon. Jeśli zacznie pracować dorywczo w sklepie elektronicznym będzie zarabiał  \( 12\,\mathrm{PLN} \) na godzinę i dostanie \( 20\% \) zniżki na smartfona kupionego w sklepie. Obliczył, że za  \( 24 \) godziny pracy nie zarobi nawet połowy sumy, którą musi zapłacić za smartfona. Inny pracodawca płaci \( 15\,\mathrm{PLN} \) za godzinę. Jeśli Piotrek będzie pracował dorywczo u tego sprzedawcy nie będzie miał prawa do zniżki w sklepie elektronicznym, ale będzie miał możliwość zakupu smartfona w e-sklepie za cenę o \( 60\,\mathrm{PLN} \) niższą niż w sklepie elektronicznym. Poza tym za \( 20 \) godzin pracy zarobi więcej niż jedną trzecią ceny smartfona w sklepie elektronicznym. Określ najdokładniej jak potrafisz cenę smartfona w sklepie elektronicznym.}
{wyższa niż  \( 720\,\mathrm{PLN} \) ale niższa niż  \( 960\,\mathrm{PLN} \)}{wyższa niż  \( 720\,\mathrm{PLN} \) i niższą niż  \( 1\,800\,\mathrm{PLN} \)}{wyższa niż  \( 480\,\mathrm{PLN} \) ale niższa niż \( 960\,\mathrm{PLN} \)}{wyższa niż  \( 480\,\mathrm{PLN} \) ale niższa niż  \( 1\,080\,\mathrm{PLN} \)}{}{}}
\LANGcs{
\Question{
Petr by si rád koupil nový mobil. Pokud by nastoupil na brigádu k prodejci elektro, dostal by odměnu \( 120\,\mathrm{CZK} \) za odpracovanou hodinu a \( 20\% \) slevu na mobil, který by si v této prodejně koupil. Spočítal si, že za \( 24 \) odpracovaných hodin by si nevydělal ani polovinu potřebných peněz. Jiný zaměstnavatel platí \( 150\,\mathrm{CZK} \) za hodinu. Pokud by Petr nastoupil na brigádu u něj, slevu u prodejce elektro by nedostal, ale mohl by si mobil koupit v e-shopu, v němž se prodává o \( 600\,\mathrm{CZK} \) levněji než u prodejce elektro. Za \( 20 \) hodin by si vydělal více než třetinu peněz potřebných ke koupi mobilu v e-shopu. Kolik mohl stát mobil v prodejně elektro?}
{více než \( 7\,200\,\mathrm{CZK} \) a méně než \( 9\,600\,\mathrm{CZK} \)}{více než \( 7\,200\,\mathrm{CZK} \) a méně než \( 10\,800\,\mathrm{CZK} \)}{více než \( 4\,800\,\mathrm{CZK} \) a méně než \( 9\,600\,\mathrm{CZK} \)}{více než \( 4\,800\,\mathrm{CZK} \) a méně než \( 10\,800\,\mathrm{CZK} \)}{}{}}
\LANGsk{
\Question{
Petr by si rád kúpil nový mobil. Ak by nastúpil na brigádu k predajcovi elektra, dostal by odmenu \( 120\,\mathrm{CZK} \) za odpracovanú hodinu a \( 20\% \) zľavu na mobil, ktorý by si v tejto predajni kúpil. Spočítal si, že za \( 24 \) odpracovaných hodín by si nezarobil ani polovicu potrebných peňazí. Iný zamestnávateľ platí \( 150\,\mathrm{CZK} \) za hodinu. Ak by Peter nastúpil na brigádu u neho, zľavu u predajcu elektra by nedostal, ale mohol by si mobil kúpiť v e-shope, v ktorom sa predáva o  \( 600\,\mathrm{CZK} \) lacnejšie než u predajcu elektra. Za \( 20 \) odpracovaných hodín by si zarobil viac než tretinu peňazí potrebných ku kúpe mobilu v predajni elektra. Koľko mohol stáť mobil v predajni elektra?}
{viac než \( 7\,200\,\mathrm{CZK} \) a menej než \( 9\,600\,\mathrm{CZK} \)}{viac než \( 7\,200\,\mathrm{CZK} \) a menej než \( 10\,800\,\mathrm{CZK} \)}{viac než \( 4\,800\,\mathrm{CZK} \) a menej než \( 9\,600\,\mathrm{CZK} \)}{viac než \( 4\,800\,\mathrm{CZK} \) a menej než \( 10\,800\,\mathrm{CZK} \)}{}{}}
\LANGes{
\Question{
A Pedro le gustaría comprarse un nuevo teléfono inteligente. Si comienza un trabajo temporal en la tienda de electrodomésticos, le pagarán \( 120\,\mathrm{CZK} \) por hora y, además, recibirá un descuento de \( 20\% \) para el móvil comprando en la tienda. He calculado que por \( 24 \) horas de trabajo no ganaría ni la mitad del precio del teléfono. Otro empleador paga \( 150\,\mathrm{CZK} \) por hora. Si Pedro consigue un trabajo temporal del otro empleador, no recibirá el descuento en la tienda de electrodomésticos, sin embargo, podrá comprarse el teléfono inteligente en la tienda online por un precio \( 600\,\mathrm{CZK} \) más bajo que en la tienda de electrodomésticos y por \( 20 \) horas de trabajo ganará más de un tercio del precio del teléfono inteligente de la tienda de electrodomésticos. Determina el precio posible del teléfono inteligente en la tienda de electrodomésticos.}
{más de \( 7\,200\,\mathrm{CZK} \) y menos de \( 9\,600\,\mathrm{CZK} \)}{más de \( 7\,200\,\mathrm{CZK} \) y menos de \( 10\,800\,\mathrm{CZK} \)}{más de \( 4\,800\,\mathrm{CZK} \) y menos de \( 9\,600\,\mathrm{CZK} \)}{más de \( 4\,800\,\mathrm{CZK} \) y menos de \( 10\,800\,\mathrm{CZK} \)}{}{}}
\End

\Data\ID{1003034503}
\Updated{2020-08-23 08:08:48}
\Level{C}
\SubArea{Systems of equations and inequalities}
\LANGen{
\Question{
Students registered for sports camps. For the biking camp registered by \( 18 \) students more than for the boating camp. After some time one of the students switched his registration from the boating camp to the biking camp. Now, there is two times more bikers than boaters. How many students registered  originally for the boating camp?}
{\( 21 \)}{\( 39 \)}{\( 20 \)}{\( 15 \)}{}{}}
\LANGpl{
\Question{
Uczniowie zapisali się na obóz sportowy. Na obóz rowerowy zapisało się o \( 18 \) uczniów więcej niż na obóz żeglarski. Po pewnym czasie jeden z uczniów przepisał się z obozu żeglarskiego na obóz rowerowy. Teraz jest dwa razy więcej rowerzystów niż żeglarzy. Ilu uczniów zapisało się początkowo na obóz żeglarski?}
{\( 21 \)}{\( 39 \)}{\( 20 \)}{\( 15 \)}{}{}}
\LANGcs{
\Question{
Studenti se hlásili na sportovní kurz. Cyklistický kurz si vybralo o  \( 18 \) studentů více, než vodácký. Po nějaké době jeden student převedl přihlášku z vodáckého kurzu na cyklistický. Nyní je cyklistů dvakrát více, než vodáků. Kolik studentů se původně hlásilo na vodácký kurz?}
{\( 21 \)}{\( 39 \)}{\( 20 \)}{\( 15 \)}{}{}}
\LANGsk{
\Question{
Študenti sa prihlasovali na športové tábory. Na cyklistický tábor sa prihlásilo o \( 18 \) študentov viac než na vodácky tábor. Po nejakej dobe jeden študent presunul svoju prihlášku z vodáckeho tábora na cyklistický tábor. Teraz je cyklistov dvakrát viac než vodákov. Koľko študentov sa pôvodne hlásilo na vodácky tábor?}
{\( 21 \)}{\( 39 \)}{\( 20 \)}{\( 15 \)}{}{}}
\LANGes{
\Question{
Algunos estudiantes se han matriculado en los campamentos deportivos. Para el campamento de ciclismo se han matriculado \( 18 \) estudiantes más que para el campamento de navegación. Después de algún tiempo, uno de los estudiantes ha cambiado su inscripción del campamento de navegación al campamento de ciclismo. Ahora hay dos veces más ciclistas que navegantes. ¿Cuántos estudiantes se han matriculado originalmente en el campamento de navegación?}
{\( 21 \)}{\( 39 \)}{\( 20 \)}{\( 15 \)}{}{}}
\End

\Data\ID{1003034504}
\Updated{2020-08-23 08:08:38}
\Level{C}
\SubArea{Systems of equations and inequalities}
\LANGen{
\Question{
A double-digit number equals to the four multiple of its digit sum. Determine the digit in the unit position if you know that the digit of tens is by \( 4 \) lower than the digit of units.}
{\( 8 \)}{\( 4 \)}{\( 6 \)}{\( 0 \)}{}{}}
\LANGpl{
\Question{
Dwucyfrowa liczba jest równa czterokrotności sumy jej cyfr. Określ cyfrę w pozycji jedności, jeśli wiesz, że cyfra dziesiątek jest o \( 4 \) mniejsza od cyfry jedności.}
{\( 8 \)}{\( 4 \)}{\( 6 \)}{\( 0 \)}{}{}}
\LANGcs{
\Question{
Dvouciferné číslo se rovná čtyřnásobku jeho ciferného součtu. Určete počet jednotek tohoto čísla, platí-li, že počet desítek je o \( 4 \) menší, než počet jednotek.}
{\( 8 \)}{\( 4 \)}{\( 6 \)}{\( 0 \)}{}{}}
\LANGsk{
\Question{
Dvojciferné číslo sa rovná štvornásobku jeho ciferného súčtu. Určte počet jednotiek tohto čísla, ak platí, že počet desiatok je o \( 4 \) menší, než počet jednotiek.}
{\( 8 \)}{\( 4 \)}{\( 6 \)}{\( 0 \)}{}{}}
\LANGes{
\Question{
Un número de dos dígitos equivale al cuatro múltiplo de la suma de sus dígitos. Determina el dígito en la posición de la unidad si se sabe que el dígito de las decenas es \( 4 \) veces menor que el dígito de las unidades.}
{\( 8 \)}{\( 4 \)}{\( 6 \)}{\( 0 \)}{}{}}
\End

\Data\ID{1003034505}
\Updated{2020-08-23 08:08:18}
\Level{C}
\SubArea{Systems of equations and inequalities}
\LANGen{
\Question{
The March price of a T-shirt and shorts was \( 600\,\mathrm{CZK} \) together. In April there was on store price adjustment. The price of the shorts decreased by \( 10\% \) and the price of the T-shirt increased by \( 10\% \). So the April price of both together the shorts and the T-shirt was by \( 20\,\mathrm{CZK} \) lower. What was the April price of the T-shirt?}
{\( 220\,\mathrm{CZK} \)}{\( 200\,\mathrm{CZK} \)}{\( 180\,\mathrm{CZK} \)}{\( 400\,\mathrm{CZK} \)}{}{}}
\LANGpl{
\Question{
Marcowa cena podkoszulka i krótkich spodenek wyniosła  \( 60\,\mathrm{PLN} \) za zestaw.  W kwietniu ceny w sklepie uległy zmianom. Cena krótkich spodenek obniżyła się o \( 10\% \) a cena podkoszulka wzrosła o \( 10\% \). Tak więc kwietniowa cena zestawu podkoszulek i krótkie spodenki była o \( 2\,\mathrm{PLN} \) niższa. Jaka była kwietniowa cena podkoszulka?}
{\( 22\,\mathrm{PLN} \)}{\( 20\,\mathrm{PLN} \)}{\( 18\,\mathrm{PLN} \)}{\( 40\,\mathrm{PLN} \)}{}{}}
\LANGcs{
\Question{
V březnu stálo tričko a kraťasy celkem  \( 600\,\mathrm{CZK} \). V dubnu ale došlo k úpravám cen. Kraťasy byly zlevněny o \( 10\,\% \) a tričko o \( 10\,\% \) zdraženo. Úpravou cen byl dubnový nákup trička a kraťas o \( 20\,\mathrm{CZK} \) levnější. Určete dubnovou cenu trička.}
{\( 220\,\mathrm{CZK} \)}{\( 200\,\mathrm{CZK} \)}{\( 180\,\mathrm{CZK} \)}{\( 400\,\mathrm{CZK} \)}{}{}}
\LANGsk{
\Question{
V marci stálo tričko a kraťasy celkom \( 600\,\mathrm{CZK} \). V apríli došlo k úpravám cien. Kraťasy boli zlacnené o \( 10\% \) a tričko o \( 10\% \) zdražené. Úpravou cien bol aprílový nákup trička a kraťas o \( 20\,\mathrm{CZK} \) lacnejší. Aká bola aprílová cena trička?}
{\( 220\,\mathrm{CZK} \)}{\( 200\,\mathrm{CZK} \)}{\( 180\,\mathrm{CZK} \)}{\( 400\,\mathrm{CZK} \)}{}{}}
\LANGes{
\Question{
En marzo una camiseta y unos pantalones cortos costaron \( 600\,\mathrm{CZK} \) en total. En abril los precios han cambiado . El precio de los pantalones cortos han disminuido en un \( 10\% \) y el precio de la camiseta ha aumentado en un \( 10\% \). Por lo tanto, en abril el precio total de los pantalones cortos y de la camiseta ha sido un \( 20\,\mathrm{CZK} \) más bajo. ¿Cuánto ha costado la camiseta en abril?}
{\( 220\,\mathrm{CZK} \)}{\( 200\,\mathrm{CZK} \)}{\( 180\,\mathrm{CZK} \)}{\( 400\,\mathrm{CZK} \)}{}{}}
\End

\Data\ID{1003034506}
\Updated{2020-07-15 03:07:21}
\Level{C}
\SubArea{Systems of equations and inequalities}
\LANGen{
\Question{
Kamil is able to mow a meadow in \( 12 \) hours. Zdeněk has a better lawn mower and he is able to mow the same meadow in \( 8 \) hours. They have agreed that Kamil starts to mow alone sooner and Zdeněk will join him later so that the total time of mowing is \( 9 \) hours. How long will they mow together?}
{\( 2 \) hours}{\( 7 \) hours}{\( 6 \) hours}{\( 3 \) hours}{}{}}
\LANGpl{
\Question{
Kamil potrafi skosić trawę w \( 12 \) godzin. Zbyszek ma lepszą kosiarkę i jest w stanie skosić tę samą łąkę w \( 8 \) godzin.  Ustalili, że Kamil najpierw zacznie kosić sam, a Zbyszek dołączy do niego później, tak, aby całkowity czas koszenia wyniósł \( 9 \) godzin. Jak długo będą kosić razem?}
{\( 2 \) godziny}{\( 7 \) godzin}{\( 6 \) godzin}{\( 3 \) godziny}{}{}}
\LANGcs{
\Question{
Kamil je schopen posekat louku za \( 12 \) hodin. Zdeněk má  lepší  sekačku a stejnou louku by sám posekal za \( 8 \) hodin. Domluvili se, že Kamil začne sekat sám a Zdeněk se přidá později tak, aby s posekáním louky byli hotovi za \( 9 \) hodin. Jak dlouho budou sekat oba společně?}
{\( 2 \) hodiny}{\( 7 \) hodin}{\( 6 \) hodin}{\( 3 \) hodiny}{}{}}
\LANGsk{
\Question{
Kamil je schopný pokosiť lúku za \( 12 \) hodín. Zdeněk má lepšiu kosačku a tú istú lúku by sám pokosil za \( 8 \) hodín. Dohodli sa, že Kamil začne kosiť sám a Zdeněk sa pripojí k nemu neskôr tak, aby s pokosením lúky boli hotoví za \( 9 \) hodín. Ako dlho budú kosiť obidvaja spoločne?}
{\( 2 \) hodiny}{\( 7 \) hodín}{\( 6 \) hodín}{\( 3 \) hodiny}{}{}}
\LANGes{
\Question{
Juan es capaz de segar un prado en \( 12 \) horas. Jorge tiene un cortacespéd mejor y el mismo prado lo segaría en \( 8 \) horas. Han acordado que Juan va a empezar a segar solo y que Jorge se unirá a él más tarde para que el tiempo total de siega sea \( 9 \) horas. ¿Cuánto tiempo cortarán juntos?}
{\( 2 \) horas}{\( 7 \) horas}{\( 6 \) horas}{\( 3 \) horas}{}{}}
\End

\Data\ID{1003060501}
\Updated{2020-07-15 03:07:12}
\Level{C}
\SubArea{Systems of equations and inequalities}
\LANGen{
\Question{
Which of the following systems of equations has no solution?}
{\(\begin{aligned}
2x-y+3z&=2 \\
6x-3y+9z&=4 \\
x+y+z&=1
\end{aligned} \)}{\(\begin{aligned}
2x-y+3z&=2 \\
6x-3y+9z&=6 \\
x+y+z&=1 \\
\end{aligned} \)}{\(\begin{aligned}
2x-y+3z&=2 \\
6x-2y+z&=4 \\
x+y+z&=1
\end{aligned} \)}{\(\begin{aligned}
2x-y+3z&=2 \\
6x-3y+9z&=6 \\
4x-4y+6z&=4
\end{aligned} \)}{}{}}
\LANGpl{
\Question{
Który z podanych układów równań nie ma rozwiązania?}
{\(\begin{aligned}
2x-y+3z&=2 \\
6x-3y+9z&=4 \\
x+y+z&=1
\end{aligned} \)}{\(\begin{aligned}
2x-y+3z&=2 \\
6x-3y+9z&=6 \\
x+y+z&=1 \\
\end{aligned} \)}{\(\begin{aligned}
2x-y+3z&=2 \\
6x-2y+z&=4 \\
x+y+z&=1
\end{aligned} \)}{\(\begin{aligned}
2x-y+3z&=2 \\
6x-3y+9z&=6 \\
4x-4y+6z&=4
\end{aligned} \)}{}{}}
\LANGcs{
\Question{
Která z následujících soustav rovnic nemá žádné řešení?}
{\(\begin{aligned}
2x-y+3z&=2 \\
6x-3y+9z&=4 \\
x+y+z&=1
\end{aligned} \)}{\(\begin{aligned}
2x-y+3z&=2 \\
6x-3y+9z&=6 \\
x+y+z&=1 \\
\end{aligned} \)}{\(\begin{aligned}
2x-y+3z&=2 \\
6x-2y+z&=4 \\
x+y+z&=1
\end{aligned} \)}{\(\begin{aligned}
2x-y+3z&=2 \\
6x-3y+9z&=6 \\
4x-4y+6z&=4
\end{aligned} \)}{}{}}
\LANGsk{
\Question{
Ktorá z nasledujúcich sústav rovníc nemá žiadne riešenie?}
{\(\begin{aligned}
2x-y+3z&=2 \\
6x-3y+9z&=4 \\
x+y+z&=1
\end{aligned} \)}{\(\begin{aligned}
2x-y+3z&=2 \\
6x-3y+9z&=6 \\
x+y+z&=1 \\
\end{aligned} \)}{\(\begin{aligned}
2x-y+3z&=2 \\
6x-2y+z&=4 \\
x+y+z&=1
\end{aligned} \)}{\(\begin{aligned}
2x-y+3z&=2 \\
6x-3y+9z&=6 \\
4x-4y+6z&=4
\end{aligned} \)}{}{}}
\LANGes{
\Question{
¿Cuál de los siguientes sistemas de ecuaciones no tiene solución?}
{\(\begin{aligned}
2x-y+3z&=2 \\
6x-3y+9z&=4 \\
x+y+z&=1
\end{aligned} \)}{\(\begin{aligned}
2x-y+3z&=2 \\
6x-3y+9z&=6 \\
x+y+z&=1 \\
\end{aligned} \)}{\(\begin{aligned}
2x-y+3z&=2 \\
6x-2y+z&=4 \\
x+y+z&=1
\end{aligned} \)}{\(\begin{aligned}
2x-y+3z&=2 \\
6x-3y+9z&=6 \\
4x-4y+6z&=4
\end{aligned} \)}{}{}}
\End

\Data\ID{1003060502}
\Updated{2020-07-15 03:07:12}
\Level{C}
\SubArea{Systems of equations and inequalities}
\LANGen{
\Question{
The system of equations is given by:
\[ \begin{aligned}
x+y-2z&=0, \\
x+2y+3z&=0, \\
-2x+y+z&=2.
\end{aligned} \]
To which of the following systems is it equivalent?  (Note:  An algorithm for solving a system of linear equations by transformation the system into this form (row echelon form) is known as Gaussian elimination or as row reduction.)}
{\( \begin{aligned}
x+y-2z&=0 \\
y+5z&=0 \\
18z&=-2
\end{aligned} \)}{\( \begin{aligned}
x+y-2z&=0 \\
y-5z&=0 \\
12z&=2
\end{aligned} \)}{\( \begin{aligned}
x+y-2z&=0 \\
y+5z&=0 \\
18z&=2
\end{aligned} \)}{\( \begin{aligned}
x+y-2z&=0 \\
y+z&=0 \\
6z&=2
\end{aligned} \)}{}{}}
\LANGpl{
\Question{
Dany jest układ równań:
\[ \begin{aligned}
x+y-2z&=0, \\
x+2y+3z&=0, \\
-2x+y+z&=2. 
\end{aligned} \]
Który z poniższych systemów jest równoważny? (Uwaga: Algorytm rozwiązywania układu równań liniowych przez przekształcenie systemu w tę postać (forma rzutu rzędów) jest znany jako eliminacja Gaussa lub redukcja rzędów.)}
{\( \begin{aligned}
x+y-2z&=0 \\
y+5z&=0 \\
18z&=-2
\end{aligned} \)}{\( \begin{aligned}
x+y-2z&=0 \\
y-5z&=0 \\
12z&=2
\end{aligned} \)}{\( \begin{aligned}
x+y-2z&=0 \\
y+5z&=0 \\
18z&=2
\end{aligned} \)}{\( \begin{aligned}
x+y-2z&=0 \\
y+z&=0 \\
6z&=2
\end{aligned} \)}{}{}}
\LANGcs{
\Question{
Je dána soustava rovnic:
\[ \begin{aligned}
x+y-2z&=0, \\
x+2y+3z&=0, \\
-2x+y+z&=2. 
\end{aligned} \]
Která z následujících soustav je s ní ekvivalentní? (Poznámka: Způsob řešení soustav lineárních rovnic úpravou soustavy na tento tvar označujeme jako Gaussovu eliminační metodu.)}
{\( \begin{aligned}
x+y-2z&=0 \\
y+5z&=0 \\
18z&=-2
\end{aligned} \)}{\( \begin{aligned}
x+y-2z&=0 \\
y-5z&=0 \\
12z&=2
\end{aligned} \)}{\( \begin{aligned}
x+y-2z&=0 \\
y+5z&=0 \\
18z&=2
\end{aligned} \)}{\( \begin{aligned}
x+y-2z&=0 \\
y+z&=0 \\
6z&=2
\end{aligned} \)}{}{}}
\LANGsk{
\Question{
Je daná sústava rovníc:
\[ \begin{aligned}
x+y-2z&=0, \\
x+2y+3z&=0, \\
-2x+y+z&=2. 
\end{aligned} \]
Ktorá z nasledujúcich sústav je s ňou ekvivalentná?  (Poznámka:  Spôsob riešenia sústav lineárnych rovníc úpravou sústavy na tento tvar označujeme ako Gaussovu eliminačnú metódu.)}
{\( \begin{aligned}
x+y-2z&=0 \\
y+5z&=0 \\
18z&=-2
\end{aligned} \)}{\( \begin{aligned}
x+y-2z&=0 \\
y-5z&=0 \\
12z&=2
\end{aligned} \)}{\( \begin{aligned}
x+y-2z&=0 \\
y+5z&=0 \\
18z&=2
\end{aligned} \)}{\( \begin{aligned}
x+y-2z&=0 \\
y+z&=0 \\
6z&=2
\end{aligned} \)}{}{}}
\LANGes{
\Question{
Dado el sistema:
\[ \begin{aligned}
x+y-2z&=0, \\
x+2y+3z&=0, \\
-2x+y+z&=2.
\end{aligned} \]
¿A cuál de los siguientes sistemas es equivalente?  (Nota:  un algoritmo que sirve para resolver un sistema de ecuaciones lineales mediante la transformación del sistema en esta forma se conoce como el método de reducción de Gauss.)}
{\( \begin{aligned}
x+y-2z&=0 \\
y+5z&=0 \\
18z&=-2
\end{aligned} \)}{\( \begin{aligned}
x+y-2z&=0 \\
y-5z&=0 \\
12z&=2
\end{aligned} \)}{\( \begin{aligned}
x+y-2z&=0 \\
y+5z&=0 \\
18z&=2
\end{aligned} \)}{\( \begin{aligned}
x+y-2z&=0 \\
y+z&=0 \\
6z&=2
\end{aligned} \)}{}{}}
\End

\Data\ID{1003060503}
\Updated{2020-07-15 03:07:12}
\Level{C}
\SubArea{Systems of equations and inequalities}
\LANGen{
\Question{
The system of equations is given by:
\[ \begin{aligned}
x-y-z&=0, \\
2x-y+3z&=1, \\
-3x+2y+z&=2.
\end{aligned} \]
To which of the following systems is it equivalent?  (Note: An algorithm for solving a system of linear equations by transformation the  system  into this form (row echelon form) is known as Gaussian elimination or as row reduction.)}
{\( \begin{aligned}
x-y-z&=0 \\
y+5z&=1 \\
3z&=3
\end{aligned} \)}{\( \begin{aligned}
x-y-z&=0 \\
y+5z&=-1 \\
3z&=-1
\end{aligned} \)}{\( \begin{aligned}
x-y-z&=0 \\
-3y-z&=-1 \\
5z&=5
\end{aligned} \)}{\( \begin{aligned}
x-y-z&=0 \\
-3y-z&=-1 \\
5z&=-7
\end{aligned} \)}{}{}}
\LANGpl{
\Question{
Dany jest układ równań:
\[ \begin{aligned}
x-y-z&=0, \\
2x-y+3z&=1, \\
-3x+2y+z&=2.
\end{aligned} \]
Który z poniższych systemów jest równoważny? (Uwaga: Algorytm rozwiązywania układu równań liniowych przez przekształcenie systemu w tę postać (forma rzutu rzędów) jest znany jako eliminacja Gaussa lub redukcja rzędów.)}
{\( \begin{aligned}
x-y-z&=0 \\
y+5z&=1 \\
3z&=3
\end{aligned} \)}{\( \begin{aligned}
x-y-z&=0 \\
y+5z&=-1 \\
3z&=-1
\end{aligned} \)}{\( \begin{aligned}
x-y-z&=0 \\
-3y-z&=-1 \\
5z&=5
\end{aligned} \)}{\( \begin{aligned}
x-y-z&=0 \\
-3y-z&=-1 \\
5z&=-7
\end{aligned} \)}{}{}}
\LANGcs{
\Question{
Je dána soustava rovnic:
\[ \begin{aligned}
x-y-z&=0, \\
2x-y+3z&=1, \\
-3x+2y+z&=2.
\end{aligned} \]
Která z následujících soustav je s ní ekvivalentní? (Poznámka: Způsob řešení soustav lineárních rovnic úpravou soustavy na tento  tvar označujeme jako Gaussovu eliminační metodu.)}
{\( \begin{aligned}
x-y-z&=0 \\
y+5z&=1 \\
3z&=3
\end{aligned} \)}{\( \begin{aligned}
x-y-z&=0 \\
y+5z&=-1 \\
3z&=-1
\end{aligned} \)}{\( \begin{aligned}
x-y-z&=0 \\
-3y-z&=-1 \\
5z&=5
\end{aligned} \)}{\( \begin{aligned}
x-y-z&=0 \\
-3y-z&=-1 \\
5z&=-7
\end{aligned} \)}{}{}}
\LANGsk{
\Question{
Je daná sústava rovníc:
\[ \begin{aligned}
x-y-z&=0, \\
2x-y+3z&=1, \\
-3x+2y+z&=2.
\end{aligned} \]
Ktorá z nasledujúcich sústav je s ňou ekvivalentná? (Poznámka: Spôsob riešenia sústav lineárnych rovníc úpravou sústavy na tento tvar označujeme ako  Gaussovu eliminačnú metódu.)}
{\( \begin{aligned}
x-y-z&=0 \\
y+5z&=1 \\
3z&=3
\end{aligned} \)}{\( \begin{aligned}
x-y-z&=0 \\
y+5z&=-1 \\
3z&=-1
\end{aligned} \)}{\( \begin{aligned}
x-y-z&=0 \\
-3y-z&=-1 \\
5z&=5
\end{aligned} \)}{\( \begin{aligned}
x-y-z&=0 \\
-3y-z&=-1 \\
5z&=-7
\end{aligned} \)}{}{}}
\LANGes{
\Question{
Dado el sistema:
\[ \begin{aligned}
x-y-z&=0, \\
2x-y+3z&=1, \\
-3x+2y+z&=2.
\end{aligned} \]
¿A cuál de los siguientes sistemas es equivalente? (Nota: un algoritmo que sirve para resolver un sistema de ecuaciones lineales mediante la transformación del sistema en esta forma se conoce como el método de reducción de Gauss).}
{\( \begin{aligned}
x-y-z&=0 \\
y+5z&=1 \\
3z&=3
\end{aligned} \)}{\( \begin{aligned}
x-y-z&=0 \\
y+5z&=-1 \\
3z&=-1
\end{aligned} \)}{\( \begin{aligned}
x-y-z&=0 \\
-3y-z&=-1 \\
5z&=5
\end{aligned} \)}{\( \begin{aligned}
x-y-z&=0 \\
-3y-z&=-1 \\
5z&=-7
\end{aligned} \)}{}{}}
\End

\Data\ID{1003060504}
\Updated{2020-07-15 03:07:12}
\Level{C}
\SubArea{Systems of equations and inequalities}
\LANGen{
\Question{
Four systems of equations are given. How many of the given systems have infinitely 
many solutions?
\[
\begin{array}{c|c}
\text{\( \begin{aligned}
4x-6y+10z&=8    \\
-2x+3y-5z&=4     \\
x+y+z&=1  
\end{aligned}\)}&
\text{\(
 \begin{aligned}
4x-6y+10z&=8\\
6x-9y+15z&=12\\
 x+y+z&=1\\
\end{aligned}\)}
\\\hline
\text{\(\begin{aligned}
4x-6y+10z&=8\\
 -2x+3y+5z&=4\\
 x+y+z&=1\\
\end{aligned}\)}&
\text{\(
\begin{aligned}
 x+y+z&=1 \\
2x+2y+2z&=2 \\
 -\frac x2-\frac y2-\frac z2&=-\frac12
\end{aligned}\)}
\end{array}
\]}
{\( 2 \)}{\( 1 \)}{\( 3 \)}{\( 4 \)}{}{}}
\LANGpl{
\Question{
Dane są cztery układy równań. Ile z podanych układów mają nieskończenie wiele rozwiązań?
\[
\begin{array}{c|c}
\text{\( \begin{aligned}
4x-6y+10z&=8    \\
-2x+3y-5z&=4     \\
x+y+z&=1  
\end{aligned}\)}&
\text{\(
 \begin{aligned}
4x-6y+10z&=8\\
6x-9y+15z&=12\\
 x+y+z&=1\\
\end{aligned}\)}
\\\hline
\text{\(\begin{aligned}
4x-6y+10z&=8\\
 -2x+3y+5z&=4\\
 x+y+z&=1\\
\end{aligned}\)}&
\text{\(
\begin{aligned}
 x+y+z&=1 \\
2x+2y+2z&=2 \\
 -\frac x2-\frac y2-\frac z2&=-\frac12
\end{aligned}\)}
\end{array}
\]}
{\( 2 \)}{\( 1 \)}{\( 3 \)}{\( 4 \)}{}{}}
\LANGcs{
\Question{
Jsou dány čtyři soustavy rovnic. Kolik z uvedených soustav má  nekonečně mnoho řešení?
\[
\begin{array}{c|c}
\text{\( \begin{aligned}
4x-6y+10z&=8    \\
-2x+3y-5z&=4     \\
x+y+z&=1  
\end{aligned}\)}&
\text{\(
 \begin{aligned}
4x-6y+10z&=8\\
6x-9y+15z&=12\\
 x+y+z&=1\\
\end{aligned}\)}
\\\hline
\text{\(\begin{aligned}
4x-6y+10z&=8\\
 -2x+3y+5z&=4\\
 x+y+z&=1\\
\end{aligned}\)}&
\text{\(
\begin{aligned}
 x+y+z&=1 \\
2x+2y+2z&=2 \\
 -\frac x2-\frac y2-\frac z2&=-\frac12
\end{aligned}\)}
\end{array}
\]}
{\( 2 \)}{\( 1 \)}{\( 3 \)}{\( 4 \)}{}{}}
\LANGsk{
\Question{
Sú dané štyri sústavy rovníc. Koľko z uvedených sústav má nekonečne veľa riešení?
\[
\begin{array}{c|c}
\text{\( \begin{aligned}
4x-6y+10z&=8    \\
-2x+3y-5z&=4     \\
x+y+z&=1  
\end{aligned}\)}&
\text{\(
 \begin{aligned}
4x-6y+10z&=8\\
6x-9y+15z&=12\\
 x+y+z&=1\\
\end{aligned}\)}
\\\hline
\text{\(\begin{aligned}
4x-6y+10z&=8\\
 -2x+3y+5z&=4\\
 x+y+z&=1\\
\end{aligned}\)}&
\text{\(
\begin{aligned}
 x+y+z&=1 \\
2x+2y+2z&=2 \\
 -\frac x2-\frac y2-\frac z2&=-\frac12
\end{aligned}\)}
\end{array}
\]}
{\( 2 \)}{\( 1 \)}{\( 3 \)}{\( 4 \)}{}{}}
\LANGes{
\Question{
Dados cuatro sistemas de ecuaciones. ¿Cuántos de los sistemas dados tienen infinitas soluciones?
\[
\begin{array}{c|c}
\text{\( \begin{aligned}
4x-6y+10z&=8    \\
-2x+3y-5z&=4     \\
x+y+z&=1  
\end{aligned}\)}&
\text{\(
 \begin{aligned}
4x-6y+10z&=8\\
6x-9y+15z&=12\\
 x+y+z&=1\\
\end{aligned}\)}
\\\hline
\text{\(\begin{aligned}
4x-6y+10z&=8\\
 -2x+3y+5z&=4\\
 x+y+z&=1\\
\end{aligned}\)}&
\text{\(
\begin{aligned}
 x+y+z&=1 \\
2x+2y+2z&=2 \\
 -\frac x2-\frac y2-\frac z2&=-\frac12
\end{aligned}\)}
\end{array}
\]}
{\( 2 \)}{\( 1 \)}{\( 3 \)}{\( 4 \)}{}{}}
\End

\Data\ID{1003083004}
\Updated{2020-07-15 03:07:12}
\Level{C}
\SubArea{Systems of equations and inequalities}
\LANGen{
\Question{
Find the real number \( a \) such that the following system of equations has no solution. 
\[ \begin{aligned} \frac25x-\frac a4y&=4 \\ -\frac x4 + \frac{5y}8&=\frac52 \end{aligned}\]}
{\( 4 \)}{\( -\frac52 \)}{Such a real number \( a \) does not exist.}{\( -4 \)}{}{}}
\LANGpl{
\Question{
Znajdź taką liczbę rzeczywistą \( a \), dla której  podany układ nie ma rozwiązania.
\[ \begin{aligned} \frac25x-\frac a4y&=4 \\ -\frac x4 + \frac{5y}8&=\frac52 \end{aligned}\]}
{\( 4 \)}{\( -\frac52 \)}{Taka liczba rzeczywista \( a \) nie istnieje.}{\( -4 \)}{}{}}
\LANGcs{
\Question{
Jakou hodnotu musí mít reálný koeficient \( a \), aby následující soustava rovnic neměla řešení?
\[ \begin{aligned} \frac25x-\frac a4y&=4 \\ -\frac x4 + \frac{5y}8&=\frac52 \end{aligned}\]}
{\( 4 \)}{\( -\frac52 \)}{Takové reálné číslo \( a \) neexistuje.}{\( -4 \)}{}{}}
\LANGsk{
\Question{
Akú hodnotu musí mať reálny koeficient \( a \), aby nasledujúca sústava rovníc nemala riešenie?   
\[ \begin{aligned} \frac25x-\frac a4y&=4 \\ -\frac x4 + \frac{5y}8&=\frac52 \end{aligned}\]}
{\( 4 \)}{\( -\frac52 \)}{Také reálne číslo \( a \) neexistuje.}{\( -4 \)}{}{}}
\LANGes{
\Question{
Halla el número real \( a \) de modo que el siguiente sistema de ecuaciones no tenga solución. 
\[ \begin{aligned} \frac25x-\frac a4y&=4 \\ -\frac x4 + \frac{5y}8&=\frac52 \end{aligned}\]}
{\( 4 \)}{\( -\frac52 \)}{El número real así \( a \) no existe.}{\( -4 \)}{}{}}
\End

\Data\ID{1003085406}
\Updated{2020-07-15 03:07:12}
\Level{C}
\SubArea{Systems of equations and inequalities}
\LANGen{
\Question{
A diagonal of a rectangle has length of \( 30\,\mathrm{cm} \) and the perimeter of the rectangle is \( 84\,\mathrm{cm} \). Find the difference of the length and the width of the rectangle in centimeters.}
{\( 6 \)}{\( 8 \)}{\( 9 \)}{\( 10 \)}{}{}}
\LANGpl{
\Question{
Przekątna prostokąta ma długość \( 30\,\mathrm{cm} \), a obwód tego prostokąta wynosi \( 84\,\mathrm{cm} \). Oblicz różnicę długości i szerokości tego prostokąta w centymetrach.}
{\( 6 \)}{\( 8 \)}{\( 9 \)}{\( 10 \)}{}{}}
\LANGcs{
\Question{
Úhlopříčka obdélníku má délku \( 30\,\mathrm{cm} \), jeho obvod měří  \( 84\,\mathrm{cm} \). Určete rozdíl délek stran obdélníku v centimetrech.}
{\( 6 \)}{\( 8 \)}{\( 9 \)}{\( 10 \)}{}{}}
\LANGsk{
\Question{
Uhlopriečka obdĺžnika má dĺžku \( 30\,\mathrm{cm} \) a obvod obdĺžnika je \( 84\,\mathrm{cm} \). Určte rozdiel dĺžok strán obdĺžnika v centimetroch.}
{\( 6 \)}{\( 8 \)}{\( 9 \)}{\( 10 \)}{}{}}
\LANGes{
\Question{
La diagonal de un rectángulo mide \( 30\,\mathrm{cm} \) y el perímetro \( 84\,\mathrm{cm} \). Halla la diferencia de los lados en centímetros.}
{\( 6 \)}{\( 8 \)}{\( 9 \)}{\( 10 \)}{}{}}
\End

\Data\ID{1003160801}
\Updated{2020-07-15 03:07:21}
\Level{C}
\SubArea{Systems of equations and inequalities}
\LANGen{
\Question{
Use the substitution method to find the solution \( [x;y] \) of the following system of equations.
\[ \begin{aligned}
\frac2{x+4}-\frac1{2-y}=-6 \\
\frac1{x+4}+\frac5{2-y}=8
\end{aligned} \]}
{\( \left[-\frac92;\frac32\right] \)}{\( [-2;2] \)}{\( [2;10] \)}{\( \left[-\frac92;3\right] \)}{}{}}
\LANGpl{
\Question{
Zastosuj metodę substytucji, aby znaleźć rozwiązanie \( [x;y] \) podanego układu nierówności. 
\[ \begin{aligned}
\frac2{x+4}-\frac1{2-y}=-6 \\
\frac1{x+4}+\frac5{2-y}=8
\end{aligned} \]}
{\( \left[-\frac92;\frac32\right] \)}{\( [-2;2] \)}{\( [2;10] \)}{\( \left[-\frac92;3\right] \)}{}{}}
\LANGcs{
\Question{
Řešením následující soustavy rovnic je uspořádaná dvojice \( [x;y] \). Pomocí vhodné substituce najděte toto řešení.
\[ \begin{aligned}
\frac2{x+4}-\frac1{2-y}=-6 \\
\frac1{x+4}+\frac5{2-y}=8
\end{aligned} \]}
{\( \left[-\frac92;\frac32\right] \)}{\( [-2;2] \)}{\( [2;10] \)}{\( \left[-\frac92;3\right] \)}{}{}}
\LANGsk{
\Question{
Pomocou vhodnej substitúcie nájdite riešenie \( [x;y] \) nasledujúcej sústavy rovníc.
\[ \begin{aligned}
\frac2{x+4}-\frac1{2-y}=-6 \\
\frac1{x+4}+\frac5{2-y}=8
\end{aligned} \]}
{\( \left[-\frac92;\frac32\right] \)}{\( [-2;2] \)}{\( [2;10] \)}{\( \left[-\frac92;3\right] \)}{}{}}
\LANGes{
\Question{
Utiliza el método de sustitución para hallar la solución \( [x;y] \) del siguiente sistema de ecuaciones.
\[ \begin{aligned}
\frac2{x+4}-\frac1{2-y}=-6 \\
\frac1{x+4}+\frac5{2-y}=8
\end{aligned} \]}
{\( \left[-\frac92;\frac32\right] \)}{\( [-2;2] \)}{\( [2;10] \)}{\( \left[-\frac92;3\right] \)}{}{}}
\End

\Data\ID{1003160802}
\Updated{2020-07-15 03:07:21}
\Level{C}
\SubArea{Systems of equations and inequalities}
\LANGen{
\Question{
Use the substitution method to find the solution \( [x;y] \) of the following system of equations.
\[ \begin{aligned}
\frac{2x}{x+3}-3\cdot\frac{y+2}y=2 \\
\frac x{x+3}+2\cdot\frac{y+2}y=8
\end{aligned} \]}
{\( [-4;2] \)}{\( [4;2] \)}{\( [2;-4] \)}{\( [-1;2] \)}{}{}}
\LANGpl{
\Question{
Zastosuj metodę substytucji, aby znaleźć rozwiązanie \( [x;y] \) podanego układu nierówności.
\[ \begin{aligned}
\frac{2x}{x+3}-3\cdot\frac{y+2}y=2 \\
\frac x{x+3}+2\cdot\frac{y+2}y=8
\end{aligned} \]}
{\( [-4;2] \)}{\( [4;2] \)}{\( [2;-4] \)}{\( [-1;2] \)}{}{}}
\LANGcs{
\Question{
Řešením následující soustavy rovnic je uspořádaná dvojice \( [x;y] \). Pomocí vhodné substituce najděte toto řešení.
\[ \begin{aligned}
\frac{2x}{x+3}-3\cdot\frac{y+2}y=2 \\
\frac x{x+3}+2\cdot\frac{y+2}y=8
\end{aligned} \]}
{\( [-4;2] \)}{\( [4;2] \)}{\( [2;-4] \)}{\( [-1;2] \)}{}{}}
\LANGsk{
\Question{
Pomocou vhodnej substitúcie nájdite riešenie \( [x;y] \) nasledujúcej sústavy rovníc.
\[ \begin{aligned}
\frac{2x}{x+3}-3\cdot\frac{y+2}y=2 \\
\frac x{x+3}+2\cdot\frac{y+2}y=8
\end{aligned} \]}
{\( [-4;2] \)}{\( [4;2] \)}{\( [2;-4] \)}{\( [-1;2] \)}{}{}}
\LANGes{
\Question{
Utiliza el método de sustitución para hallar la solución \( [x;y] \) del siguiente sistema de ecuaciones.
\[ \begin{aligned}
\frac{2x}{x+3}-3\cdot\frac{y+2}y=2 \\
\frac x{x+3}+2\cdot\frac{y+2}y=8
\end{aligned} \]}
{\( [-4;2] \)}{\( [4;2] \)}{\( [2;-4] \)}{\( [-1;2] \)}{}{}}
\End

\Data\ID{1003160803}
\Updated{2020-07-15 03:07:21}
\Level{C}
\SubArea{Systems of equations and inequalities}
\LANGen{
\Question{
Use the substitution method to find the solution \( [x;y] \) of the following system of equations.
\[ \begin{aligned}
\frac{x+y}x+\frac1{x+y}=1 \\
\frac{2\cdot(x+y)}x-\frac1{x+y}=-7
\end{aligned} \]}
{\( \left[-\frac16;\frac12\right] \)}{\( [-2;3] \)}{\( \left[-\frac12;-\frac12\right] \)}{\( \left[\frac12;\frac3{-2}\right] \)}{}{}}
\LANGpl{
\Question{
Zastosuj metodę substytucji, aby znaleźć rozwiązanie \( [x;y] \) podanego układu równań.
\[ \begin{aligned}
\frac{x+y}x+\frac1{x+y}=1 \\
\frac{2\cdot(x+y)}x-\frac1{x+y}=-7
\end{aligned} \]}
{\( \left[-\frac16;\frac12\right] \)}{\( [-2;3] \)}{\( \left[-\frac12;-\frac12\right] \)}{\( \left[\frac12;\frac3{-2}\right] \)}{}{}}
\LANGcs{
\Question{
Řešením následující soustavy rovnic je uspořádaná dvojice  \( [x;y] \). Pomocí vhodné substituce najděte toto řešení.
\[ \begin{aligned}
\frac{x+y}x+\frac1{x+y}=1 \\
\frac{2\cdot(x+y)}x-\frac1{x+y}=-7
\end{aligned} \]}
{\( \left[-\frac16;\frac12\right] \)}{\( [-2;3] \)}{\( \left[-\frac12;-\frac12\right] \)}{\( \left[\frac12;\frac3{-2}\right] \)}{}{}}
\LANGsk{
\Question{
Pomocou vhodnej substitúcie nájdite riešenie \( [x;y] \) nasledujúcej sústavy rovníc.
\[ \begin{aligned}
\frac{x+y}x+\frac1{x+y}=1 \\
\frac{2\cdot(x+y)}x-\frac1{x+y}=-7
\end{aligned} \]}
{\( \left[-\frac16;\frac12\right] \)}{\( [-2;3] \)}{\( \left[-\frac12;-\frac12\right] \)}{\( \left[\frac12;\frac3{-2}\right] \)}{}{}}
\LANGes{
\Question{
Utiliza el método de sustitución para hallar la solución \( [x;y] \) del siguiente sistema de ecuaciones.
\[ \begin{aligned}
\frac{x+y}x+\frac1{x+y}=1 \\
\frac{2\cdot(x+y)}x-\frac1{x+y}=-7
\end{aligned} \]}
{\( \left[-\frac16;\frac12\right] \)}{\( [-2;3] \)}{\( \left[-\frac12;-\frac12\right] \)}{\( \left[\frac12;\frac3{-2}\right] \)}{}{}}
\End

\Data\ID{1103034507}
\Updated{2020-08-23 08:08:26}
\Level{C}
\SubArea{Systems of equations and inequalities}
\IMG{1103034507b.pdf}
\LANGen{
\Question{
Consider a balance scales comprising of beam with unequal length of arms where the fulcrum is very close to one end of the beam. (Such scales are called steelyard. For instance, it is often used to weigh a catch in fisheries.) The load is hanged on the shorter arm, while the balance about the fulcrum is obtained by sliding the counterweight along the longer arm. (See the picture.) Suppose the distance of the load hanging point from the fulcrum is fixed at \( 5\,\mathrm{cm} \). If the weight of the load is \( 80\,\mathrm{N} \), the balance is achieved as the counterweight is moved to the very end of the longer arm. If the weight of the load is \( 60\,\mathrm{N} \), the balance is achieved when the counterweight is moved to the distance of \( 30\,\mathrm{cm} \) from the fulcrum. What is the length of the beam? 
\[ \]
Hint: The steelyard is based on the law of the lever. For balanced lever is: \( F_1\cdot a=F_2\cdot b \), where \( F_1 \) is the weight of the load in the distance \( a \) from the fulcrum and \( F_2 \) is the weight of the counterweight in the distance \( b \) from the fulcrum.}
{\( 45\,\mathrm{cm} \)}{\( 54\,\mathrm{cm} \)}{\( 40\,\mathrm{cm} \)}{\( 35\,\mathrm{cm} \)}{}{}}
\LANGpl{
\Question{
Rozważmy wagę bilansową składającą się z belki o nierównej długości ramion, gdzie punkt podparcia jest bardzo blisko jednego końca belki. (Takie łuski są nazywane bezkręgowymi, na przykład często są używane do ważenia połowu w łowiskach.) Ładunek jest zawieszany na krótszym ramieniu, podczas gdy równowaga wokół punktu oparcia jest uzyskiwana przez przesuwanie przeciwwagi wzdłuż dłuższego ramienia. (Zobacz obrazek.) Załóżmy, że odległość punktu zawieszenia ładunku od punktu podparcia wynosi \( 5\, \mathrm {cm} \). Jeśli ciężar ładunku wynosi \( 80\, \mathrm {N} \), równowaga zostaje osiągnięta, gdy przeciwwaga zostanie przesunięta na sam koniec dłuższego ramienia. Jeśli ciężar ładunku wynosi \( 60\,\mathrm {N} \), równowaga zostaje osiągnięta, gdy przeciwwaga zostanie przesunięta na odległość \(30\, \mathrm {cm} \) od punktu podparcia. Jaka jest długość belki?

\[ \]
Wskazówka: Bezmian opiera się na prawie dźwigni.  Dla dźwigni zrównoważonej jest: \( F_1\cdot a=F_2\cdot b \) gdzie \(F_1\) jest wagą ładunku w odległości \(a\) od punktu podparcia, a \(F_2\) jest ciężarem  przeciwwagi w odległości \(b\) od punktu podparcia.}
{\( 45\,\mathrm{cm} \)}{\( 54\,\mathrm{cm} \)}{\( 40\,\mathrm{cm} \)}{\( 35\,\mathrm{cm} \)}{}{}}
\LANGcs{
\Question{
Máme nerovnoramenné váhy s proměnnou délkou jednoho z ramen (takové váhy se často označují jako přezmeny a využívají se např. v rybářství pro vážení vylovených ryb). Na jedné straně  vodorovné "tyče" (vahadla) je zavěšeno břemeno a někde na druhé straně je závaží.  Stačí jediné závaží, které se posouvá po delším rameni páky tak dlouho, až nastane rovnováha. Břemeno se zavěšuje vždy \( 5\,\mathrm{cm} \) od bodu závěsu vahadla (viz obrázek). Má-li břemeno tíhu \( 80\,\mathrm{N} \), dosáhneme rovnováhy, když posuneme vyrovnávací závaží až na konec vahadla. Má-li břemeno tíhu \( 60\,\mathrm{N} \), rovnováha nastane, když závaží bude od bodu závěsu vzdáleno \( 30\,\mathrm{cm} \). Jak dlouhé je vahadlo?
\[ \]
Nápověda: Váhy představují dvojzvratnou páku. Platí podmínka rovnováhy: \( F_1\cdot a=F_2\cdot b \), kde \( F_1 \) je tíha břemene ve vzdálenosti \( a \) od bodu závěsu a \( F_2 \) je tíha závaží ve vzdálenosti \( b \) od bodu závěsu.}
{\( 45\,\mathrm{cm} \)}{\( 54\,\mathrm{cm} \)}{\( 40\,\mathrm{cm} \)}{\( 35\,\mathrm{cm} \)}{}{}}
\LANGsk{
\Question{
Máme nerovnoramenné váhy s nerovnakou dĺžkou ramien. (Také váhy sa nazývajú minciere a využívajú sa napr. v rybárstve pre váženie vylovených rýb.) Na jednej strane vodorovnej tyče (vahadla) je zavesené teleso a niekde na druhej strane je závažie, ktoré sa posúva po dlhšom ramene páky tak dlho, až nastane rovnováha. (Pozri obrázok.) Teleso sa zavesí  \( 5\,\mathrm{cm} \) od bodu opory vahadla. Ak má teleso tiaž \( 80\,\mathrm{N} \), dosiahneme rovnováhu, keď posunieme vyrovnávacie závažie až na koniec vahadla. Ak má teleso tiaž \( 60\,\mathrm{N} \), rovnováha nastane, keď závažie bude vzdialené \( 30\,\mathrm{cm} \) od oporného bodu. Aké dlhé je vahadlo? 
\[ \]
Nápoveda: Mincier je založený na princípe páky. Pre vahadlo platí: \( F_1\cdot a=F_2\cdot b \), kde \( F_1 \) je tiaž telesa vo vzdialenosti \( a \) od oporného bodu a \( F_2 \) je tiaž závažia vo vzdialenosti  \( b \) od oporného bodu.}
{\( 45\,\mathrm{cm} \)}{\( 54\,\mathrm{cm} \)}{\( 40\,\mathrm{cm} \)}{\( 35\,\mathrm{cm} \)}{}{}}
\LANGes{
\Question{
Considera una balanza compuesta por una viga con brazos de longitud desigual donde el punto de apoyo está muy cerca de un extremo de la viga. La carga se cuelga en el brazo más corto, mientras tanto el equilibrio sobre el punto de apoyo se obtiene deslizando el contrapeso a lo largo del brazo más largo. (Mira la imagen.) Supongamos que la distancia del punto de suspensión de la carga desde el punto de apoyo se fija en \( 5\,\mathrm{cm} \). Si el peso de la carga es \( 80\,\mathrm{N} \), el equilibrio se logra a medida que el contrapeso se mueve hasta el final del brazo más largo. Si el peso de la carga es \( 60\,\mathrm{N} \), el equilibrio se logra cuando el contrapeso se mueve a una distancia de \( 30\,\mathrm{cm} \) desde el punto de apoyo. ¿Cuál es la longitud de la viga? 
\[ \]
Sugerencia: La romana se basa en la ley de la palanca. Para la palanca equilibrada vale: \( F_1\cdot a=F_2\cdot b \), donde \( F_1 \) es el peso de la carga a una distancia \( a \) desde el punto de apoyo y \( F_2 \) es el peso del contrapeso a una distancia \( b \) desde el punto de apoyo.}
{\( 45\,\mathrm{cm} \)}{\( 54\,\mathrm{cm} \)}{\( 40\,\mathrm{cm} \)}{\( 35\,\mathrm{cm} \)}{}{}}
\End

\Data\ID{1103085403}
\Updated{2020-07-15 03:07:12}
\Level{C}
\SubArea{Systems of equations and inequalities}
\IMG{1103157103.pdf}
\LANGen{
\Question{
Two resistors of unknown resistances \( R_1 \) and \( R_2 \), where \( R_1 < R_2 \) are connected in series (Figure A) and total resistance of the circuit is \( R_S=64\,\Omega \). If the resistors are connected in parallel (Figure B), the total resistance \( R_P=15\,\Omega \). Find \( R_1 \).}
{\( 24 \)}{\( 22 \)}{\( 12 \)}{\( 15 \)}{}{}}
\LANGpl{
\Question{
Dwa rezystory o oporach \( R_1 \) i \( R_2 \), gdzie \( R_1 < R_2 \) są połączone szeregowo (obraz A) a całkowity opór obwodu jest równy \( R_S=64\,\Omega \). Jeśli rezystory połączymy równolegle (obraz B), całkowity opór wynosi \( R_P=15\,\Omega \). Wyznacz \( R_1 \).}
{\( 24 \)}{\( 22 \)}{\( 12 \)}{\( 15 \)}{}{}}
\LANGcs{
\Question{
Máme dva rezistory s neznámými odpory \( R_1 \) a \( R_2 \), přičemž \( R_1 < R_2 \). Zapojíme-li tyto rezistory sériově (viz obr. A), mají celkový odpor \( R_S=64\,\Omega \). Zapojíme-li odpory paralelně (viz obr. B), mají celkový odpor \( R_P=15\,\Omega \). Určete odpor \( R_1 \).}
{\( 24 \)}{\( 22 \)}{\( 12 \)}{\( 15 \)}{}{}}
\LANGsk{
\Question{
Máme dva rezistory s neznámymi odpormi  \( R_1 \) a \( R_2 \), kde \( R_1 < R_2 \). Ak zapojíme tieto rezistory  sériovo (viz obr. A) majú celkový odpor \( R_S=64\,\Omega \). Ak zapojíme rezistory paralelne (viz obr. B), majú celkový odpor \( R_P=15\,\Omega \). Určte odpor \( R_1 \).}
{\( 24 \)}{\( 22 \)}{\( 12 \)}{\( 15 \)}{}{}}
\LANGes{
\Question{
Tenemos dos resistores de resistencias desconocidas \( R_1 \) y \( R_2 \), donde \( R_1 < R_2 \) están conectados en serie (Imagen A) y la resistencia total del circuito es \( R_S=64\,\Omega \). Si los resistores están conectados en paralelo (Imagen B), la resistencia total es \( R_P=15\,\Omega \). Calcula \( R_1 \).}
{\( 24 \)}{\( 22 \)}{\( 12 \)}{\( 15 \)}{}{}}
\End

\Data\ID{1103162606}
\Updated{2020-07-15 03:07:21}
\Level{C}
\SubArea{Systems of equations and inequalities}
\IMG{1103162606.pdf}
\LANGen{
\Question{
The figure below shows two semicircles \( \widehat{AB} \) and \( \widehat{DA} \).  Find the dilation scale factor if you know that the semicircle \( \widehat{DA} \) is the image of the semicircle \( \widehat{AB} \)  in dilation with the centre \( C \).}
{\( -0.5 \)}{\( 0.5 \)}{\( -0.25 \)}{\( 0.75 \)}{}{}}
\LANGpl{
\Question{
Rysunek poniżej przedstawia dwa półkola \( \widehat{AB} \) i \( \widehat{DA} \). Wyznacz wskaźnik dylatacji jeśli wiesz, ze półkole \( \widehat{DA} \) to obraz półkola \( \widehat{AB} \)  ze środkiem dylatacji  \( C \).}
{\( -0{,}5 \)}{\( 0{,}5 \)}{\( -0{,}25 \)}{\( 0{,}75 \)}{}{}}
\LANGcs{
\Question{
Jsou dány dvě půlkružnice  \( \widehat{AB} \) a \( \widehat{DA} \) (viz obrázek).  Určete koeficient stejnolehlosti se středem v bodě \( C \), která zobrazuje půlkružnici \( \widehat{AB} \) na půlkružnici \( \widehat{DA} \).}
{\( -0{,}5 \)}{\( 0{,}5 \)}{\( -0{,}25 \)}{\( 0{,}75 \)}{}{}}
\LANGsk{
\Question{
Na obrázku sú znázornené 2 polkružnice \( \widehat{AB} \) a \( \widehat{DA} \).  Určte koeficient rovnoľahlosti, ak viete, že polkružnica \( \widehat{DA} \) je obrazom polkružnice \( \widehat{AB} \)  v rovnoľahlosti so stredom \( C \).}
{\( -0{,}5 \)}{\( 0{,}5 \)}{\( -0{,}25 \)}{\( 0{,}75 \)}{}{}}
\LANGes{
\Question{
La siguiente imagen representa dos semicircunferencias \( \widehat{AB} \) y \( \widehat{DA} \). Halla la razón de homotecia con el centro \( C \), que representa la semicircunferencia \( \widehat{DA} \) en la semicircunferencia \( \widehat{AB} \).}
{\( -0.5 \)}{\( 0.5 \)}{\( -0.25 \)}{\( 0.75 \)}{}{}}
\End

\Data\ID{1103206301}
\Updated{2020-07-15 03:07:10}
\Level{C}
\SubArea{Systems of equations and inequalities}
\LANGen{
\Question{
Which one of the figures shows the solution of the following system of inequalities?
\[ \begin{aligned}
2x-y-8&\leq 0 \\
-x-y+2&\geq0 \\
x-1&\geq0 \\
\end{aligned} \]}
{}{}{}{}{}{}}
\LANGpl{
\Question{
Który z poniższych wykresów przedstawia rozwiązanie podanego układu nierówności?
\[ \begin{aligned}
2x-y-8&\leq 0 \\
-x-y+2&\geq0 \\
x-1&\geq0 \\
\end{aligned} \]}
{}{}{}{}{}{}}
\LANGcs{
\Question{
Vyberte obrázek, na němž je grafické řešení soustavy nerovnic.
\[ \begin{aligned}
2x-y-8&\leq 0 \\
-x-y+2&\geq0 \\
x-1&\geq0 \\
\end{aligned} \]}
{}{}{}{}{}{}}
\LANGsk{
\Question{
Vyberte obrázok, na ktorom je grafické riešenie danej sústavy nerovníc. 
\[ \begin{aligned}
2x-y-8&\leq 0 \\
-x-y+2&\geq0 \\
x-1&\geq0 \\
\end{aligned} \]}
{}{}{}{}{}{}}
\LANGes{
\Question{
NA}
{}{}{}{}{}{}}

\IMGanswer{1103206301a.pdf}
\IMGanswer{1103206301b.pdf}
\IMGanswer{1103206301c.pdf}
\IMGanswer{1103206301d.pdf}\End

\Data\ID{9000009909}
\Updated{2020-07-15 03:07:34}
\Level{C}
\SubArea{Systems of equations and inequalities}
\LANGen{
\Question{
Consider the system
\[\begin{aligned}
 y & = \frac{k} 
{x}, & &
 \\y & = a, & &
\end{aligned}\]
where \(a\),
\(k\) are real
parameters and \(x\),
\(y\) are real
variables. Find the conditions which ensure that the system has a unique solution on
\(\mathbb{R}^{-}\times \mathbb{R}^{-}\).}
{\(a < 0\) and
\(k > 0\)}{\(a < 0\) and
\(k < 0\)}{\(a > 0\) and
\(k < 0\)}{\(a > 0\) and
\(k > 0\)}{}{}}
\LANGpl{
\Question{
Rozważmy układ
\[\begin{aligned}
 y & = \frac{k} 
{x}, & &
 \\y & = a, & &
\end{aligned}\]
gdzie \(a\),
\(k\) są rzeczywistymi parametrami, a  \(x\),
\(y\) są rzeczywistymi zmiennymi. Określ warunki, dla których dany układ ma tylko jedno rozwiązanie 
\(\mathbb{R}^{-}\times \mathbb{R}^{-}\).}
{\(a < 0\) i
\(k > 0\)}{\(a < 0\) i
\(k < 0\)}{\(a > 0\) i
\(k < 0\)}{\(a > 0\) i
\(k > 0\)}{}{}}
\LANGcs{
\Question{
Je dána soustava rovnic \(y = \frac{k} 
{x}\),
\(y = a\), kde
\(a,\ k\) jsou reálné
parametry a \(x,\ y\) 
proměnné. Jaká musí být znaménka obou parametrů, aby soustava měla jediné
řešení v \(\mathbb{R}^{-}\times \mathbb{R}^{-}\)?}
{\(a < 0 \wedge k > 0\)}{\(a < 0 \wedge k < 0\)}{\(a > 0 \wedge k < 0\)}{\(a > 0 \wedge k > 0\)}{}{}}
\LANGsk{
\Question{
Je daná sústava rovníc
\[\begin{aligned}
 y & = \frac{k} 
{x}, & &
 \\y & = a, & &
\end{aligned}\]
kde \(a\),
\(k\) sú reálne
parametre a \(x\),
\(y\) sú reálne
premenné. Za akých podmienok má sústava jediné riešenie v
\(\mathbb{R}^{-}\times \mathbb{R}^{-}\)?}
{\(a < 0\) a
\(k > 0\)}{\(a < 0\) a
\(k < 0\)}{\(a > 0\) a
\(k < 0\)}{\(a > 0\) a
\(k > 0\)}{}{}}
\LANGes{
\Question{
Dado el sistema
\[\begin{aligned}
 y & = \frac{k} 
{x}, & &
 \\y & = a, & &
\end{aligned}\]
donde \(a\),
\(k\) son parámetros reales \(x\),
\(y\) son variables reales. Halla las condiciones para que el sistema tenga la única solución en
\(\mathbb{R}^{-}\times \mathbb{R}^{-}\).}
{\(a < 0\) and
\(k > 0\)}{\(a < 0\) and
\(k < 0\)}{\(a > 0\) and
\(k < 0\)}{\(a > 0\) and
\(k > 0\)}{}{}}
\End

\Data\ID{9000019901}
\Updated{2020-07-15 03:07:34}
\Level{C}
\SubArea{Systems of equations and inequalities}
\LANGen{
\Question{
Find the rank of the following matrix
\(A\).
\[ 
A = \left( \MATRIX{ 
-3& 3 & 3\cr 
 2 & -2 & 2
\cr 
-1& 1 & 1 } \right )
\]}
{\(\mathop{\mathrm{rank}}(A) = 2\)}{\(\mathop{\mathrm{rank}}(A) = 3\)}{\(\mathop{\mathrm{rank}}(A) = 1\)}{\(\mathop{\mathrm{rank}}(A) = 0\)}{}{}}
\LANGpl{
\Question{
Wyznacz rząd \( r(A) \) macierzy 
\(A\).
\[ 
A = \left( \MATRIX{ 
-3& 3 & 3\cr 
 2 & -2 & 2
\cr 
-1& 1 & 1 } \right )
\]}
{\(\mathop{\mathrm{r}}(A) = 2\)}{\(\mathop{\mathrm{r}}(A) = 3\)}{\(\mathop{\mathrm{r}}(A) = 1\)}{\(\mathop{\mathrm{r}}(A) = 0\)}{}{}}
\LANGcs{
\Question{
Určete hodnost matice \(A\).
\[ 
A = \left( \MATRIX{ 
-3& 3 & 3\cr 
 2 & -2 & 2
\cr 
-1& 1 & 1 } \right )
\]}
{\(h(A) = 2\)}{\(h(A) = 3\)}{\(h(A) = 1\)}{\(h(A) = 0\)}{}{}}
\LANGsk{
\Question{
Určte hodnosť matice
\(A\).
\[ 
A = \left( \MATRIX{ 
-3& 3 & 3\cr 
 2 & -2 & 2
\cr 
-1& 1 & 1 } \right )
\]}
{\(h(A) = 2\)}{\(h(A) = 3\)}{\(h(A) = 1\)}{\(h(A) = 0\)}{}{}}
\LANGes{
\Question{
Halla el rango de la siguiente matriz
\(A\).
\[ 
A = \left( \MATRIX{ 
-3& 3 & 3\cr 
 2 & -2 & 2
\cr 
-1& 1 & 1 } \right )
\]}
{\(\mathop{\mathrm{rango}}(A) = 2\)}{\(\mathop{\mathrm{rango}}(A) = 3\)}{\(\mathop{\mathrm{rango}}(A) = 1\)}{\(\mathop{\mathrm{rango}}(A) = 0\)}{}{}}
\End

\Data\ID{9000019902}
\Updated{2020-07-15 03:07:34}
\Level{C}
\SubArea{Systems of equations and inequalities}
\LANGen{
\Question{
Find the rank of the following matrix
\(B\).
\[ 
B = \left( \MATRIX{ 
3& 3& -2& 0\cr 
0& 1 & 2 & 1
\cr 
0& 0& -3& 0 } \right )
\]}
{\(\mathop{\mathrm{rank}}(B) = 3\)}{\(\mathop{\mathrm{rank}}(B) = 2\)}{\(\mathop{\mathrm{rank}}(B) = 1\)}{\(\mathop{\mathrm{rank}}(B) = 0\)}{}{}}
\LANGpl{
\Question{
Wyznacz rząd \( r(B) \) macierzy
\(B\).
\[ 
B = \left( \MATRIX{ 
3& 3& -2& 0\cr 
0& 1 & 2 & 1
\cr 
0& 0& -3& 0 } \right )
\]}
{\(\mathop{\mathrm{r}}(B) = 3\)}{\(\mathop{\mathrm{r}}(B) = 2\)}{\(\mathop{\mathrm{r}}(B) = 1\)}{\(\mathop{\mathrm{r}}(B) = 0\)}{}{}}
\LANGcs{
\Question{
Určete hodnost matice \(B\).
\[ 
B = \left( \MATRIX{ 
3& 3& -2& 0\cr 
0& 1 & 2 & 1
\cr 
0& 0& -3& 0 } \right )
\]}
{\(h(B) = 3\)}{\(h(B) = 2\)}{\(h(B) = 1\)}{\(h(B) = 0\)}{}{}}
\LANGsk{
\Question{
Určte hodnosť matice 
\(B\).
\[ 
B = \left( \MATRIX{ 
3& 3& -2& 0\cr 
0& 1 & 2 & 1
\cr 
0& 0& -3& 0 } \right )
\]}
{\(h(B) = 3\)}{\(h(B) = 2\)}{\(h(B) = 1\)}{\(h(B) = 0\)}{}{}}
\LANGes{
\Question{
Halla el rango de la siguiente matriz
\(B\).
\[ 
B = \left( \MATRIX{ 
3& 3& -2& 0\cr 
0& 1 & 2 & 1
\cr 
0& 0& -3& 0 } \right )
\]}
{\(\mathop{\mathrm{rango}}(B) = 3\)}{\(\mathop{\mathrm{rango}}(B) = 2\)}{\(\mathop{\mathrm{rango}}(B) = 1\)}{\(\mathop{\mathrm{rango}}(B) = 0\)}{}{}}
\End

\Data\ID{9000019903}
\Updated{2020-07-15 03:07:34}
\Level{C}
\SubArea{Systems of equations and inequalities}
\LANGen{
\Question{
Identify a true statement related to the following matrix
\(A\).
\[ 
A = \left( \MATRIX{ 
-2& 3 & 10& 5 & -5\cr 
 6 & 11 & -7 & 2 & -3
\cr 
-7& 15& -6& 2 & 4\cr 
-8 & 1 & 13 & -5 & 0
 } \right )
\]}
{\(A\) is
a \(4\times 5\) matrix
and \(a_{(3,\, 2)} = 15\).}{\(A\) is
a \(4\times 5\) matrix
and \(a_{(2,\, 3)} = 15\).}{\(A\) is
a \(5\times 4\) matrix
and \(a_{(3,\, 2)} = -7\).}{\(A\) is
a \(5\times 4\) matrix
and \(a_{(3,\, 2)} = 15\).}{}{}}
\LANGpl{
\Question{
Które z poniższych zdań odnoszących się do macierzy
\(A\) jest zgodne z prawdą?
\[ 
A = \left( \MATRIX{ 
-2& 3 & 10& 5 & -5\cr 
 6 & 11 & -7 & 2 & -3
\cr 
-7& 15& -6& 2 & 4\cr 
-8 & 1 & 13 & -5 & 0
 } \right )
\]}
{\(A\) jest 
macierzą \(4\times 5\) 
i \(a_{(3,\, 2)} = 15\).}{\(A\) jest macierzą \(4\times 5\) 
i \(a_{(2,\, 3)} = 15\).}{\(A\) jest macierzą \(5\times 4\) 
i \(a_{(3,\, 2)} = -7\).}{\(A\) jest macierzą \(5\times 4\) 
i \(a_{(3,\, 2)} = 15\).}{}{}}
\LANGcs{
\Question{
Je dána matice \(A\).
Vyberte správné tvrzení.
\[ 
A = \left( \MATRIX{ 
-2& 3 & 10& 5 & -5\cr 
 6 & 11 & -7 & 2 & -3
\cr 
-7& 15& -6& 2 & 4\cr 
-8 & 1 & 13 & -5 & 0
 } \right )
\]}
{Matice je typu \((4,\, 5)\) 
a prvek \(a_{(3,\, 2)} = 15\).}{Matice je typu \((4,\, 5)\) 
a prvek \(a_{(2,\, 3)} = 15\).}{Matice je typu \((5,\, 4)\) 
a prvek \(a_{(3,\, 2)} = -7\).}{Matice je typu \((5,\, 4)\) 
a prvek \(a_{(3,\, 2)} = 15\).}{}{}}
\LANGsk{
\Question{
Je daná matica
\(A\). Vyberte správne tvrdenie.
\[ 
A = \left( \MATRIX{ 
-2& 3 & 10& 5 & -5\cr 
 6 & 11 & -7 & 2 & -3
\cr 
-7& 15& -6& 2 & 4\cr 
-8 & 1 & 13 & -5 & 0
 } \right )
\]}
{Matica \(A\) je typu
 \(4\times 5\) 
a \(a_{(3,\, 2)} = 15\).}{Matica \(A\) je typu
 \(4\times 5\) 
a \(a_{(2,\, 3)} = 15\).}{Matica \(A\) je typu
 \(5\times 4\) 
a \(a_{(3,\, 2)} = -7\).}{Matica \(A\) je typu
 \(5\times 4\) 
a \(a_{(3,\, 2)} = 15\).}{}{}}
\LANGes{
\Question{
Identifica la proposición lógica verdadera relacionada con la siguiente matriz
\(A\).
\[ 
A = \left( \MATRIX{ 
-2& 3 & 10& 5 & -5\cr 
 6 & 11 & -7 & 2 & -3
\cr 
-7& 15& -6& 2 & 4\cr 
-8 & 1 & 13 & -5 & 0
 } \right )
\]}
{\(A\) es \(4\times 5\) matriz
y \(a_{(3,\, 2)} = 15\).}{\(A\) es \(4\times 5\) matriz
y \(a_{(2,\, 3)} = 15\).}{\(A\) es \(5\times 4\) matriz
y \(a_{(3,\, 2)} = -7\).}{\(A\) es \(5\times 4\) matriz
y \(a_{(3,\, 2)} = 15\).}{}{}}
\End

\Data\ID{9000019904}
\Updated{2020-08-23 08:08:03}
\Level{C}
\SubArea{Systems of equations and inequalities}
\LANGen{
\Question{
The coefficient matrix of a \(3\times 3\) 
linear system is \(A\) and
the augmented matrix \(A'\).
Find \(\mathop{\mathrm{rank}}(A)\) 
and \(\mathop{\mathrm{rank}}(A')\).
\[ 
 A = \begin{pmatrix}
 -1 & 3 & 2 \\
 0 & 4 & -5 \\
 0 & 0 & 2
 \end{pmatrix} \qquad
 A' = \left(\begin{array}{ccc|c}
 -1 & 3 & 2 & 5 \\
 0 & 4 & -5 & 10\\
 0 & 0 & 2 & 0
 \end{array}\right)
\]}
{\(\mathop{\mathrm{rank}}(A) = 3,\ \mathop{\mathrm{rank}}(A') = 3\)}{\(\mathop{\mathrm{rank}}(A) = 2,\ \mathop{\mathrm{rank}}(A') = 3\)}{\(\mathop{\mathrm{rank}}(A) = 3,\ \mathop{\mathrm{rank}}(A') = 2\)}{\(\mathop{\mathrm{rank}}(A) = 2,\ \mathop{\mathrm{rank}}(A') = 2\)}{}{}}
\LANGpl{
\Question{
Macierzą współczynników układu liniowego  \(3\) równań z \(3\) niewiadomymi jest macierz  \(A\) i macierzą rozszerzoną tego układu jest macierz \(A'\). Wyznacz rząd 
\(\mathop{\mathrm{r}}(A)\) macierzy \(A\) i rząd \(\mathop{\mathrm{r}}(A')\) macierzy \(A'\).
\[ 
 A = \begin{pmatrix}
 -1 & 3 & 2 \\
 0 & 4 & -5 \\
 0 & 0 & 2
 \end{pmatrix} \qquad
 A' = \left(\begin{array}{ccc|c}
 -1 & 3 & 2 & 5 \\
 0 & 4 & -5 & 10\\
 0 & 0 & 2 & 0
 \end{array}\right)
\]}
{\(\mathop{\mathrm{r}}(A) = 3,\ \mathop{\mathrm{r}}(A') = 3\)}{\(\mathop{\mathrm{r}}(A) = 2,\ \mathop{\mathrm{r}}(A') = 3\)}{\(\mathop{\mathrm{r}}(A) = 3,\ \mathop{\mathrm{r}}(A') = 2\)}{\(\mathop{\mathrm{r}}(A) = 2,\ \mathop{\mathrm{r}}(A') = 2\)}{}{}}
\LANGcs{
\Question{
Je dána soustava tří rovnic o třech neznámých, jejíž matice soustavy je
\(A\) a rozšířená matice
soustavy je \(A'\). Určete hodnost \(h(A)\)
matice soustavy \(A\) a hodnost \(h(A')\)
rozšířené matice soustavy \(A'\).
\[ 
 A = \begin{pmatrix}
 -1 & 3 & 2 \\
 0 & 4 & -5 \\
 0 & 0 & 2
 \end{pmatrix} \qquad
 A' = \left(\begin{array}{ccc|c}
 -1 & 3 & 2 & 5 \\
 0 & 4 & -5 & 10\\
 0 & 0 & 2 & 0
 \end{array}\right)
\]}
{\(h(A) = 3,\ h(A') = 3\)}{\(h(A) = 2,\ h(A') = 3\)}{\(h(A) = 3,\ h(A') = 2\)}{\(h(A) = 2,\ h(A') = 2\)}{}{}}
\LANGsk{
\Question{
Je daná sústava \(3\) rovníc o \(3\) neznámych, ktorej matica sústavy je \(A\) a rozšírená matica sústavy je \(A'\).
Určte hodnosť \(h(A)\) matice sústavy \(A\)
a hodnosť \(h(A')\) rozšírenej matice sústavy \(A'\).
\[ 
 A = \begin{pmatrix}
 -1 & 3 & 2 \\
 0 & 4 & -5 \\
 0 & 0 & 2
 \end{pmatrix} \qquad
 A' = \left(\begin{array}{ccc|c}
 -1 & 3 & 2 & 5 \\
 0 & 4 & -5 & 10\\
 0 & 0 & 2 & 0
 \end{array}\right)
\]}
{\(h(A) = 3,\ h(A') = 3\)}{\(h(A) = 2,\ h(A') = 3\)}{\(h(A) = 3,\ h(A') = 2\)}{\(h(A) = 2,\ h(A') = 2\)}{}{}}
\LANGes{
\Question{
Dado el siguiente sistema matricial.
Halla \(\mathop{\mathrm{rango}}(A)\) 
y \(\mathop{\mathrm{rango}}(A')\).
\[ 
 A = \begin{pmatrix}
 -1 & 3 & 2 \\
 0 & 4 & -5 \\
 0 & 0 & 2
 \end{pmatrix} \qquad
 A' = \left(\begin{array}{ccc|c}
 -1 & 3 & 2 & 5 \\
 0 & 4 & -5 & 10\\
 0 & 0 & 2 & 0
 \end{array}\right)
\]}
{\(\mathop{\mathrm{rango}}(A) = 3,\ \mathop{\mathrm{rango}}(A') = 3\)}{\(\mathop{\mathrm{rango}}(A) = 2,\ \mathop{\mathrm{rango}}(A') = 3\)}{\(\mathop{\mathrm{rango}}(A) = 3,\ \mathop{\mathrm{rango}}(A') = 2\)}{\(\mathop{\mathrm{rango}}(A) = 2,\ \mathop{\mathrm{rango}}(A') = 2\)}{}{}}
\End

\Data\ID{9000019905}
\Updated{2020-07-15 03:07:34}
\Level{C}
\SubArea{Systems of equations and inequalities}
\LANGen{
\Question{
Let \(A\) 
and \(A'\) 
be the coefficient matrix and the augmented matrix of the following linear system,
respectively. Find the ranks of these matrices.
\[ 
\begin{array}{cl}
\phantom{ -} 3x + 5y +\phantom{ 2}z =\phantom{ -}10& 
\\ 
 - 2x - 3y + 2z = -10& 
\\ 
\phantom{ - 2}x +\phantom{ 2}y - 5z =\phantom{ -}10& \end{array}
\]}
{\(\mathop{\mathrm{rank}}(A) = 2,\ \mathop{\mathrm{rank}}(A') = 2\)}{\(\mathop{\mathrm{rank}}(A) = 3,\ \mathop{\mathrm{rank}}(A') = 3\)}{\(\mathop{\mathrm{rank}}(A) = 3,\ \mathop{\mathrm{rank}}(A') = 2\)}{\(\mathop{\mathrm{rank}}(A) = 2,\ \mathop{\mathrm{rank}}(A') = 3\)}{}{}}
\LANGpl{
\Question{
Dany jest układ trzech równań liniowych z trzema niewiadomymi. Wyznacz rząd 
\(\mathop{\mathrm{r}}(A)\) macierzy \(A\) i rząd \(\mathop{\mathrm{r}}(A')\) macierzy \(A'\). 
\[ 
\begin{array}{cl}
\phantom{ -} 3x + 5y +\phantom{ 2}z =\phantom{ -}10& 
\\ 
 - 2x - 3y + 2z = -10& 
\\ 
\phantom{ - 2}x +\phantom{ 2}y - 5z =\phantom{ -}10& \end{array}
\]}
{\(\mathop{\mathrm{r}}(A) = 2,\ \mathop{\mathrm{r}}(A') = 2\)}{\(\mathop{\mathrm{r}}(A) = 3,\ \mathop{\mathrm{r}}(A') = 3\)}{\(\mathop{\mathrm{r}}(A) = 3,\ \mathop{\mathrm{r}}(A') = 2\)}{\(\mathop{\mathrm{r}}(A) = 2,\ \mathop{\mathrm{r}}(A') = 3\)}{}{}}
\LANGcs{
\Question{
Je dána soustava tří rovnic o třech neznámých. Určete hodnost \(h(A)\)
matice soustavy \(A\) a  hodnost \(h(A')\)
rozšířené matice soustavy \(A'\).
\[ 
\begin{array}{cl}
\phantom{ -} 3x + 5y +\phantom{ 2}z =\phantom{ -}10& 
\\ 
 - 2x - 3y + 2z = -10& 
\\ 
\phantom{ - 2}x +\phantom{ 2}y - 5z =\phantom{ -}10& \end{array}
\]}
{\(h(A) = 2,\ h(A') = 2\)}{\(h(A) = 3,\ h(A') = 3\)}{\(h(A) = 3,\ h(A') = 2\)}{\(h(A) = 2,\ h(A') = 3\)}{}{}}
\LANGsk{
\Question{
Je daná sústava troch rovníc o troch neznámych. Určte hodnosť  \(h(A)\) matice sústavy \(A\)
a hodnosť  \(h(A')\) rozšírenej matice sústavy \(A'\).
\[ 
\begin{array}{cl}
\phantom{ -} 3x + 5y +\phantom{ 2}z =\phantom{ -}10& 
\\ 
 - 2x - 3y + 2z = -10& 
\\ 
\phantom{ - 2}x +\phantom{ 2}y - 5z =\phantom{ -}10& \end{array}
\]}
{\(h(A) = 2,\  h(A') = 2\)}{\(h(A) = 3,\  h(A') = 3\)}{\(h(A) = 3,\  h(A') = 2\)}{\(h(A) = 2,\  h(A') = 3\)}{}{}}
\LANGes{
\Question{
Dado el sistema de tres ecuaciones con tres incógnitas. Halla el rango \((A)\)
de la matriz del sistema \(A\) y el rango \((A')\)
de la matriz aumentada del sistema \(A'\).
\[ 
\begin{array}{cl}
\phantom{ -} 3x + 5y +\phantom{ 2}z =\phantom{ -}10& 
\\ 
 - 2x - 3y + 2z = -10& 
\\ 
\phantom{ - 2}x +\phantom{ 2}y - 5z =\phantom{ -}10& \end{array}
\]}
{\(rango(A) = 2,\ rango(A') = 2\)}{\(rango(A) = 3,\ rango(A') = 3\)}{\(rango(A) = 3,\ rango(A') = 2\)}{\(rango(A) = 2,\ rango(A') = 3\)}{}{}}
\End

\Data\ID{9000019906}
\Updated{2020-07-15 03:07:34}
\Level{C}
\SubArea{Systems of equations and inequalities}
\LANGen{
\Question{
Consider a linear system of four equations with four unknowns. The rank of the coefficient
matrix \(A\) is
\(\mathop{\mathrm{rank}}(A) = 3\). The rank of the
augmented matrix \(A'\) 
is \(\mathop{\mathrm{rank}}(A') = 4\).
Identify a true statement on this system.}
{The system does not have any solution.}{The system has infinitely many solutions.}{The system has a unique solution.}{It is not possible to draw any conclusion from this information.}{}{}}
\LANGpl{
\Question{
Dany jest układ liniowy czterech równań z czterema niewiadomymi. Rząd macierzy 
 \(A\) jest
\(\mathop{\mathrm{r}}(A) = 3\). Rząd rozszerzonej macierzy \(A'\) 
wynosi \(\mathop{\mathrm{r}}(A') = 4\).
Które z poniższych zdań dotyczących tego układu jest prawdziwe?}
{Układ nie ma rozwiązania.}{Układ ma nieskończenie wiele rozwiązań.}{Układ ma tylko jedno rozwiązanie.}{Brak wystarczających informacji.}{}{}}
\LANGcs{
\Question{
Je dána soustava čtyř rovnic o čtyřech neznámých. Hodnost matice soustavy
\(A\) je
\(h(A) = 3\), hodnost rozšířené
matice soustavy \(A'\) 
je \(h(A') = 4\).
Kolik má daná soustava rovnic řešení?}
{žádné řešení}{nekonečně mnoho řešení}{právě jedno řešení}{nelze určit počet řešení}{}{}}
\LANGsk{
\Question{
Je daná sústava štyroch rovníc o štyroch neznámych. Hodnosť matice sústavy \(A\) je 
\(h(A) = 3\), hodnosť rozšírenej matice sústavy 
 \(A'\) 
je  \(h(A') = 4\).
Koľko riešení má daná sústava rovníc?}
{Sústava nemá žiadne riešenie.}{Sústava má nekonečne veľa riešení.}{Sústava má práve jedno riešenie.}{Nedá sa určiť počet riešení.}{}{}}
\LANGes{
\Question{
Dado el sistema de cuatro ecuaciones con cuatro incgnitas. El rango de la matriz del sistema
\(A\) es
\(rango(A) = 3\), el rango de la matriz aumentada del sistema \(A'\) 
es \(rango(A') = 4\).
¿Cuántas soluciones tiene el sistema de ecuaciones?}
{no tiene solución}{tiene infinitas soluciones}{tiene solo una solución}{no es posible hallar el número de soluciones}{}{}}
\End

\Data\ID{9000019907}
\Updated{2020-07-15 03:07:34}
\Level{C}
\SubArea{Systems of equations and inequalities}
\LANGen{
\Question{
The augmented matrix of a system of three equations with
three unknowns is row equivalent with the following matrix
\(A'\). Find
the solution of the system.
\[ 
 A' = \left(\begin{array}{ccc|c}
 1 & 2 & 4 & 0\\
 0 & 2 & 7 & 7\\
 0 & 0 & 7 & 35
 \end{array}\right)
\]}
{\([8;-14;5]\)}{\([-62;21;5]\)}{\([8;14;-5]\)}{\([-22;-21;5]\)}{}{}}
\LANGpl{
\Question{
Rozszerzona macierz układu trzech równań z trzema niewiadomymi jest równoważna wierszowo z nasz następującą macierzą
\(A'\). 
Znajdź rozwiązanie tego układu.
\[ 
 A' = \left(\begin{array}{ccc|c}
 1 & 2 & 4 & 0\\
 0 & 2 & 7 & 7\\
 0 & 0 & 7 & 35
 \end{array}\right)
\]}
{\([8;-14;5]\)}{\([-62;21;5]\)}{\([8;14;-5]\)}{\([-22;-21;5]\)}{}{}}
\LANGcs{
\Question{
Rozšířená matice soustavy tří rovnic o třech neznámých je ekvivalentní
s maticí \(A'\).
Určete správné řešení soustavy.
\[ 
 A' = \left(\begin{array}{ccc|c}
 1 & 2 & 4 & 0\\
 0 & 2 & 7 & 7\\
 0 & 0 & 7 & 35
 \end{array}\right)
\]}
{\([8;-14;5]\)}{\([-62;21;5]\)}{\([8;14;-5]\)}{\([-22;-21;5]\)}{}{}}
\LANGsk{
\Question{
Rozšírená matica sústavy troch rovníc o troch neznámych je ekvivalentná s maticou  
\(A'\). Určte správne riešenie sústavy.
\[ 
 A' = \left(\begin{array}{ccc|c}
 1 & 2 & 4 & 0\\
 0 & 2 & 7 & 7\\
 0 & 0 & 7 & 35
 \end{array}\right)
\]}
{\([8;-14;5]\)}{\([-62;21;5]\)}{\([8;14;-5]\)}{\([-22;-21;5]\)}{}{}}
\LANGes{
\Question{
La matriz aumentada de un sistema de tres ecuaciones con tres incógnitas equivale a la matriz
\(A'\). Halla la solución del sistema.
\[ 
 A' = \left(\begin{array}{ccc|c}
 1 & 2 & 4 & 0\\
 0 & 2 & 7 & 7\\
 0 & 0 & 7 & 35
 \end{array}\right)
\]}
{\([8;-14;5]\)}{\([-62;21;5]\)}{\([8;14;-5]\)}{\([-22;-21;5]\)}{}{}}
\End

\Data\ID{9000019908}
\Updated{2020-07-15 03:07:34}
\Level{C}
\SubArea{Systems of equations and inequalities}
\LANGen{
\Question{
The augmented matrix of a system of three equations with
three unknowns is row equivalent with the following matrix
\(A'\). Find
the solution of the system.
\[ 
 A' = \left(\begin{array}{ccc|c}
 -1 & 0 & 1 &-1\\
 0 & 7 & 2 & -1\\
 0 & 0 & 30 & 6
 \end{array}\right)
\]}
{\(\left [\frac{6}
{5};-\frac{1}
{5}; \frac{1} 
{5}\right ]_{}\)}{\(\left [\frac{1}
{5};-\frac{1}
{5}; \frac{6} 
{5}\right ]\)}{\(\left [\frac{1}
{5};-\frac{6}
{5};-\frac{1}
{5}\right ]\)}{\(\left [-\frac{6}
{5}; \frac{1} 
{5}; \frac{1} 
{5}\right ]\)}{}{}}
\LANGpl{
\Question{
Rozszerzona macierz układu trzech równań z trzema niewiadomymi jest równoważna wierszowo z nasz następującą macierzą
\(A'\). Znajdź rozwiązanie tego układu.
\[ 
 A' = \left(\begin{array}{ccc|c}
 -1 & 0 & 1 &-1\\
 0 & 7 & 2 & -1\\
 0 & 0 & 30 & 6
 \end{array}\right)
\]}
{\(\left [\frac{6}
{5};-\frac{1}
{5}; \frac{1} 
{5}\right ]_{}\)}{\(\left [\frac{1}
{5};-\frac{1}
{5}; \frac{6} 
{5}\right ]\)}{\(\left [\frac{1}
{5};-\frac{6}
{5};-\frac{1}
{5}\right ]\)}{\(\left [-\frac{6}
{5}; \frac{1} 
{5}; \frac{1} 
{5}\right ]\)}{}{}}
\LANGcs{
\Question{
Je dána soustava tří rovnic o třech neznámých,
jejíž rozšířená matice soustavy je ekvivalentní s maticí
\(A'\).
Vyberte správné řešení soustavy rovnic.
\[ 
 A' = \left(\begin{array}{ccc|c}
 -1 & 0 & 1 &-1\\
 0 & 7 & 2 & -1\\
 0 & 0 & 30 & 6
 \end{array}\right)
\]}
{\(\left [\frac{6}
{5};-\frac{1}
{5}; \frac{1} 
{5}\right ]_{}\)}{\(\left [\frac{1}
{5};-\frac{1}
{5}; \frac{6} 
{5}\right ]\)}{\(\left [\frac{1}
{5};-\frac{6}
{5};-\frac{1}
{5}\right ]\)}{\(\left [-\frac{6}
{5}; \frac{1} 
{5}; \frac{1} 
{5}\right ]\)}{}{}}
\LANGsk{
\Question{
Rozšírená matica sústavy troch rovníc o troch neznámych je ekvivalentná s maticou
\(A'\). Určte správne riešenie sústavy rovníc.
\[ 
 A' = \left(\begin{array}{ccc|c}
 -1 & 0 & 1 &-1\\
 0 & 7 & 2 & -1\\
 0 & 0 & 30 & 6
 \end{array}\right)
\]}
{\(\left [\frac{6}
{5};-\frac{1}
{5}; \frac{1} 
{5}\right ]_{}\)}{\(\left [\frac{1}
{5};-\frac{1}
{5}; \frac{6} 
{5}\right ]\)}{\(\left [\frac{1}
{5};-\frac{6}
{5};-\frac{1}
{5}\right ]\)}{\(\left [-\frac{6}
{5}; \frac{1} 
{5}; \frac{1} 
{5}\right ]\)}{}{}}
\LANGes{
\Question{
La matriz aumentada de un sistema de tres ecuaciones con tres incógnitas equivale a la matriz
\(A'\). Halla la solución del sistema.
\[ 
 A' = \left(\begin{array}{ccc|c}
 -1 & 0 & 1 &-1\\
 0 & 7 & 2 & -1\\
 0 & 0 & 30 & 6
 \end{array}\right)
\]}
{\(\left [\frac{6}
{5};-\frac{1}
{5}; \frac{1} 
{5}\right ]_{}\)}{\(\left [\frac{1}
{5};-\frac{1}
{5}; \frac{6} 
{5}\right ]\)}{\(\left [\frac{1}
{5};-\frac{6}
{5};-\frac{1}
{5}\right ]\)}{\(\left [-\frac{6}
{5}; \frac{1} 
{5}; \frac{1} 
{5}\right ]\)}{}{}}
\End

\Data\ID{9000019909}
\Updated{2020-07-15 03:07:34}
\Level{C}
\SubArea{Systems of equations and inequalities}
\LANGen{
\Question{
The augmented matrix of a system of three equations with three unknowns is the following matrix
\(M'\). Identify the matrix which
is row equivalent to \(M'\).
\[ 
 M' = \left(\begin{array}{ccc|c}
 1 & 2 & 4 & 14\\
 -1 & 0 & 3 & 7\\
 3 & 1 & -2 & 42
 \end{array}\right)
\]}
{\(\left(\begin{array}{ccc|c}
1 & 2 & 4 & 14\\
0 & 2 & 7 & 21\\
0 & 0 & 7 & 105
\end{array}\right)\)}{\(\left(\begin{array}{ccc|c}
1 & 2 & 4 & 14\\
0 & 2 & 7 & 21\\
0 & 0 & -8 & 70
\end{array}\right)\)}{\(\left(\begin{array}{ccc|c}
1 & 2 & 4 & 14\\
0 & 2 & 7 & 21\\
0 & 0 & -29 & -147
\end{array}\right)\)}{\(\left(\begin{array}{ccc|c}
1 & 2 & 4 & 14\\
0 & 2 & 1 & 7\\
0 & 0 & -23 & 35
\end{array}\right)\)}{}{}}
\LANGpl{
\Question{
Rozszerzoną macierzą układu trzech równań z trzema niewiadomymi jest następująca macierz
\(M'\). Która z podanych macierzy jest wierszowo równoważna z \(M'\)?
\[ 
 M' = \left(\begin{array}{ccc|c}
 1 & 2 & 4 & 14\\
 -1 & 0 & 3 & 7\\
 3 & 1 & -2 & 42
 \end{array}\right)
\]}
{\(\left(\begin{array}{ccc|c}
1 & 2 & 4 & 14\\
0 & 2 & 7 & 21\\
0 & 0 & 7 & 105
\end{array}\right)\)}{\(\left(\begin{array}{ccc|c}
1 & 2 & 4 & 14\\
0 & 2 & 7 & 21\\
0 & 0 & -8 & 70
\end{array}\right)\)}{\(\left(\begin{array}{ccc|c}
1 & 2 & 4 & 14\\
0 & 2 & 7 & 21\\
0 & 0 & -29 & -147
\end{array}\right)\)}{\(\left(\begin{array}{ccc|c}
1 & 2 & 4 & 14\\
0 & 2 & 1 & 7\\
0 & 0 & -23 & 35
\end{array}\right)\)}{}{}}
\LANGcs{
\Question{
Je dána soustava tří rovnic o třech neznámých, jejíž rozšířená matice
soustavy je \(M'\).
Vyberte matici, která je ekvivalentní s maticí
\(M'\).
\[
 M' = \left(\begin{array}{ccc|c}
 1 & 2 & 4 & 14\\
 -1 & 0 & 3 & 7\\
 3 & 1 & -2 & 42
 \end{array}\right)
\]}
{\(\left(\begin{array}{ccc|c}
1 & 2 & 4 & 14\\
0 & 2 & 7 & 21\\
0 & 0 & 7 & 105
\end{array}\right)\)}{\(\left(\begin{array}{ccc|c}
1 & 2 & 4 & 14\\
0 & 2 & 7 & 21\\
0 & 0 & -8 & 70
\end{array}\right)\)}{\(\left(\begin{array}{ccc|c}
1 & 2 & 4 & 14\\
0 & 2 & 7 & 21\\
0 & 0 & -29 & -147
\end{array}\right)\)}{\(\left(\begin{array}{ccc|c}
1 & 2 & 4 & 14\\
0 & 2 & 1 & 7\\
0 & 0 & -23 & 35
\end{array}\right)\)}{}{}}
\LANGsk{
\Question{
Je daná sústava troch rovníc o troch neznámych, ktorých rozšírená matica sústavy je 
\(M'\). Vyberte maticu, ktorá je ekvivalentná s maticou \(M'\).
\[ 
 M' = \left(\begin{array}{ccc|c}
 1 & 2 & 4 & 14\\
 -1 & 0 & 3 & 7\\
 3 & 1 & -2 & 42
 \end{array}\right)
\]}
{\(\left(\begin{array}{ccc|c}
1 & 2 & 4 & 14\\
0 & 2 & 7 & 21\\
0 & 0 & 7 & 105
\end{array}\right)\)}{\(\left(\begin{array}{ccc|c}
1 & 2 & 4 & 14\\
0 & 2 & 7 & 21\\
0 & 0 & -8 & 70
\end{array}\right)\)}{\(\left(\begin{array}{ccc|c}
1 & 2 & 4 & 14\\
0 & 2 & 7 & 21\\
0 & 0 & -29 & -147
\end{array}\right)\)}{\(\left(\begin{array}{ccc|c}
1 & 2 & 4 & 14\\
0 & 2 & 1 & 7\\
0 & 0 & -23 & 35
\end{array}\right)\)}{}{}}
\LANGes{
\Question{
La matriz aumentada de un sistema de tres ecuaciones con tres incógnitas es la siguiente matriz
\(M'\). Identifica la matriz que es equivalente a \(M'\).
\[ 
 M' = \left(\begin{array}{ccc|c}
 1 & 2 & 4 & 14\\
 -1 & 0 & 3 & 7\\
 3 & 1 & -2 & 42
 \end{array}\right)
\]}
{\(\left(\begin{array}{ccc|c}
1 & 2 & 4 & 14\\
0 & 2 & 7 & 21\\
0 & 0 & 7 & 105
\end{array}\right)\)}{\(\left(\begin{array}{ccc|c}
1 & 2 & 4 & 14\\
0 & 2 & 7 & 21\\
0 & 0 & -8 & 70
\end{array}\right)\)}{\(\left(\begin{array}{ccc|c}
1 & 2 & 4 & 14\\
0 & 2 & 7 & 21\\
0 & 0 & -29 & -147
\end{array}\right)\)}{\(\left(\begin{array}{ccc|c}
1 & 2 & 4 & 14\\
0 & 2 & 1 & 7\\
0 & 0 & -23 & 35
\end{array}\right)\)}{}{}}
\End

\Data\ID{9000019910}
\Updated{2020-07-15 03:07:34}
\Level{C}
\SubArea{Systems of equations and inequalities}
\LANGen{
\Question{
The augmented matrix of a system of three equations
with three unknowns is row equivalent to the following matrix
\(A'\). Find
the solution of the system.
\[ 
 A' = \left(\begin{array}{ccc|c}
 -1 & -6 & 1 &-20\\
 0 & 5 & 4 & -12\\
 0 & 0 & 0 & -8
 \end{array}\right)
\]}
{no solution}{\(\left [-\frac{172}
 {5} ;-\frac{12}
 {5} ;0\right ]\)}{\([-12t;4t;-8t],\ t\in \mathbb{R}\)}{\(\left [-12;4;-8\right ]\)}{}{}}
\LANGpl{
\Question{
Rozszerzona macierz układu trzech równań z trzema niewiadomymi jest wierszowo równoważna z następującą macierzą
\(A'\). Wyznacz rozwiązanie tego układu.
\[ 
 A' = \left(\begin{array}{ccc|c}
 -1 & -6 & 1 &-20\\
 0 & 5 & 4 & -12\\
 0 & 0 & 0 & -8
 \end{array}\right)
\]}
{brak rozwiązania}{\(\left [-\frac{172}
 {5} ;-\frac{12}
 {5} ;0\right ]\)}{\([-12t;4t;-8t],\ t\in \mathbb{R}\)}{\(\left [-12;4;-8\right ]\)}{}{}}
\LANGcs{
\Question{
Je dána soustava tří rovnic o třech neznámých,
jejíž rozšířená matice soustavy je ekvivalentní s maticí
\(A'\).
Vyberte správné řešení soustavy rovnic.
\[ 
 A' = \left(\begin{array}{ccc|c}
 -1 & -6 & 1 &-20\\
 0 & 5 & 4 & -12\\
 0 & 0 & 0 & -8
 \end{array}\right)
\]}
{nemá řešení}{\(\left [-\frac{172}
 {5} ;-\frac{12}
 {5} ;0\right ]\)}{\([-12t;4t;-8t],\ t\in \mathbb{R}\)}{\(\left [-12;4;-8\right ]\)}{}{}}
\LANGsk{
\Question{
Je daná sústava troch rovníc o troch neznámych, ktorej rozšírená matica sústavy je ekvivalentná s maticou 
\(A'\). Vyberte správne riešenie sústavy rovníc.
\[ 
 A' = \left(\begin{array}{ccc|c}
 -1 & -6 & 1 &-20\\
 0 & 5 & 4 & -12\\
 0 & 0 & 0 & -8
 \end{array}\right)
\]}
{nemá riešenie}{\(\left [-\frac{172}
 {5} ;-\frac{12}
 {5} ;0\right ]\)}{\([-12t;4t;-8t],\ t\in \mathbb{R}\)}{\(\left [-12;4;-8\right ]\)}{}{}}
\LANGes{
\Question{
La matriz aumentada de un sistema de tres ecuaciones con tres incógnitas equivale a la matriz
\(A'\). Halla la solución del sistema.
\[ 
 A' = \left(\begin{array}{ccc|c}
 -1 & -6 & 1 &-20\\
 0 & 5 & 4 & -12\\
 0 & 0 & 0 & -8
 \end{array}\right)
\]}
{no hay solución}{\(\left [-\frac{172}
 {5} ;-\frac{12}
 {5} ;0\right ]\)}{\([-12t;4t;-8t],\ t\in \mathbb{R}\)}{\(\left [-12;4;-8\right ]\)}{}{}}
\End

\Data\ID{9000020901}
\Updated{2020-07-15 03:07:34}
\Level{C}
\SubArea{Systems of equations and inequalities}
\IMG{000209_01-figure0.pdf}
\LANGen{
\Question{
The solution of the given set of equations can be interpreted as the intersection of the curves shown in the figure. Find the solution of the
system in \(\mathbb{R}\times \mathbb{R}\).
\[ \begin{alignedat}{80}
 &2x^{2} & - &3y &^{2} & = 2 &4 & & & & & & & &
 \\ &2x & - &3y & & = &0 & & & & & & & &
 \\\end{alignedat}\]}
{\([-6;-4],\ [6;4]\)}{\([-6;-4]\)}{\([6;4]\)}{no solution}{}{}}
\LANGpl{
\Question{
Poniższy rysunek przedstawia graficzne rozwiązanie następującego układu. Wyznacz rozwiązanie układu w 
\(\mathbb{R}\times \mathbb{R}\).
\[ \begin{alignedat}{80}
 &2x^{2} & - &3y &^{2} & = 2 &4 & & & & & & & &
 \\ &2x & - &3y & & = &0 & & & & & & & &
 \\\end{alignedat}\]}
{\([-6;-4],\ [6;4]\)}{\([-6;-4]\)}{\([6;4]\)}{nie ma rozwiązania}{}{}}
\LANGcs{
\Question{
Řešení soustavy níže uvedených rovnic lze interpretovat jako hledání
průsečíku křivek zobrazených na obrázku. Určete řešení soustavy v
\(\mathbb{R}\times \mathbb{R}\).
\[ \begin{alignedat}{80}
 &2x^{2} & - &3y &^{2} & = 2 &4 & & & & & & & &
 \\ &2x & - &3y & & = &0 & & & & & & & &
 \\\end{alignedat}\]}
{\([-6;-4],\ [6;4]\)}{\([-6;-4]\)}{\([6;4]\)}{nemá řešení}{}{}}
\LANGsk{
\Question{
Obrázok zobrazuje grafické riešenie danej sústavy rovníc. Určte riešenie sústavy v  \(\mathbb{R}\times \mathbb{R}\).
\[ \begin{alignedat}{80}
 &2x^{2} & - &3y &^{2} & = 2 &4 & & & & & & & &
 \\ &2x & - &3y & & = &0 & & & & & & & &
 \\\end{alignedat}\]}
{\([-6;-4],\ [6;4]\)}{\([-6;-4]\)}{\([6;4]\)}{nemá riešenie}{}{}}
\LANGes{
\Question{
La solución del siguiente sistema de ecuaciones se puede interpretar como el punto de intersección de las curvas representadas en la imagen. Halla la solución del sistema en \(\mathbb{R}\times \mathbb{R}\).
\[ \begin{alignedat}{80}
 &2x^{2} & - &3y &^{2} & = 2 &4 & & & & & & & &
 \\ &2x & - &3y & & = &0 & & & & & & & &
 \\\end{alignedat}\]}
{\([-6;-4],\ [6;4]\)}{\([-6;-4]\)}{\([6;4]\)}{no tiene solución}{}{}}
\End

\Data\ID{9000020902}
\Updated{2020-07-15 03:07:34}
\Level{C}
\SubArea{Systems of equations and inequalities}
\IMG{000209_02-figure0.pdf}
\LANGen{
\Question{
The solution of the given set of equations can be interpreted as the intersection of the curves shown in the figure. Find the solution of the
system in \(\mathbb{R}\times \mathbb{R}\).
\[ \begin{alignedat}{80}
 &4x^{2} & + &y &^{2} & = &20 & & & & & & & & &
 \\ &2x & + &y & & = &6 & & & & & & & & &
 \\\end{alignedat}\]}
{\([1;4],\ [2;2]\)}{\([2;2]\)}{\([1;4]\)}{no solution}{}{}}
\LANGpl{
\Question{
Poniższy rysunek przedstawia graficzne rozwiązanie następującego układu. Wyznacz rozwiązanie układu w 
 \(\mathbb{R}\times \mathbb{R}\).
\[ \begin{alignedat}{80}
 &4x^{2} & + &y &^{2} & = &20 & & & & & & & & &
 \\ &2x & + &y & & = &6 & & & & & & & & &
 \\\end{alignedat}\]}
{\([1;4],\ [2;2]\)}{\([2;2]\)}{\([1;4]\)}{nie ma rozwiązania}{}{}}
\LANGcs{
\Question{
Řešení soustavy níže uvedených rovnic lze interpretovat jako hledání
průsečíku křivek zobrazených na obrázku. Určete řešení soustavy v
\(\mathbb{R}\times \mathbb{R}\).
\[ \begin{alignedat}{80}
 &4x^{2} & + &y &^{2} & = &20 & & & & & & & & &
 \\ &2x & + &y & & = &6 & & & & & & & & &
 \\\end{alignedat}\]}
{\([1;4],\ [2;2]\)}{\([2;2]\)}{\([1;4]\)}{nemá řešení}{}{}}
\LANGsk{
\Question{
Obrázok zobrazuje grafické riešenie danej sústavy rovníc. Určte riešenie sústavy v \(\mathbb{R}\times \mathbb{R}\).
\[ \begin{alignedat}{80}
 &4x^{2} & + &y &^{2} & = &20 & & & & & & & & &
 \\ &2x & + &y & & = &6 & & & & & & & & &
 \\\end{alignedat}\]}
{\([1;4],\ [2;2]\)}{\([2;2]\)}{\([1;4]\)}{nemá riešenie}{}{}}
\LANGes{
\Question{
La solución del siguiente sistema de ecuaciones se puede interpretar como el punto de intersección de las curvas representadas en la imagen. Halla la solución del sistema en \(\mathbb{R}\times \mathbb{R}\).
\[ \begin{alignedat}{80}
 &4x^{2} & + &y &^{2} & = &20 & & & & & & & & &
 \\ &2x & + &y & & = &6 & & & & & & & & &
 \\\end{alignedat}\]}
{\([1;4],\ [2;2]\)}{\([2;2]\)}{\([1;4]\)}{no tiene solución}{}{}}
\End

\Data\ID{9000020903}
\Updated{2020-07-15 03:07:34}
\Level{C}
\SubArea{Systems of equations and inequalities}
\LANGen{
\Question{
Identify a true statement related to the solution of the following system in
\(\mathbb{R}\times \mathbb{R}\).
\[ \begin{alignedat}{80}
 &x^{2} & + &4 & &y^{2} & - & &2x & = &15 & & & & & & & & & & & &
 \\ &x & - & & &y & + & &1 & = &0 & & & & & & & & & & & &
 \\\end{alignedat}\]}
{The system has two solutions.}{The system has a unique solution.}{The system does not have any solution.}{The system has infinitely many solutions.}{}{}}
\LANGpl{
\Question{
Które z poniższych zdań odnoszących się do rozwiązania następującego układu w 
\(\mathbb{R}\times \mathbb{R}\) jest prawdziwe?
\[ \begin{alignedat}{80}
 &x^{2} & + &4 & &y^{2} & - & &2x & = &15 & & & & & & & & & & & &
 \\ &x & - & & &y & + & &1 & = &0 & & & & & & & & & & & &
 \\\end{alignedat}\]}
{Układ ma dwa rozwiązania.}{Układ ma tylko jedno rozwiązanie.}{Układ nie ma rozwiązania.}{Układ ma nieskończenie wiele rozwiązań.}{}{}}
\LANGcs{
\Question{
Rozhodněte o počtu řešení soustavy dvou rovnic v
\(\mathbb{R}\times \mathbb{R}\).
\[ \begin{alignedat}{80}
 &x^{2} & + &4 & &y^{2} & - & &2x & = &15 & & & & & & & & & & & &
 \\ &x & - & & &y & + & &1 & = &0 & & & & & & & & & & & &
 \\\end{alignedat}\]}
{dvě řešení}{jedno řešení}{žádné řešení}{nekonečně mnoho řešení}{}{}}
\LANGsk{
\Question{
Rozhodnite o počte riešení sústavy dvoch rovníc v
\(\mathbb{R}\times \mathbb{R}\).
\[ \begin{alignedat}{80}
 &x^{2} & + &4 & &y^{2} & - & &2x & = &15 & & & & & & & & & & & &
 \\ &x & - & & &y & + & &1 & = &0 & & & & & & & & & & & &
 \\\end{alignedat}\]}
{Sústava má dve riešenia.}{Sústava má jedno riešenie.}{Sústava nemá žiadne riešenie.}{Sústava má nekonečne veľa riešení.}{}{}}
\LANGes{
\Question{
Identifica la proposición lógica verdadera relacionada con la solución del siguiente sistema en
\(\mathbb{R}\times \mathbb{R}\).
\[ \begin{alignedat}{80}
 &x^{2} & + &4 & &y^{2} & - & &2x & = &15 & & & & & & & & & & & &
 \\ &x & - & & &y & + & &1 & = &0 & & & & & & & & & & & &
 \\\end{alignedat}\]}
{El sistema tiene dos soluciones.}{El sistema tiene solo una solución.}{El sistema no tiene solución.}{El sistema tiene infinitas soluciones.}{}{}}
\End

\Data\ID{9000020904}
\Updated{2020-07-15 03:07:34}
\Level{C}
\SubArea{Systems of equations and inequalities}
\LANGen{
\Question{
Find the condition on the parameter
\(c\in \mathbb{R}\) 
which ensures that the following system has two solutions in
\(\mathbb{R}\times \mathbb{R}\).
\[ \begin{alignedat}{80}
 &x^{2} & + &y^{2} & = 2 & & & & & &
 \\ &x & + &c & = y & & & & & &
 \\\end{alignedat}\]}
{\(|c| < 2\)}{\(|c| = 2\)}{\(|c| > 2\)}{\(c = 2\)}{}{}}
\LANGpl{
\Question{
Określ warunek parametru
\(c\in \mathbb{R}\), dla którego 
następujący układ ma dwa rozwiązania w 
\(\mathbb{R}\times \mathbb{R}\).
\[ \begin{alignedat}{80}
 &x^{2} & + &y^{2} & = 2 & & & & & &
 \\ &x & + &c & = y & & & & & &
 \\\end{alignedat}\]}
{\(|c| < 2\)}{\(|c| = 2\)}{\(|c| > 2\)}{\(c = 2\)}{}{}}
\LANGcs{
\Question{
Pro které \(c\in \mathbb{R}\) 
má soustava dvou rovnic dvě řešení v
\(\mathbb{R}\times \mathbb{R}\)?
\[ \begin{alignedat}{80}
 &x^{2} & + &y^{2} & = 2 & & & & & &
 \\ &x & + &c & = y & & & & & &
 \\\end{alignedat}\]}
{\(|c| < 2\)}{\(|c| = 2\)}{\(|c| > 2\)}{\(c = 2\)}{}{}}
\LANGsk{
\Question{
Pre ktoré
\(c\in \mathbb{R}\) 
má daná sústava dvoch rovníc dve riešenia v 
\(\mathbb{R}\times \mathbb{R}\)?
\[ \begin{alignedat}{80}
 &x^{2} & + &y^{2} & = 2 & & & & & &
 \\ &x & + &c & = y & & & & & &
 \\\end{alignedat}\]}
{\(|c| < 2\)}{\(|c| = 2\)}{\(|c| > 2\)}{\(c = 2\)}{}{}}
\LANGes{
\Question{
Suponiendo que
\(c\in \mathbb{R}\) 
halla la condición para que el siguiente sistema tenga dos soluciones en
\(\mathbb{R}\times \mathbb{R}\).
\[ \begin{alignedat}{80}
 &x^{2} & + &y^{2} & = 2 & & & & & &
 \\ &x & + &c & = y & & & & & &
 \\\end{alignedat}\]}
{\(|c| < 2\)}{\(|c| = 2\)}{\(|c| > 2\)}{\(c = 2\)}{}{}}
\End

\Data\ID{9000020905}
\Updated{2020-07-15 03:07:34}
\Level{C}
\SubArea{Systems of equations and inequalities}
\LANGen{
\Question{
Find the condition on the parameter
\(c\in \mathbb{R}\) 
which ensures that the following system has a unique solution in
\(\mathbb{R}\times \mathbb{R}\).
\[ \begin{alignedat}{80}
 &x^{2} & + &y^{2} & = 2 & & & & & &
 \\ &x & + &c & = y & & & & & &
 \\\end{alignedat}\]}
{\(|c| = 2\)}{\(|c| > 2\)}{\(|c| < 2\)}{\(c = 2\)}{}{}}
\LANGpl{
\Question{
Określ warunek parametru
\(c\in \mathbb{R}\), dla którego
następujący układ ma tylko jedno rozwiązanie w
\(\mathbb{R}\times \mathbb{R}\).
\[ \begin{alignedat}{80}
 &x^{2} & + &y^{2} & = 2 & & & & & &
 \\ &x & + &c & = y & & & & & &
 \\\end{alignedat}\]}
{\(|c| = 2\)}{\(|c| > 2\)}{\(|c| < 2\)}{\(c = 2\)}{}{}}
\LANGcs{
\Question{
Pro které \(c\in \mathbb{R}\) 
má soustava dvou rovnic jedno řešení v
\(\mathbb{R}\times \mathbb{R}\)?
\[ \begin{alignedat}{80}
 &x^{2} & + &y^{2} & = 2 & & & & & &
 \\ &x & + &c & = y & & & & & &
 \\\end{alignedat}\]}
{\(|c| = 2\)}{\(|c| > 2\)}{\(|c| < 2\)}{\(c = 2\)}{}{}}
\LANGsk{
\Question{
Pre ktoré 
\(c\in \mathbb{R}\) 
má sústava dvoch rovníc jedno riešenie v 
\(\mathbb{R}\times \mathbb{R}\)?
\[ \begin{alignedat}{80}
 &x^{2} & + &y^{2} & = 2 & & & & & &
 \\ &x & + &c & = y & & & & & &
 \\\end{alignedat}\]}
{\(|c| = 2\)}{\(|c| > 2\)}{\(|c| < 2\)}{\(c = 2\)}{}{}}
\LANGes{
\Question{
Suponiendo que
\(c\in \mathbb{R}\) 
halla la condición para que el siguiente sistema tenga dos soluciones en
\(\mathbb{R}\times \mathbb{R}\).
\[ \begin{alignedat}{80}
 &x^{2} & + &y^{2} & = 2 & & & & & &
 \\ &x & + &c & = y & & & & & &
 \\\end{alignedat}\]}
{\(|c| = 2\)}{\(|c| > 2\)}{\(|c| < 2\)}{\(c = 2\)}{}{}}
\End

\Data\ID{9000020906}
\Updated{2020-08-23 08:08:12}
\Level{C}
\SubArea{Systems of equations and inequalities}
\LANGen{
\Question{
Identify an equation which can be obtained from the following system by
eliminating one of the variables.
\[ \begin{alignedat}{80}
 &y^{2} & - &2 &x & + &3 & = 0 & & & & & & & &
 \\ &x & - & &y & - &1 & = 0 & & & & & & & &
 \\\end{alignedat}\]}
{\((y - 1)^{2} = 0\)}{\((y + 1)^{2} = 0\)}{\((x - 4)^{2} = 0\)}{\((x + 2)^{2} = 0\)}{}{}}
\LANGpl{
\Question{
Wybierz równanie, które otrzymamy po usunięciu jednej ze zmiennych z następującego układu.
\[ \begin{alignedat}{80}
 &y^{2} & - &2 &x & + &3 & = 0 & & & & & & & &
 \\ &x & - & &y & - &1 & = 0 & & & & & & & &
 \\\end{alignedat}\]}
{\((y - 1)^{2} = 0\)}{\((y + 1)^{2} = 0\)}{\((x - 4)^{2} = 0\)}{\((x + 2)^{2} = 0\)}{}{}}
\LANGcs{
\Question{
Vyberte rovnici o jedné neznámé, na kterou vede soustava dvou rovnic o
dvou neznámých.
\[ \begin{alignedat}{80}
 &y^{2} & - &2 &x & + &3 & = 0 & & & & & & & &
 \\ &x & - & &y & - &1 & = 0 & & & & & & & &
 \\\end{alignedat}\]}
{\((y - 1)^{2} = 0\)}{\((y + 1)^{2} = 0\)}{\((x - 4)^{2} = 0\)}{\((x + 2)^{2} = 0\)}{}{}}
\LANGsk{
\Question{
Vyberte rovnicu o jednej neznámej, ktorú možno získať z danej sústavy dvoch rovníc o dvoch neznámych.
\[ \begin{alignedat}{80}
 &y^{2} & - &2 &x & + &3 & = 0 & & & & & & & &
 \\ &x & - & &y & - &1 & = 0 & & & & & & & &
 \\\end{alignedat}\]}
{\((y - 1)^{2} = 0\)}{\((y + 1)^{2} = 0\)}{\((x - 4)^{2} = 0\)}{\((x + 2)^{2} = 0\)}{}{}}
\LANGes{
\Question{
Identifica la ecuación que se puede obtener del siguiente sistema eliminando una de las variables.
\[ \begin{alignedat}{80}
 &y^{2} & - &2 &x & + &3 & = 0 & & & & & & & &
 \\ &x & - & &y & - &1 & = 0 & & & & & & & &
 \\\end{alignedat}\]}
{\((y - 1)^{2} = 0\)}{\((y + 1)^{2} = 0\)}{\((x - 4)^{2} = 0\)}{\((x + 2)^{2} = 0\)}{}{}}
\End

\Data\ID{9000020907}
\Updated{2020-07-15 03:07:34}
\Level{C}
\SubArea{Systems of equations and inequalities}
\LANGen{
\Question{
Identify a true statement related to the solution of the following system in
\(\mathbb{R}\times \mathbb{R}\).
\[ \begin{alignedat}{80}
 &2x^{2} & - & &y^{2} & - &2 &x & - 5 & = 0 & & & & & & & & & &
 \\ & & & &3x & - & &y & - 5 & = 0 & & & & & & & & & &
 \\\end{alignedat}\]}
{The system has no solution.}{The system has two solutions.}{The system has a unique solution.}{None of the above conclusions can be obtained.}{}{}}
\LANGpl{
\Question{
Które z poniższych zdań odnoszących się do rozwiązania następującego układu  w 
\(\mathbb{R}\times \mathbb{R}\) jest prawdziwe?
\[ \begin{alignedat}{80}
 &2x^{2} & - & &y^{2} & - &2 &x & - 5 & = 0 & & & & & & & & & &
 \\ & & & &3x & - & &y & - 5 & = 0 & & & & & & & & & &
 \\\end{alignedat}\]}
{Układ nie ma rozwiązania.}{Układ ma dwa rozwiązania.}{Układ ma tylko jedno rozwiązanie.}{Żaden z powyższych wniosków nie jest zgodny z prawdą.}{}{}}
\LANGcs{
\Question{
Rozhodněte o počtu řešení soustavy dvou rovnic v
\(\mathbb{R}\times \mathbb{R}\).
\[ \begin{alignedat}{80}
 &2x^{2} & - & &y^{2} & - &2 &x & - 5 & = 0 & & & & & & & & & &
 \\ & & & &3x & - & &y & - 5 & = 0 & & & & & & & & & &
 \\\end{alignedat}\]}
{žádné řešení}{dvě řešení}{jedno řešení}{nelze rozhodnout}{}{}}
\LANGsk{
\Question{
Rozhodnite o počte riešení sústavy dvoch rovníc v
\(\mathbb{R}\times \mathbb{R}\).
\[ \begin{alignedat}{80}
 &2x^{2} & - & &y^{2} & - &2 &x & - 5 & = 0 & & & & & & & & & &
 \\ & & & &3x & - & &y & - 5 & = 0 & & & & & & & & & &
 \\\end{alignedat}\]}
{Sústava nemá žiadne riešenie.}{Sústava má dve riešenia.}{Sústava má práve jedno riešenie.}{Nedá sa rozhodnúť.}{}{}}
\LANGes{
\Question{
Identifica la proposición lógica verdadera relacionada con la solución del siguiente sistema en
\(\mathbb{R}\times \mathbb{R}\).
\[ \begin{alignedat}{80}
 &2x^{2} & - & &y^{2} & - &2 &x & - 5 & = 0 & & & & & & & & & &
 \\ & & & &3x & - & &y & - 5 & = 0 & & & & & & & & & &
 \\\end{alignedat}\]}
{El sistema no tiene solución.}{El sistema tiene dos soluciones.}{El sistema tiene solo una solución.}{No es posible identificarla.}{}{}}
\End

\Data\ID{9000020908}
\Updated{2020-07-15 03:07:34}
\Level{C}
\SubArea{Systems of equations and inequalities}
\LANGen{
\Question{
Assuming that the real parameter \(c\) 
satisfies \(c > 16\),
solve the system and identify a true statement.
\[ \begin{alignedat}{80}
 &y^{2} & - &4x & & = 0 & & & & & &
 \\8 &x & - &4y & + c & = 0 & & & & & &
 \\\end{alignedat}\]}
{The system has no solution.}{The system has two solutions.}{The system has a unique solution.}{The system has infinitely many solutions.}{}{}}
\LANGpl{
\Question{
Zakładając, że rzeczywisty parametr \(c\) 
spełnia \(c > 16\),
rozwiąż podany układ i wybierz stwierdzenie zgodne z prawdą.
\[ \begin{alignedat}{80}
 &y^{2} & - &4x & & = 0 & & & & & &
 \\8 &x & - &4y & + c & = 0 & & & & & &
 \\\end{alignedat}\]}
{Układ nie ma rozwiązania.}{Układ ma dwa rozwiązania.}{Układ ma tylko jedno rozwiązanie.}{Układ ma nieskończenie wiele rozwiązań.}{}{}}
\LANGcs{
\Question{
Rozhodněte o počtu řešení soustavy dvou rovnic v
\(\mathbb{R}\times \mathbb{R}\) pro
parametr \(c > 16\).
\[ \begin{alignedat}{80}
 &y^{2} & - &4x & & = 0 & & & & & &
 \\8 &x & - &4y & + c & = 0 & & & & & &
 \\\end{alignedat}\]}
{žádné řešení}{dvě řešení}{jedno řešení}{nekonečně mnoho řešení}{}{}}
\LANGsk{
\Question{
Rozhodnite o počte riešení sústavy dvoch rovníc pre reálny parameter \(c\),  
 \(c > 16\).
\[ \begin{alignedat}{80}
 &y^{2} & - &4x & & = 0 & & & & & &
 \\8 &x & - &4y & + c & = 0 & & & & & &
 \\\end{alignedat}\]}
{Sústava nem žiadne riešenie.}{Sústava má dve riešenia.}{Sústava má práve jedno riešenie.}{Sústava má nekonečne veľa riešení.}{}{}}
\LANGes{
\Question{
Suponiendo que el parámetro real \(c\) 
satisface \(c > 16\),
resuelve el sistema e identifica la proposición lógica verdadera.
\[ \begin{alignedat}{80}
 &y^{2} & - &4x & & = 0 & & & & & &
 \\8 &x & - &4y & + c & = 0 & & & & & &
 \\\end{alignedat}\]}
{El sistema no tiene solución.}{El sistema tiene dos soluciones.}{El sistema tiene solo una solución.}{El sistema tiene infinitas soluciones.}{}{}}
\End

\Data\ID{9000020910}
\Updated{2020-07-15 03:07:34}
\Level{C}
\SubArea{Systems of equations and inequalities}
\LANGen{
\Question{
The perimeter of a rectangle is \(28\, \mathrm{cm}\).
The diagonal of this rectangle is \(10\, \mathrm{cm}\).
Find the sides of the rectangle.}
{\(8\, \mathrm{cm}\) 
and \(6\, \mathrm{cm}\)}{\(7\, \mathrm{cm}\) 
and \(7\, \mathrm{cm}\)}{\(9\, \mathrm{cm}\) 
and \(5\, \mathrm{cm}\)}{\(7\, \mathrm{cm}\) 
and \(3\, \mathrm{cm}\)}{}{}}
\LANGpl{
\Question{
Obwód prostokąta wynosi  \(28\, \mathrm{cm}\). Przekątna tego prostokąta jest równa  \(10\, \mathrm{cm}\).
Znajdź długości jego boków.}
{\(8\, \mathrm{cm}\) 
i \(6\, \mathrm{cm}\)}{\(7\, \mathrm{cm}\) 
i \(7\, \mathrm{cm}\)}{\(9\, \mathrm{cm}\) 
i \(5\, \mathrm{cm}\)}{\(7\, \mathrm{cm}\) 
i \(3\, \mathrm{cm}\)}{}{}}
\LANGcs{
\Question{
Obdélník má obvod \(28\, \mathrm{cm}\) 
a úhlopříčku \(10\, \mathrm{cm}\).
Určete rozměry obdélníku.}
{\(8\, \mathrm{cm}\) a
\(6\, \mathrm{cm}\)}{\(7\, \mathrm{cm}\) a
\(7\, \mathrm{cm}\)}{\(9\, \mathrm{cm}\) a
\(5\, \mathrm{cm}\)}{\(7\, \mathrm{cm}\) a
\(3\, \mathrm{cm}\)}{}{}}
\LANGsk{
\Question{
Obvod obdĺžnika je \(28\, \mathrm{cm}\).
Uhlopriečka tohto obdĺžnika je  \(10\, \mathrm{cm}\).
Určte rozmery obdĺžnika.}
{\(8\, \mathrm{cm}\) 
a \(6\, \mathrm{cm}\)}{\(7\, \mathrm{cm}\) 
a \(7\, \mathrm{cm}\)}{\(9\, \mathrm{cm}\) 
a \(5\, \mathrm{cm}\)}{\(7\, \mathrm{cm}\) 
a \(3\, \mathrm{cm}\)}{}{}}
\LANGes{
\Question{
El perímetro de un rectángulo mide \(28\, \mathrm{cm}\).
La diagonal de este rectángulo mide \(10\, \mathrm{cm}\).
¿Cuáles son las medidas del rectángulo?}
{\(8\, \mathrm{cm}\) 
y \(6\, \mathrm{cm}\)}{\(7\, \mathrm{cm}\) 
y \(7\, \mathrm{cm}\)}{\(9\, \mathrm{cm}\) 
y \(5\, \mathrm{cm}\)}{\(7\, \mathrm{cm}\) 
y \(3\, \mathrm{cm}\)}{}{}}
\End

\Data\ID{9000022901}
\Updated{2020-08-23 08:08:41}
\Level{C}
\SubArea{Systems of equations and inequalities}
\LANGen{
\Question{
An arrow has been shot at the angle
\(60^{\circ }\) at the
speed \(10\, \mathrm{m}\, \mathrm{s}^{-1}\).
Find the time when the height equals to the horizontal distance from the take-off
point. Hint: The position is given by the equations
\(x = v_{0}t\cdot \cos \alpha \),
\(y = v_{0}t\cdot \sin \alpha -\frac{1} 
{2}gt^{2}\). Use
\(g = 10\, \mathrm{m}\, \mathrm{s}^{-2}\) as an
acceleration of gravity.}
{\(\left (\sqrt{3} - 1\right )\, \mathrm{s}\)}{\(\left (\sqrt{3} + 1\right )\, \mathrm{s}\)}{\(\sqrt{3}\, \mathrm{s}\)}{\(\left (\sqrt{2} - 1\right )\, \mathrm{s}\)}{\(\left (\sqrt{2} + 1\right )\, \mathrm{s}\)}{}}
\LANGpl{
\Question{
Strzała została wystrzelona pod kątem
\(60^{\circ }\) z prędkością
 \(10\, \mathrm{m}\, \mathrm{s}^{-1}\).
Określ czas kiedy wysokość będzie równa horyzontalnej odległości od punktu wystrzału.  
 Wskazówka: Pozycję określają równania
\(x = v_{0}t\cdot \cos \alpha \),
\(y = v_{0}t\cdot \sin \alpha -\frac{1} 
{2}gt^{2}\). Użyj 
\(g = 10\, \mathrm{m}\, \mathrm{s}^{-2}\) jako przyspieszenia ziemskiego.}
{\(\left (\sqrt{3} - 1\right )\, \mathrm{s}\)}{\(\left (\sqrt{3} + 1\right )\, \mathrm{s}\)}{\(\sqrt{3}\, \mathrm{s}\)}{\(\left (\sqrt{2} - 1\right )\, \mathrm{s}\)}{\(\left (\sqrt{2} + 1\right )\, \mathrm{s}\)}{}}
\LANGcs{
\Question{
Jak dlouho bude trvat, než šíp vystřelený rychlostí
\(10\, \mathrm{m}\, \mathrm{s}^{-1}\) pod
úhlem \(60^{\circ }\),
bude stejně vysoko, jako daleko (ve vodorovném směru)? Nápověda: Poloha vrženého tělesa v daném okamžiku je popsána rovnicemi
\(x = v_{0}t\cdot \cos \alpha \),
\(y = v_{0}t\cdot \sin \alpha -\frac{1} 
{2}gt^{2}\).
Za tíhové zrychlení dosazujte zaokrouhlenou hodnotu
\(g = 10\, \mathrm{m}\, \mathrm{s}^{-2}\).}
{\(\left (\sqrt{3} - 1\right )\, \mathrm{s}\)}{\(\left (\sqrt{3} + 1\right )\, \mathrm{s}\)}{\(\sqrt{3}\, \mathrm{s}\)}{\(\left (\sqrt{2} - 1\right )\, \mathrm{s}\)}{\(\left (\sqrt{2} + 1\right )\, \mathrm{s}\)}{}}
\LANGsk{
\Question{
Ako  dlho bude trvať, než šíp vystrelený pod uhlom
\(60^{\circ }\) rýchlosťou
 \(10\, \mathrm{m}\, \mathrm{s}^{-1}\) bude rovnako vysoko, ako ďaleko (vo vodorovnom smere)?  Pomôcka: Poloha vrhnutého telesa v danom okamžiku je popísaná rovnicami
\(x = v_{0}t\cdot \cos \alpha \),
\(y = v_{0}t\cdot \sin \alpha -\frac{1} 
{2}gt^{2}\). Použite
\(g = 10\, \mathrm{m}\, \mathrm{s}^{-2}\) pre
gravitačné zrýchlenie.}
{\(\left (\sqrt{3} - 1\right )\, \mathrm{s}\)}{\(\left (\sqrt{3} + 1\right )\, \mathrm{s}\)}{\(\sqrt{3}\, \mathrm{s}\)}{\(\left (\sqrt{2} - 1\right )\, \mathrm{s}\)}{\(\left (\sqrt{2} + 1\right )\, \mathrm{s}\)}{}}
\LANGes{
\Question{
Se ha disparado una flecha con un ángulo de
\(60^{\circ }\) respecto a la horizontal y con una velocidad de \(10\, \mathrm{m}\, \mathrm{s}^{-1}\).
Halla el momento en el cual la altura sea igual que la distancia horizontal desde el punto de disparo. Sugerencia: La posición viene dada por la ecuación
\(x = v_{0}t\cdot \cos \alpha \),
\(y = v_{0}t\cdot \sin \alpha -\frac{1} 
{2}gt^{2}\). Utiliza
\(g = 10\, \mathrm{m}\, \mathrm{s}^{-2}\) como la aceleración de la gravedad.}
{\(\left (\sqrt{3} - 1\right )\, \mathrm{s}\)}{\(\left (\sqrt{3} + 1\right )\, \mathrm{s}\)}{\(\sqrt{3}\, \mathrm{s}\)}{\(\left (\sqrt{2} - 1\right )\, \mathrm{s}\)}{\(\left (\sqrt{2} + 1\right )\, \mathrm{s}\)}{}}
\End

\Data\ID{9000022904}
\Updated{2020-07-15 03:07:34}
\Level{C}
\SubArea{Systems of equations and inequalities}
\LANGen{
\Question{
Find the values of a real parameter
\(t\) which
ensure that the following system has a unique solution.
\[ \begin{alignedat}{80}
 2x & + &y & + &t & = - &2 & & & & & & & &
 \\ - 4x & - 2 &y & + &1 & = &0 & & & & & & & &
 \\\end{alignedat}\]}
{\(t\in \emptyset \)}{\(t\in \mathbb{R}\)}{\(t = 3\)}{\(t = 1\)}{\(t\in \mathbb{R}\setminus \{3\}\)}{}}
\LANGpl{
\Question{
Wyznacz wartości rzeczywistego parametru
\(t\), dla którego 
następujący układ ma tylko jedno rozwiązanie.
\[ \begin{alignedat}{80}
 2x & + &y & + &t & = - &2 & & & & & & & &
 \\ - 4x & - 2 &y & + &1 & = &0 & & & & & & & &
 \\\end{alignedat}\]}
{\(t\in \emptyset \)}{\(t\in \mathbb{R}\)}{\(t = 3\)}{\(t = 1\)}{\(t\in \mathbb{R}\setminus \{3\}\)}{}}
\LANGcs{
\Question{
Pro které hodnoty reálného parametru
\(t\) bude
mít níže uvedená soustava rovnic právě jedno řešení?
\[ \begin{alignedat}{80}
 2x & + &y & + &t & = - &2 & & & & & & & &
 \\ - 4x & - 2 &y & + &1 & = &0 & & & & & & & &
 \\\end{alignedat}\]}
{\(t\in\emptyset\)}{\(\mathop{\forall }t\in \mathbb{R}\)}{\(t = 3\)}{\(t = 1\)}{\(\mathop{\forall }t\in \mathbb{R}\setminus \{3\}\)}{}}
\LANGsk{
\Question{
Pre ktoré hodnoty reálneho parametra
\(t\) bude mať uvedená sústava rovníc práve jedno riešenie?
\[ \begin{alignedat}{80}
 2x & + &y & + &t & = - &2 & & & & & & & &
 \\ - 4x & - 2 &y & + &1 & = &0 & & & & & & & &
 \\\end{alignedat}\]}
{\(t\in \emptyset \)}{\(t\in \mathbb{R}\)}{\(t = 3\)}{\(t = 1\)}{\(t\in \mathbb{R}\setminus \{3\}\)}{}}
\LANGes{
\Question{
Encuentra el valor del parámetro real \(t\) para que el siguiente sistema tenga solo una solución.
\[ \begin{alignedat}{80}
 2x & + &y & + &t & = - &2 & & & & & & & &
 \\ - 4x & - 2 &y & + &1 & = &0 & & & & & & & &
 \\\end{alignedat}\]}
{\(t\in \emptyset \)}{\(t\in \mathbb{R}\)}{\(t = 3\)}{\(t = 1\)}{\(t\in \mathbb{R}\setminus \{3\}\)}{}}
\End

\Data\ID{9000022905}
\Updated{2020-07-15 03:07:34}
\Level{C}
\SubArea{Systems of equations and inequalities}
\LANGen{
\Question{
Find the values of a real parameter
\(t\) which
ensure that the following system has a unique solution.
\[ \begin{alignedat}{80}
 tx & + &y & + &3 & = 0 & & & & & &
 \\4x & - 2 &y & + &1 & = 0 & & & & & &
 \\\end{alignedat}\]}
{\(t\in \mathbb{R}\setminus \{ - 2\}\)}{\(t\in \mathbb{R}\)}{\(t = -2\)}{\(t\in \emptyset \)}{}{}}
\LANGpl{
\Question{
Wyznacz wartości rzeczywistego parametru
\(t\), dla którego 
podany układ ma tylko jedno rozwiązanie.
\[ \begin{alignedat}{80}
 tx & + &y & + &3 & = 0 & & & & & &
 \\4x & - 2 &y & + &1 & = 0 & & & & & &
 \\\end{alignedat}\]}
{\(t\in \mathbb{R}\setminus \{ - 2\}\)}{\(t\in \mathbb{R}\)}{\(t = -2\)}{\(t\in \emptyset \)}{}{}}
\LANGcs{
\Question{
Pro které hodnoty reálného parametru
\(t\) bude
mít níže uvedená soustava rovnic právě jedno řešení?
\[ \begin{alignedat}{80}
 tx & + &y & + &3 & = 0 & & & & & &
 \\4x & - 2 &y & + &1 & = 0 & & & & & &
 \\\end{alignedat}\]}
{\(\mathop{\forall }t\in \mathbb{R}\setminus \{ - 2\}\)}{\(\mathop{\forall }t\in \mathbb{R}\)}{\(t = -2\)}{\(t\in \emptyset\)}{}{}}
\LANGsk{
\Question{
Pre ktoré hodnoty reálneho parametra
\(t\) bude mať uvedená sústava rovníc práve jedno riešenie?
\[ \begin{alignedat}{80}
 tx & + &y & + &3 & = 0 & & & & & &
 \\4x & - 2 &y & + &1 & = 0 & & & & & &
 \\\end{alignedat}\]}
{\(t\in \mathbb{R}\setminus \{ - 2\}\)}{\(t\in \mathbb{R}\)}{\(t = -2\)}{\(t\in \emptyset \)}{}{}}
\LANGes{
\Question{
Encuentra el valor del parámetro real \(t\) para que el siguiente sistema tenga solo una solución.
\[ \begin{alignedat}{80}
 tx & + &y & + &3 & = 0 & & & & & &
 \\4x & - 2 &y & + &1 & = 0 & & & & & &
 \\\end{alignedat}\]}
{\(t\in \mathbb{R}\setminus \{ - 2\}\)}{\(t\in \mathbb{R}\)}{\(t = -2\)}{\(t\in \emptyset \)}{}{}}
\End

\Data\ID{9000022906}
\Updated{2020-07-15 03:07:34}
\Level{C}
\SubArea{Systems of equations and inequalities}
\LANGen{
\Question{
Find the values of a real parameter
\(t\) 
which ensure that the following system has a unique solution
\([a,b]\) such that
both \(a\) 
and \(b\) 
are positive real numbers.
\[ \begin{alignedat}{80}
 a & - &tb & = - &2 & & & & & &
 \\a & + 2 &tb & = &0 & & & & & &
 \\\end{alignedat}\]}
{\(t\in \emptyset \)}{\(t\in \mathbb{R}^{+}\)}{\(t\in \mathbb{R}^{-}\)}{\(t = 0\)}{\(t\in \mathbb{R}\)}{}}
\LANGpl{
\Question{
Wyznacz wartości rzeczywistego parametru
\(t\), dla którego
podany układ ma tylko jedno rozwiązanie
\([a,b]\) takie, że 
zarówno \(a\) 
i \(b\) 
liczbami rzeczywistymi dodatnimi.
\[ \begin{alignedat}{80}
 a & - &tb & = - &2 & & & & & &
 \\a & + 2 &tb & = &0 & & & & & &
 \\\end{alignedat}\]}
{\(t\in \emptyset \)}{\(t\in \mathbb{R}^{+}\)}{\(t\in \mathbb{R}^{-}\)}{\(t = 0\)}{\(t\in \mathbb{R}\)}{}}
\LANGcs{
\Question{
Pro které hodnoty reálného parametru
\(t\) 
bude mít níže uvedená soustava právě jedno řešení
\([a,b]\) takové,
že \(a\) i
\(b\) 
budou kladná čísla?
\[ \begin{alignedat}{80}
 a & - &tb & = - &2 & & & & & &
 \\a & + 2 &tb & = &0 & & & & & &
 \\\end{alignedat}\]}
{\(t\in \emptyset\)}{\(\mathop{\forall }t\in \mathbb{R}^{+}\)}{\(\mathop{\forall }t\in \mathbb{R}^{-}\)}{\(t = 0\)}{\(\mathop{\forall }t\in \mathbb{R}\)}{}}
\LANGsk{
\Question{
Pre ktoré hodnoty reálneho parametra
\(t\) bude mať uvedená sústava práve jedno riešenie 
\([a,b]\) také, že 
 \(a\) 
aj \(b\) 
sú kladné reálne čísla?
\[ \begin{alignedat}{80}
 a & - &tb & = - &2 & & & & & &
 \\a & + 2 &tb & = &0 & & & & & &
 \\\end{alignedat}\]}
{\(t\in \emptyset \)}{\(t\in \mathbb{R}^{+}\)}{\(t\in \mathbb{R}^{-}\)}{\(t = 0\)}{\(t\in \mathbb{R}\)}{}}
\LANGes{
\Question{
Encuentra el valor del parámetro real \(t\)  para que el siguiente sistema tenga solo una solución \([a,b]\) suponiendo que \(a\) y \(b\) son números positivos.
\[ \begin{alignedat}{80}
 a & - &tb & = - &2 & & & & & &
 \\a & + 2 &tb & = &0 & & & & & &
 \\\end{alignedat}\]}
{\(t\in \emptyset \)}{\(t\in \mathbb{R}^{+}\)}{\(t\in \mathbb{R}^{-}\)}{\(t = 0\)}{\(t\in \mathbb{R}\)}{}}
\End

\Data\ID{9000022907}
\Updated{2020-07-15 03:07:34}
\Level{C}
\SubArea{Systems of equations and inequalities}
\LANGen{
\Question{
Solve the following system of equations and identify a true statement.
\[ \begin{alignedat}{80}
 |x - 2| & + &y & = &2 & & & & & &
 \\ - 2|5 + x| &- 3 &y & = - &5 & & & & & &
 \\\end{alignedat}\]}
{The system has two solutions. Both solutions satisfy
\(y < 0\).}{The system has a unique solution. This solution satisfies
\(y > 0\).}{The system has two solutions. Both solutions satisfy
\(y > 0\).}{The system has more than two solutions.}{The system does not have any solution.}{}}
\LANGpl{
\Question{
Rozwiąż podany układ i wybierz zdanie zgodne z prawdą. 
\[ \begin{alignedat}{80}
 |x - 2| & + &y & = &2 & & & & & &
 \\ - 2|5 + x| &- 3 &y & = - &5 & & & & & &
 \\\end{alignedat}\]}
{Układ ma dwa rozwiązania. Obydwa rozwiązania spełniają warunek
\(y < 0\).}{Układ ma tylko jedno rozwiązanie, które spełnia warunek
\(y > 0\).}{Układ ma dwa rozwiązania. Obydwa rozwiązania spełniają warunek
\(y > 0\).}{Układ ma więcej niż dwa rozwiązania.}{Układ nie ma rozwiązania.}{}}
\LANGcs{
\Question{
Vyřešte níže uvedenou soustavu rovnic a vyberte správná tvrzení o
řešení soustavy.
\[ \begin{alignedat}{80}
 |x - 2| & + &y & = &2 & & & & & &
 \\ - 2|5 + x| &- 3 &y & = - &5 & & & & & &
 \\\end{alignedat}\]}
{Soustava má právě dvě řešení. Pro obě řešení platí, že
\(y < 0\).}{Soustava má právě jedno řešení, pro které platí, že
\(y > 0\).}{Soustava má právě dvě řešení. Pro obě řešení platí, že
\(y > 0\).}{Soustava má více než dvě řešení.}{Soustava nemá řešení.}{}}
\LANGsk{
\Question{
Vyriešte danú sústavu rovníc a vyberte správne tvrdenie o riešení sústavy.
\[ \begin{alignedat}{80}
 |x - 2| & + &y & = &2 & & & & & &
 \\ - 2|5 + x| &- 3 &y & = - &5 & & & & & &
 \\\end{alignedat}\]}
{Sústava má práve dve riešenia. Pre obe riešenia platí, že
\(y < 0\).}{Sústava má práve jedno riešenie, pre ktoré platí, že 
\(y > 0\).}{Sústava má práve dve riešenia. Pre obe riešenia platí, že 
\(y > 0\).}{Sústava má viac než dve riešenia.}{Sústava nemá riešenie.}{}}
\LANGes{
\Question{
Resuelve el siguiente sistema de ecuaciones e identifica la proposición lógica.
\[ \begin{alignedat}{80}
 |x - 2| & + &y & = &2 & & & & & &
 \\ - 2|5 + x| &- 3 &y & = - &5 & & & & & &
 \\\end{alignedat}\]}
{El sistema tiene dos soluciones. Para ambas soluciones vale
\(y < 0\).}{El sistema tiene solo una solución que satisface
\(y > 0\).}{El sistema tiene dos soluciones. Para ambas soluciones vale
\(y > 0\).}{El sistema tiene más de dos soluciones.}{El sistema no tiene ninguna solución.}{}}
\End

\Data\ID{9000031101}
\Updated{2020-08-23 08:08:59}
\Level{C}
\SubArea{Systems of equations and inequalities}
\LANGen{
\Question{
Solve the following system of equations and identify a correct statement.
\[\begin{aligned}
 (x - 3)^{2} + (y - 1)^{2} = 1 & &
 \\2x^{2} + 2y^{2} - 12x - 4y + 18 = 0 & &
\end{aligned}\]}
{The system has more than two solutions.}{The system does not have any solution.}{The system has a unique solution.}{The system has two solutions.}{}{}}
\LANGpl{
\Question{
Rozwiąż następujący układ równań i wybierz poprawną odpowiedź. 
\[\begin{aligned}
 (x - 3)^{2} + (y - 1)^{2} = 1 & &
 \\2x^{2} + 2y^{2} - 12x - 4y + 18 = 0 & &
\end{aligned}\]}
{Więcej niż dwa rozwiązania.}{Nie ma rozwiązania.}{Tylko jedno rozwiązanie.}{Dwa rozwiązania.}{}{}}
\LANGcs{
\Question{
Je dána soustava rovnic:
\[\begin{aligned}
 (x - 3)^{2} + (y - 1)^{2} = 1, & &
 \\2x^{2} + 2y^{2} - 12x - 4y + 18 = 0. & &
\end{aligned}\]
Vyberte správné tvrzení.}
{Soustava má více než dvě řešení.}{Soustava nemá řešení.}{Soustava má právě jedno řešení.}{Soustava má právě dvě řešení.}{}{}}
\LANGsk{
\Question{
Je daná sústava rovníc. Vyberte pravdivé tvrdenie.
\[\begin{aligned}
 (x - 3)^{2} + (y - 1)^{2} = 1 & &
 \\2x^{2} + 2y^{2} - 12x - 4y + 18 = 0 & &
\end{aligned}\]}
{Sústava má viac než dve riešenia.}{Sústava nemá žiadne riešenie.}{Sústava má práve jedno riešenie.}{Sústava má práve dve riešenia.}{}{}}
\LANGes{
\Question{
Dado el sistema de ecuaciones:
\[\begin{aligned}
 (x - 3)^{2} + (y - 1)^{2} = 1 & &
 \\2x^{2} + 2y^{2} - 12x - 4y + 18 = 0 & &
\end{aligned}\]
Identifica la proposición lógica verdadera.}
{El sistema tiene más de dos soluciones.}{El sistema no tiene soluciones.}{El sisteme tiene una solución.}{El sistema tiene dos soluciones.}{}{}}
\End

\Data\ID{9000031102}
\Updated{2020-08-23 08:08:24}
\Level{C}
\SubArea{Systems of equations and inequalities}
\LANGen{
\Question{
Solve the following system of equations and identify a correct statement.
\[\begin{aligned}
 (x - 1)^{2} + y^{2} = 1 & &
 \\(x - 4)^{2} + y^{2} = 4 & &
\end{aligned}\]}
{The system has a unique solution \(\left [x,y\right ]\),
where \(y = 0\).}{The system does not have any solution.}{The system has a unique solution \(\left [x,y\right ]\),
where \(y > 0\).}{The system has two solutions \(\left [x_{1},y_{1}\right ]\),
\(\left [x_{2},y_{2}\right ]\), where
\(y_{1} = -y_{2}\).}{}{}}
\LANGpl{
\Question{
Rozwiąż następujący układ równań i wybierz poprawną odpowiedź. 
\[\begin{aligned}
 (x - 1)^{2} + y^{2} = 1 & &
 \\(x - 4)^{2} + y^{2} = 4 & &
\end{aligned}\]}
{Tylko jedno rozwiązanie \(\left [x,y\right ]\),
gdzie \(y = 0\).}{Nie ma rozwiązania.}{Tylko jedno rozwiązanie \(\left [x,y\right ]\),
gdzie \(y > 0\).}{Dwa rozwiązania  \(\left [x_{1},y_{1}\right ]\),
\(\left [x_{2},y_{2}\right ]\), gdzie
\(y_{1} = -y_{2}\).}{}{}}
\LANGcs{
\Question{
Je dána soustava rovnic:
\[\begin{aligned}
 (x - 1)^{2} + y^{2} = 1, & &
 \\(x - 4)^{2} + y^{2} = 4. & &
\end{aligned}\]
Vyberte správné tvrzení.}
{Soustava má právě jedno řešení \(\left [x,y\right ]\),
pro něž platí, že \(y = 0\).}{Soustava nemá řešení.}{Soustava má právě jedno řešení \(\left [x,y\right ]\),
pro něž platí, že \(y > 0\).}{Soustava má právě dvě řešení \(\left [x_{1},y_{1}\right ]\),
\(\left [x_{2},y_{2}\right ]\), v nichž
platí \(y_{1} = -y_{2}\).}{}{}}
\LANGsk{
\Question{
Je daná sústava rovníc:
\[\begin{aligned}
 (x - 1)^{2} + y^{2} = 1, & &
 \\(x - 4)^{2} + y^{2} = 4. & &
\end{aligned}\] Vyberte správne tvrdenie.}
{Sústava má práve  jedno riešenie \(\left [x,y\right ]\),
kde \(y = 0\).}{Sústava nemá riešenie.}{Sústava má prve jedno riešenie \(\left [x,y\right ]\),
kde \(y > 0\).}{Sústava má práve dve riešenia \(\left [x_{1},y_{1}\right ]\),
\(\left [x_{2},y_{2}\right ]\), kde
\(y_{1} = -y_{2}\).}{}{}}
\LANGes{
\Question{
Dado el sistema de ecuaciones:
\[\begin{aligned}
 (x - 1)^{2} + y^{2} = 1 & &
 \\(x - 4)^{2} + y^{2} = 4 & &
\end{aligned}\]
Identifica la proposición lógica verdadera.}
{El sistema tiene solo una solución \(\left [x,y\right ]\),
where \(y = 0\).}{El sistema no tiene soluciones.}{El sistema tiene solo una solución \(\left [x,y\right ]\),
suponiendo que \(y > 0\).}{El sistema tiene dos soluciones \(\left [x_{1},y_{1}\right ]\),
\(\left [x_{2},y_{2}\right ]\), suponiendo que
\(y_{1} = -y_{2}\).}{}{}}
\End

\Data\ID{9000031106}
\Updated{2020-08-23 08:08:36}
\Level{C}
\SubArea{Systems of equations and inequalities}
\LANGen{
\Question{
Solve the following system of equations and identify a correct statement.
\[\begin{aligned}
 \sqrt{x + y} & = \left |x\right | & &
 \\x + y & = 4 & &
\end{aligned}\]}
{The system has two solutions \(\left [x_{1},y_{1}\right ]\),
\(\left [x_{2},y_{2}\right ]\), where
\(x_{1} = -x_{2}\).}{The system does not have any solution.}{The system has a unique solution.}{The system has two solutions \(\left [x_{1},y_{1}\right ]\),
\(\left [x_{2},y_{2}\right ]\), where
\(x_{1} = x_{2}\).}{}{}}
\LANGpl{
\Question{
Rozwiąż następujący układ równań i wybierz poprawną odpowiedź. 
\[\begin{aligned}
 \sqrt{x + y} & = \left |x\right | & &
 \\x + y & = 4 & &
\end{aligned}\]}
{Dwa rozwiązania \(\left [x_{1},y_{1}\right ]\),
\(\left [x_{2},y_{2}\right ]\), gdzie
\(x_{1} = -x_{2}\).}{Nie ma rozwiązania.}{Tylko jedno rozwiązanie.}{Dwa rozwiązania \(\left [x_{1},y_{1}\right ]\),
\(\left [x_{2},y_{2}\right ]\), gdzie
\(x_{1} = x_{2}\).}{}{}}
\LANGcs{
\Question{
Je dána soustava rovnic:
\[\begin{aligned}
 \sqrt{x + y} & = \left |x\right |, & &
 \\x + y & = 4. & &
\end{aligned}\]
Vyberte správné tvrzení.}
{Soustava má právě dvě řešení \(\left [x_{1},y_{1}\right ]\),
\(\left [x_{2},y_{2}\right ]\), v nichž
platí \(x_{1} = -x_{2}\).}{Soustava nemá řešení.}{Soustava má právě jedno řešení.}{Soustava má právě dvě řešení \(\left [x_{1},y_{1}\right ]\),
\(\left [x_{2},y_{2}\right ]\), v nichž
platí \(x_{1} = x_{2}\).}{}{}}
\LANGsk{
\Question{
Je daná sústava rovníc:
\[\begin{aligned}
 \sqrt{x + y} & = \left |x\right |, & &
 \\x + y & = 4. & &
\end{aligned}\] Vyberte správne tvrdenie.}
{Sústava má práve dve riešenia \(\left [x_{1},y_{1}\right ]\),
\(\left [x_{2},y_{2}\right ]\), kde
\(x_{1} = -x_{2}\).}{Sústava nemá žiadne riešenie.}{Sústava má práve jedno riešenie.}{Sústava má práve dve riešenia \(\left [x_{1},y_{1}\right ]\),
\(\left [x_{2},y_{2}\right ]\), kde
\(x_{1} = x_{2}\).}{}{}}
\LANGes{
\Question{
Dado el sistema de soluciones:
\[\begin{aligned}
 \sqrt{x + y} & = \left |x\right | & &
 \\x + y & = 4 & &
\end{aligned}\]
Identifica la proposición lógica verdadera.}
{El sistema tiene dos soluciones \(\left [x_{1},y_{1}\right ]\),
\(\left [x_{2},y_{2}\right ]\), suponiendo que
\(x_{1} = -x_{2}\).}{El sistema no tiene soluciones.}{El sistema tiene solo una solución.}{El sistema tiene dos soluciones \(\left [x_{1},y_{1}\right ]\),
\(\left [x_{2},y_{2}\right ]\), suponiendo que
\(x_{1} = x_{2}\).}{}{}}
\End

\Data\ID{9000031107}
\Updated{2020-08-23 08:08:58}
\Level{C}
\SubArea{Systems of equations and inequalities}
\LANGen{
\Question{
Solve the following system of equations and identify a correct statement.
\[\begin{aligned}
 \sqrt{x} = y & &
 \\x^{2} + y^{2} = 6 & &
\end{aligned}\]}
{The system has a unique solution.}{The system does not have any solution.}{The system has two solutions.}{The system has more than two solutions.}{}{}}
\LANGpl{
\Question{
Rozwiąż następujący układ równań i wybierz poprawną odpowiedź. 
\[\begin{aligned}
 \sqrt{x} = y & &
 \\x^{2} + y^{2} = 6 & &
\end{aligned}\]}
{Tylko jedno rozwiązanie.}{Nie ma rozwiązania.}{Dwa rozwiązania.}{Więcej niż dwa rozwiązania.}{}{}}
\LANGcs{
\Question{
Je dána soustava rovnic:
\[\begin{aligned}
 \sqrt{x} = y, & &
 \\x^{2} + y^{2} = 6. & &
\end{aligned}\]
Vyberte správné tvrzení.}
{Soustava má právě jedno řešení.}{Soustava nemá řešení.}{Soustava má právě dvě řešení.}{Soustava má více než dvě řešení.}{}{}}
\LANGsk{
\Question{
Vyriešte danú sústavu rovníc a vyberte pravdivé tvrdenie.
\[\begin{aligned}
 \sqrt{x} = y & &
 \\x^{2} + y^{2} = 6 & &
\end{aligned}\]}
{Sústava má práve jedno riešenie.}{Sústava nemá riešenie.}{Sústava má práve dve riešenia.}{Sústava má viac ako dve riešenia.}{}{}}
\LANGes{
\Question{
Dado el sistema de ecuaciones:
\[\begin{aligned}
 \sqrt{x} = y & &
 \\x^{2} + y^{2} = 6 & &
\end{aligned}\]
Identifica la proposición lógica verdadera.}
{El sistema tiene solo una solución.}{El sistema no tiene soluciones.}{El sistema tiene dos soluciones.}{El sistema tiene más de dos soluciones.}{}{}}
\End

\Data\ID{9000031108}
\Updated{2020-08-23 08:08:29}
\Level{C}
\SubArea{Systems of equations and inequalities}
\LANGen{
\Question{
Solve the following system of equations and identify a correct statement.
\[\begin{aligned}
 2x^{2} - y = 2 & &
 \\\left |x\right | + y = 1 & &
\end{aligned}\]}
{The system has two solutions.}{The system does not have any solution.}{The system has a unique solution.}{The system has more than two solutions.}{}{}}
\LANGpl{
\Question{
Rozwiąż następujący układ równań i wybierz poprawną odpowiedź. 
\[\begin{aligned}
 2x^{2} - y = 2 & &
 \\\left |x\right | + y = 1 & &
\end{aligned}\]}
{Dwa rozwiązania.}{Nie ma rozwiązania.}{Tylko jedno rozwiązanie.}{Więcej niż dwa rozwiązania.}{}{}}
\LANGcs{
\Question{
Je dána soustava rovnic:
\[\begin{aligned}
 2x^{2} - y = 2, & &
 \\\left |x\right | + y = 1. & &
\end{aligned}\]
Vyberte správné tvrzení.}
{Soustava má právě dvě řešení.}{Soustava nemá řešení.}{Soustava má právě jedno řešení.}{Soustava má více než dvě řešení.}{}{}}
\LANGsk{
\Question{
Vyriešte danú sústavu rovníc a vyberte pravdivé tvrdenie.
\[\begin{aligned}
 2x^{2} - y = 2 & &
 \\\left |x\right | + y = 1 & &
\end{aligned}\]}
{Sústava má práve dve riešenia.}{Sústava nemá riešenie.}{Sústava má práve jedno riešenie.}{Sústava má viac ako dve riešenia.}{}{}}
\LANGes{
\Question{
Dado el sistema de ecuaciones:
\[\begin{aligned}
 2x^{2} - y = 2 & &
 \\\left |x\right | + y = 1 & &
\end{aligned}\]
Identifica la proposición lógica verdadera.}
{El sistema tiene dos soluciones.}{El sistema no tiene soluciones.}{El sistema tiene solo una solución.}{El sistema tiene más de dos soluciones.}{}{}}
\End

\Data\ID{9000031109}
\Updated{2020-08-23 08:08:50}
\Level{C}
\SubArea{Systems of equations and inequalities}
\LANGen{
\Question{
Solve the following system of equations and identify a correct statement.
\[\begin{aligned}
 \left |x\right | = x + y & &
 \\\left |y\right | = 1 + x & &
\end{aligned}\]}
{The system has a unique solution.}{The system does not have any solution.}{The system has two solutions.}{The system has more than two solutions.}{}{}}
\LANGpl{
\Question{
Rozwiąż następujący ukad równań i wybierz poprawną odpowiedź. 
\[\begin{aligned}
 \left |x\right | = x + y & &
 \\\left |y\right | = 1 + x & &
\end{aligned}\]}
{Tylko jedno rozwiązanie.}{Nie ma rozwiązania.}{Dwa rozwiązania.}{Więcej niż dwa rozwiązania.}{}{}}
\LANGcs{
\Question{
Je dána soustava rovnic:
\[\begin{aligned}
 \left |x\right | = x + y, & &
 \\\left |y\right | = 1 + x. & &
\end{aligned}\]
Vyberte správné tvrzení}
{Soustava má právě jedno řešení.}{Soustava nemá řešení.}{Soustava má právě dvě řešení.}{Soustava má více než dvě řešení.}{}{}}
\LANGsk{
\Question{
Vyriešte danú sústavu rovníc a vyberte pravdivé tvrdenie.
\[\begin{aligned}
 \left |x\right | = x + y & &
 \\\left |y\right | = 1 + x & &
\end{aligned}\]}
{Sústava má práve jedno riešenie.}{Sústava nemá žiadne riešenie.}{Sústava má práve dve riešenia.}{Sústava má viac ako dve riešenia.}{}{}}
\LANGes{
\Question{
Dado el sistema de ecuaciones:
\[\begin{aligned}
 \left |x\right | = x + y & &
 \\\left |y\right | = 1 + x & &
\end{aligned}\]
Identifica la proposición lógica verdadera.}
{El sistema tiene solo una solución.}{El sistema no tiene soluciones.}{El sistema tiene dos soluciones.}{El sistema tiene más de dos soluciones.}{}{}}
\End

\Data\ID{9000031110}
\Updated{2020-08-23 08:08:23}
\Level{C}
\SubArea{Systems of equations and inequalities}
\LANGen{
\Question{
Solve the following system of equations and identify a correct statement.
\[\begin{aligned}
 \left |x - 2\right | & = y & &
 \\\left |y + 2\right | & = x - 6 & &
\end{aligned}\]}
{The system has no solution.}{The system has a unique solution.}{The system has two solutions.}{The system has more than two solutions.}{}{}}
\LANGpl{
\Question{
Rozwiąż następujący układ równań i wybierz poprawną odpowiedź.
\[\begin{aligned}
 \left |x - 2\right | & = y & &
 \\\left |y + 2\right | & = x - 6 & &
\end{aligned}\]}
{Nie ma rozwiązania.}{Tylko jedno rozwiązanie.}{Dwa rozwiązania.}{Więcej niż dwa rozwiązania.}{}{}}
\LANGcs{
\Question{
Je dána soustava rovnic:
\[\begin{aligned}
 \left |x - 2\right | & = y, & &
 \\\left |y + 2\right | & = x - 6. & &
\end{aligned}\]}
{Soustava nemá řešení.}{Soustava má právě jedno řešení.}{Soustava má právě dvě řešení.}{Soustava má více než dvě řešení.}{}{}}
\LANGsk{
\Question{
Vyriešte danú sústavu rovníc a vyberte pravdivé tvrdenie.
\[\begin{aligned}
 \left |x - 2\right | & = y & &
 \\\left |y + 2\right | & = x - 6 & &
\end{aligned}\]}
{Sústava nemá žiadne riešenie.}{Sústava má práve jedno riešenie.}{Sústava má práve dve riešenia.}{Sústava má viac ako dve riešenia.}{}{}}
\LANGes{
\Question{
Dado el sistema de ecuaciones:
\[\begin{aligned}
 \left |x - 2\right | & = y & &
 \\\left |y + 2\right | & = x - 6 & &
\end{aligned}\]
Identifica la proposición lógica verdadera.}
{El sistema no tiene soluciones.}{El sistema tiene solo una solución.}{El sistema tiene dos soluciones.}{El sistema tiene más de dos soluciones.}{}{}}
\End
